%=============================================================================
%  chapter6.tex -- What This Model Cannot Say
%
%  Uses macros from thesis_main.tex Part 2. Defines no new notation.
%  Written against outputs/writing_specification.md (10 Sep 2026): the reader
%  is a first-year M.Tech Chemical Engineering student who has never opened a
%  plasma physics book. Every acronym is expanded at first use in this
%  chapter; every physical picture precedes its formalism.
%
%  NOTE TO THE COMPILER: thesis_main.tex currently carries this chapter as an
%  inline skeleton (\chapter{What This Model Cannot Say} ... \label{ch:limits}
%  with the section labels sec:opacity, sec:molecules, sec:nmax,
%  sec:open_system, sec:other_assumptions). Replace that skeleton with
%  %=============================================================================
%  chapter6.tex -- What This Model Cannot Say
%
%  Uses macros from thesis_main.tex Part 2. Defines no new notation.
%  Written against outputs/writing_specification.md (10 Sep 2026): the reader
%  is a first-year M.Tech Chemical Engineering student who has never opened a
%  plasma physics book. Every acronym is expanded at first use in this
%  chapter; every physical picture precedes its formalism.
%
%  NOTE TO THE COMPILER: thesis_main.tex currently carries this chapter as an
%  inline skeleton (\chapter{What This Model Cannot Say} ... \label{ch:limits}
%  with the section labels sec:opacity, sec:molecules, sec:nmax,
%  sec:open_system, sec:other_assumptions). Replace that skeleton with
%  %=============================================================================
%  chapter6.tex -- What This Model Cannot Say
%
%  Uses macros from thesis_main.tex Part 2. Defines no new notation.
%  Written against outputs/writing_specification.md (10 Sep 2026): the reader
%  is a first-year M.Tech Chemical Engineering student who has never opened a
%  plasma physics book. Every acronym is expanded at first use in this
%  chapter; every physical picture precedes its formalism.
%
%  NOTE TO THE COMPILER: thesis_main.tex currently carries this chapter as an
%  inline skeleton (\chapter{What This Model Cannot Say} ... \label{ch:limits}
%  with the section labels sec:opacity, sec:molecules, sec:nmax,
%  sec:open_system, sec:other_assumptions). Replace that skeleton with
%  %=============================================================================
%  chapter6.tex -- What This Model Cannot Say
%
%  Uses macros from thesis_main.tex Part 2. Defines no new notation.
%  Written against outputs/writing_specification.md (10 Sep 2026): the reader
%  is a first-year M.Tech Chemical Engineering student who has never opened a
%  plasma physics book. Every acronym is expanded at first use in this
%  chapter; every physical picture precedes its formalism.
%
%  NOTE TO THE COMPILER: thesis_main.tex currently carries this chapter as an
%  inline skeleton (\chapter{What This Model Cannot Say} ... \label{ch:limits}
%  with the section labels sec:opacity, sec:molecules, sec:nmax,
%  sec:open_system, sec:other_assumptions). Replace that skeleton with
%  \input{chapter6} -- those labels are reused below, so every existing
%  cross-reference from Chapters 3 and 5 continues to resolve.
%
%  Sources for every number, in order of appearance:
%    outputs/findings_10_four_agent_review.md   ss 3.1, 4.1-4.4, 7.3, 8, 10,
%                                               and the ADDENDUM (A.1-A.6)
%    validation/lyman_trapping/lyman_trapping.txt    (run 9-10 Sep 2026,
%                                               src/validation/verify_lyman_trapping.py)
%    thesis_tex/chapter5_C5D.tex:284    (the "It is not detachment" wording)
%    outputs/findings_09_central_quantity_misnamed.md   ss 8, 9
%=============================================================================

\chapter{What This Model Cannot Say}
\label{ch:limits}

Chapter~\ref{ch:results} handed the reader two things: a map of how wrong a
tabulated $H_\alpha/H_\beta$ inversion can be while the ground state is still
catching up, and a mechanism that explains the map's shape. The map has a ridge
in it, the ridge has a closed-form explanation, and the explanation was tested
three separate ways. At this point in the argument the reader believes a
result.

This chapter's job is to bound that belief without destroying it. Every number
in Chapter~\ref{ch:results} came out of one $43$-state matrix describing one
kind of atom under one set of assumptions, and some of those assumptions are
badly violated in exactly the corner of the map where the numbers are largest.
An examiner will find that corner. It is better to arrive there first, with
arithmetic.

\medskip
\noindent\textbf{The standard this chapter holds itself to.} A limitation that
has been \emph{computed} is worth more than a limitation that has been hedged.
Where an objection could be turned into a calculation, it was, and the
calculation is reported here with the number that would have refuted the
result. Two objections were expected to be fatal and were not; they are given
the same space as the ones that bite. Where an objection cannot be computed
inside a model with no spatial dimension, it is stated as a \emph{conditional},
``\emph{if} this holds, the error is as computed'', and not softened into a
caveat sentence. Where a check has simply not been run, this chapter says so.

\medskip
\noindent\textbf{Three abbreviations recur below} and are expanded here for
this chapter: \emph{CR} for collisional--radiative, the class of model built in
Chapter~\ref{ch:theory}; \emph{QSS} for quasi-steady-state, the Bodenstein-style
approximation that eliminates the fast excited levels; and \emph{CRE} for
collisional--radiative equilibrium, the fully settled steady state that the
diagnostic lookup table assumes and that this thesis measures the cost of
assuming. An \emph{ELM} (edge-localised mode) is the repetitive burst of
heat and particles that a tokamak edge throws at the divertor every few
milliseconds, and it is the transient every ``event duration'' below refers to;
$\taudrive = \SI{100}{\micro\second}$ throughout, the duration used in
Chapter~\ref{ch:results} and the aggressive end of the range; the
ITER-pedestal extrapolation is \SI{506}{\micro\second}
\cite[Sec.~5]{Loarte2003} (Section~\ref{sec:persistence}).

%-----------------------------------------------------------------------------
\section{Optical thickness: when the plasma cannot let its own light out}
\label{sec:opacity}
% QUESTION: the model lets every emitted photon escape the plasma -- where is
% that false, what does it do to the error map, and how much of the map
% survives it?

\subsection{The question, and the picture behind it}

Chapter~\ref{ch:atoms} gave every downward transition an Einstein $A$
coefficient: the probability per second that an excited atom drops to a lower
level and emits a photon. Those coefficients entered the CR matrix $\Lmat$
bare. Entering bare means one specific physical assumption: that every photon
emitted anywhere in the plasma leaves it and never comes back. A plasma for
which this is true is called \emph{optically thin}.

A chemical engineer has met the failure of this assumption in another costume.
Consider a reaction whose product is removed by desorption from a catalyst
surface. If the product leaves the reactor, the desorption step runs at its
intrinsic rate. If the product cannot leave, if it re-adsorbs somewhere else
in the bed and is put back onto a site, then the \emph{effective} desorption
rate is smaller than the intrinsic one, and the whole network downstream of that
step shifts. Radiation trapping is that, with photons.

Here is the specific chain. The strongest transition in hydrogen is
Lyman-$\alpha$: an atom in the $2p$ level drops to the ground state $1s$ and
emits an ultraviolet photon. If the gas around it contains many ground-state
atoms, that photon has a good chance of being absorbed by one of them, which
puts \emph{that} atom back into $2p$. Nothing has left the plasma; a
$2p$ population has simply been moved from one place to another. Averaged over
the volume, the $2p$ level decays more slowly than its $A$ coefficient says. The
standard way to write this down is to multiply the coefficient by an
\emph{escape factor} $\Theta_P$ between $0$ and $1$: $\Theta_P = 1$ is a
transparent plasma, $\Theta_P = 10^{-5}$ is one in which the transition is
effectively switched off. The formulation used here is Holstein's population
escape factor~\cite{Holstein1947}.

\textbf{Why this reaches the populations, and therefore the result.} If $2p$
lives longer, atoms pile up in $n=2$. Atoms in $n=2$ are only
\SI{10.2}{\electronvolt} from the ground state but only
\SI{3.4}{\electronvolt} from the continuum, so they are far easier to ionise
than ground-state atoms. A longer $n=2$ lifetime therefore means more
\emph{stepwise} ionisation, ionisation via an excited level rather than
directly out of the ground state, which means fewer neutrals and a faster
turnover of the ground-state reservoir. The ground-state reservoir is the whole
mechanism of this thesis. Trapping is not a cosmetic correction to a line
intensity; it acts on the slow clock.

\subsection{Where the assumption fails, in numbers}

The formalism needed is one line. A photon travelling a distance $D$ through a
gas of absorbers is attenuated by $e^{-\tau}$, where the \emph{optical depth}
$\tau = \sigma_0\,\ngs\,D$ is dimensionless: $\ngs$ is the ground-state density
in \si{\per\cubic\centi\metre}, $D$ the path length in \si{\centi\metre}, and
$\sigma_0$ the absorption cross section at line centre in
\si{\square\centi\metre}. For a Doppler-broadened line,
$\sigma_0 = 2.654\times10^{-2}\,f/(\sqrt{\pi}\,\Delta\nu_D)$, with $f$ the
oscillator strength (a dimensionless measure of how strongly the transition
couples to light) and $\Delta\nu_D$ the Doppler width set by the temperature of
the absorbing atoms. $\tau \ll 1$ is transparent; $\tau \gg 1$ is opaque.

The number that matters is $\ngs$, not $\nel$, and Chapter~\ref{ch:theory}
already established the awkward fact that these two run in opposite directions:
the plasma is most neutral where it is coldest. The worst point of the error map
is $\Te = \SI{1.0}{\electronvolt}$,
$\nel = 5.18\times10^{13}\,\si{\per\cubic\centi\metre}$, the maximum of
$\epsplat$ over all $680$ window-resolvable points of the map, cold edge
included. There the model's own CR solution gives $\ngs/\nion = 28.3$, so
quasi-neutrality fixes
$\ngs = 1.47\times10^{15}\,\si{\per\cubic\centi\metre}$. At
\SI{1}{\electronvolt} that gives
$\sigma_0 = 5.47\times10^{-14}\,\si{\square\centi\metre}$, and the Lyman-$\alpha$
optical depth per centimetre of path is

\begin{center}
\begin{tabular}{lr}
  \toprule
  $\Te$ (eV) & $\tau_{\rm Ly\alpha}$ per \si{\centi\metre} at
              $\nel = 5.18\times10^{13}\,\si{\per\cubic\centi\metre}$ \\
  \midrule
  1.00 & 80.3 \\
  1.15 & 10.6 \\
  1.60 & 0.224 \\
  2.02 & 0.0255 \\
  2.95 & 0.00163 \\
  \bottomrule
\end{tabular}
\end{center}

\noindent An optical depth of $80$ per centimetre means a Lyman-$\alpha$ photon
travels about a tenth of a millimetre before it is reabsorbed. Over
\SI{5}{\centi\metre} of plasma the line-centre depth
over the half slab is $201$ and the population escape factor is
$1.2\times10^{-3}$; roughly one photon in eight hundred gets out. Chapter~%
\ref{ch:theory} quotes $3\times10^{-5}$ at a different grid point, a different
slab thickness and a different density, so the two are not the same quantity and
the agreement in order of magnitude should not be read as a reproduction of
one by the other.
The blunt reading is that the effective $A_{2p\to1s}$ belonging in $\Lmat$ at
that point is wrong, self-consistently, by about $1.8$ orders of magnitude, or
$2.9$ orders on the thin one-shot estimate
(\texttt{verify\_lyman\_optical\_depth.py}).

Two rows down the same table the problem has evaporated. At
\SI{2.02}{\electronvolt} the depth per centimetre is $0.0255$, and at the
benchmark point (\benchTe, \benchne) it is $0.0039$; a photon there crosses
metres of plasma before anything happens to it. \textbf{The measured quantity
is not ``how opaque is this plasma'' but ``how fast does opacity switch on as
the plasma cools'',} and the answer is a factor of $5\times10^{4}$ between
\SI{1}{\electronvolt} and \SI{2.95}{\electronvolt}.

That is why the results of Chapter~\ref{ch:results} were declared for
$\Te \gtrsim \SI{2}{\electronvolt}$ in the first place
(Section~\ref{sec:two_references}). It is also, uncomfortably, where the
headline was living. Recounting the \SI{100}{\micro\second} breakdown census against optical depth
rather than against temperature: of the $202$ points at which the
\SI{100}{\micro\second} time-averaged estimate of the error exceeds \SI{10}{\percent}, an extremum count over
the $680$ window-resolvable points of the full grid, $157$ lie below
\SI{2}{\electronvolt}, $85$ have
$\tau_{\rm Ly\alpha}(\SI{5}{\centi\metre}) > 1$, and $15$ have it above $100$,
on the half-slab depth used above.

An earlier draft gave $105$ and $26$, a convention error rather than a
recomputation: those counts are reproduced only by combining the full-path
depth $\kappa_0 D$ with a Doppler width $\sqrt{kT_n/m}$, short of the correct
$\sqrt{2kT_n/m}$ by $\sqrt{2}$. With the correct width the census is $85$ and
$15$ on the half-slab depth quoted here and $100$ and $23$ on the full-path
depth. A deuterium cross-section with the full-path depth also returns $105$
and $26$, since the isotope mass enters the width as the same $\sqrt{2}$, so
the two errors are indistinguishable from the counts alone. All seven
conventions are tabulated in \texttt{validation/lyman\_optical\_depth/}.
Nothing in the defended range depends on the choice: every crossing lies below
\SI{2}{\electronvolt} except two, both at $D = \SI{20}{\centi\metre}$.

\subsection{What we expected trapping to do, and what it did}

The expectation, stated before the calculation was run, was that trapping would
demolish the map: an effective $A$ coefficient wrong by about three orders of
magnitude at the worst point ought to invalidate everything computed there and
nearby, and the project's own records warned that under strong trapping the
fast relaxation time $\taurel$ rises to \SI{2.05}{\micro\second} and the
timescale separation $\Msep$ falls by a factor $158$. That expectation was
wrong.

So the calculation was done. The Lyman series was given a self-consistent
Holstein escape factor and the entire error map rebuilt. Self-consistent matters
here: $\Theta_P$ depends on $\ngs$, and $\ngs$ depends on $\Theta_P$ through the
CR balance, so the two were iterated to a fixed point rather than evaluated once.
The fixed point converged at all $400$ grid points for every slab thickness
tested, and non-convergence was made to raise an error rather than silently
return the last iterate. All $14$ Lyman channels were trapped, the seven
resolved $n\mathrm{P}\to1\mathrm{S}$ transitions and the seven bundled ones,
not Lyman-$\alpha$ alone.

Three gates were passed before any result was looked at. The oscillator
strengths were derived from this repository's own $A$ coefficients by inverting
the Einstein relation rather than imported, and come out as $f = 0.4162$,
$0.0791$, $0.0290$, $0.0139$ for Lyman-$\alpha$ through Lyman-$\delta$, the
accepted literature values to four figures. The derived Lyman-$\alpha$
$\sigma_0$ agrees with an independent implementation to \SI{0.059}{\percent},
after a first run that failed at \SI{1156}{\percent} on a $4\pi$ factor; and
rebuilding the matrix with trapping switched off reproduces the canonical
$\Lmat$ exactly, so every trapped number below is a difference against the real
matrix (Section~\ref{sec:severe_checks}).

\textbf{One new quantity had to be invented to do this at all, and it is not in
the model.} The escape factor depends on how far a photon must travel to get
out, so it depends on the size of the plasma, the slab thickness $D$. A model
with no spatial dimension has no $D$. It was therefore swept, from
\SI{1}{\centi\metre} to \SI{20}{\centi\metre}, and never fixed; every trapped
number below carries the $D$ it was computed at. Reporting a single trapped
number without its $D$ would be reporting an assumption as a measurement.

\begin{table}[tbp]
  \centering
  \caption[Breakdown census under Lyman trapping]{Effect of radiation
  trapping on the \SI{100}{\micro\second} breakdown census. Each entry counts the grid points at
  which the time-averaged estimate of the CR-equilibrium-reference error,
  averaged over a chosen duration $\taudrive = \SI{100}{\micro\second}$, exceeds
  \SI{10}{\percent}, out of the points at which a
  quasi-steady-state window is resolvable; ``worst'' is the largest estimate
  within that scope. Rows differ only in the assumed slab thickness $D$ used to
  compute the self-consistent Holstein escape factor applied to all $14$ Lyman
  transitions. The first row is the optically thin matrix used throughout
  Chapter~\ref{ch:results}. Above \SI{2}{\electronvolt} the count does not move
  at all across a twentyfold range of $D$. The $D = \SI{20}{\centi\metre}$ row
  has a slightly smaller denominator because two points lose their resolvable
  window under trapping. Source:
  \texttt{validation/lyman\_trapping/lyman\_trapping.txt}.}
  \label{tab:trapping_census}
  \footnotesize
  \begin{tabular}{lccc}
    \toprule
    matrix & all window-resolvable & $\Te \ge \SI{2}{\electronvolt}$ &
          $\Te \ge \SI{2}{\electronvolt}$, $\nel \ge 10^{14}$ \\
    \midrule
    optically thin (canonical)          & $202/680$, worst $0.3868$ & $45/448$, worst $0.1748$ & $0/108$, worst $0.0717$ \\
    trapped, $D = \SI{1}{\centi\metre}$  & $176/680$, worst $0.2925$ & $45/448$, worst $0.1748$ & $0/108$, worst $0.0712$ \\
    trapped, $D = \SI{5}{\centi\metre}$  & $170/680$, worst $0.2550$ & $45/448$, worst $0.1746$ & $0/108$, worst $0.0694$ \\
    trapped, $D = \SI{20}{\centi\metre}$ & $160/678$, worst $0.2237$ & $45/446$, worst $0.1740$ & $0/106$, worst $0.0627$ \\
    \bottomrule
  \end{tabular}
\end{table}

\subsection{Above two electron-volts, nothing moves}

Table~\ref{tab:trapping_census} contains the strongest robustness statement in
this thesis, and it is the opposite of what was expected.

\textbf{Above \SI{2}{\electronvolt} the breakdown count does not change at all,
$45$ of $448$ points in every run, and the worst case within that set
moves by \SI{0.5}{\percent}, from $0.1748$ to $0.1740$, across a twentyfold
range in assumed slab thickness.} The escape factor at the benchmark point
(\benchTe, \benchne) runs from $0.9986$ at $D = \SI{1}{\centi\metre}$ to
$0.973$ at $D = \SI{20}{\centi\metre}$. Even at the thickest slab tested,
\SI{97}{\percent} of Lyman-$\alpha$ photons still leave: the plasma is optically
thin there in the sense that actually matters, and assuming so costs less than
the resolution of the map.

This converts a disclaimer into a measurement. The
$\Te \gtrsim \SI{2}{\electronvolt}$ restriction was chosen on an escape-factor
argument before any of this was recomputed, and it sits almost exactly at the
temperature above which trapping stops mattering: a guess that has now been
checked, and survived. The defensible claim of Chapter~\ref{ch:results}, $45$
of $448$ points above \SI{2}{\electronvolt} at \SI{100}{\micro\second} and $24$
of $448$ at \SI{506}{\micro\second} (Section~\ref{sec:persistence}), is
invariant under the one physical process most likely to have invalidated it.

\subsection{Below two electron-volts, the cold-edge number is inflated
threefold}

The same calculation is much less kind to the cold edge, and it should be,
because that is where the optical depth was $80$ per centimetre.
Table~\ref{tab:trapping_worstpoint} follows the worst point of the optically
thin map through the sweep.

\begin{table}[tbp]
  \centering
  \caption[Lyman trapping at the cold-edge worst point]{The cold-edge worst
  point of the optically thin error map, at $\Te = \SI{1.0}{\electronvolt}$ and
  $\nel = 5.18\times10^{13}\,\si{\per\cubic\centi\metre}$ under a heating step,
  is the maximum of $\epsplat$ over all $680$ window-resolvable grid points; it
  is traced here as a function of the assumed slab thickness $D$. $\epsplat$ is the plateau
  error, the fractional error a lookup table makes while the ground state is
  stale; the ELM bound is its average over an event of
  $\taudrive = \SI{100}{\micro\second}$. The plateau error falls by a factor
  $2.5$ to $3.3$, but the diagnostic's intrinsic sensitivity
  $|\ffrac{3}-\ffrac{4}|$ barely moves: what collapses is how far the ground
  state is displaced, not how strongly the line ratio responds to it. Source:
  \texttt{validation/lyman\_trapping/lyman\_trapping.txt}.}
  \label{tab:trapping_worstpoint}
  \small
  \begin{tabular}{lrrrrrrr}
    \toprule
    $D$ (cm) & $\epsplat$ & ELM bound & $\tauslow$ (s) & $\Msep$ &
    $\ffrac{3}$ & $\ffrac{4}$ & $|\ffrac{3}-\ffrac{4}|$ \\
    \midrule
    $0$ (thin) & $0.38690$ & $0.38682$ & $2.3324\times10^{-1}$ & $3.74\times10^{7}$ & $0.7728$ & $0.2862$ & $0.4866$ \\
    $1$        & $0.15748$ & $0.15739$ & $9.0651\times10^{-2}$ & $2.75\times10^{6}$ & $0.6643$ & $0.2055$ & $0.4589$ \\
    $5$        & $0.13274$ & $0.13260$ & $4.6940\times10^{-2}$ & $6.13\times10^{5}$ & $0.5967$ & $0.1793$ & $0.4174$ \\
    $20$       & $0.11570$ & $0.11550$ & $2.8500\times10^{-2}$ & $2.14\times10^{5}$ & $0.5139$ & $0.1501$ & $0.3639$ \\
    \bottomrule
  \end{tabular}
\end{table}

\textbf{The \SI{38.7}{\percent} becomes \SIrange{11.6}{15.7}{\percent}}, a
factor of $2.5$ to $3.3$, and which value inside that range it takes is set
by a slab thickness the model does not contain. The point still exceeds the
\SI{10}{\percent} threshold at every thickness tested, so the qualitative
statement that the cold corner breaks down survives intact; the \emph{number}
does not, and it must never appear without its trapped range and without the
optical depth that produced it.

The route by which the number falls is itself a result, and it is not the
obvious route. The obvious guess is that trapping desensitises the diagnostic,
that a trapped plasma simply responds less. It does not. The sensitivity
$|\ffrac{3}-\ffrac{4}|$, which is the factor in the exact decomposition of
Section~\ref{sec:mechanism} measuring how strongly the line ratio answers to its
reservoir, moves by only \SI{6}{\percent} at $D = \SI{1}{\centi\metre}$ and
\SI{25}{\percent} at $D = \SI{20}{\centi\metre}$. What collapses is the other
factor: the displacement.

The chain is the one set out at the start of this section, now with numbers on
it. Trapping lengthens the $n=2$ lifetime; the longer lifetime raises stepwise
ionisation; the ground-state population falls; and the slow clock $\tauslow$
shortens by a factor $8$ at the worst point. At $\Te = \SI{1.0}{\electronvolt}$,
$\nel = 1.93\times10^{13}\,\si{\per\cubic\centi\metre}$, $\ngs$ drops from
$2.99\times10^{14}$ to $7.26\times10^{13}\,\si{\per\cubic\centi\metre}$ between
$D = 1$ and \SI{20}{\centi\metre}. A reservoir that turns over faster is a less
stale reservoir, so $|\Delta\ln\xscale|$ roughly halves. This is exactly the
two-axis attribution of Section~\ref{sec:temperature_mechanism} arriving from a
third direction: trapping acts on the displacement, not on the sensitivity.

\subsection{Below 1.15 electron-volts the ridge location is not determined}

One more thing moves at the cold edge, and it is a location rather than a
magnitude. Taking the density at which $\epsplat$ is maximal, row by row: at
$\Te \ge \SI{1.6}{\electronvolt}$ that maximum falls on
$\nel = 1.93\times10^{13}\,\si{\per\cubic\centi\metre}$ in every run, optically
thin and trapped alike, and the values themselves barely register the change:
at the benchmark temperature $\epsplat = 0.12166$, $0.12166$, $0.12165$,
$0.12163$ for $D = 0, 1, 5$ and \SI{20}{\centi\metre}, with
$|\ffrac{3}-\ffrac{4}|$ changing in the fifth decimal place. Below
\SI{1.15}{\electronvolt} the maximum wanders across three density columns, a
factor of $7$, as a function of an assumed parameter. \textbf{The two coldest
rows of the grid cannot support a claim about where the ridge is at all.}

That result cuts in the mechanism's favour, which is worth saying explicitly
because it is easy to miss. The ridge attribution of
Section~\ref{sec:mechanism} was verified over $\Te \ge \SI{2}{\electronvolt}$,
and $\Te \ge \SI{2}{\electronvolt}$ is precisely the range in which trapping has
now been shown to be irrelevant. The agreement with Griem's boundary criterion
is therefore not an artifact of neglected opacity; it was established where
opacity does not act.

\subsection{A contradiction resolved, and what is still open}

Two statements about trapping had been circulating in this project's records:
that $\taurel$ is insensitive to Lyman-$\alpha$ trapping to within
\SI{1}{\percent}, and that under strong trapping $\taurel$ rises to
\SI{2.05}{\micro\second} with $\Msep$ falling $158$-fold. They look
contradictory. Measured directly, both are true and neither had stated its
scope: at the benchmark point (\benchTe, \benchne) $\taurel$ moves from
$2.2769\times10^{-9}$ to $2.2947\times10^{-9}\,\si{\second}$ at
$D = \SI{20}{\centi\metre}$, a change of \SI{0.78}{\percent}; at
$\Te = \SI{1.0}{\electronvolt}$,
$\nel = 1.93\times10^{13}\,\si{\per\cubic\centi\metre}$ it moves from
$8.8842\times10^{-9}$ to $2.4289\times10^{-7}\,\si{\second}$, a factor of $27.3$.
The ``within \SI{1}{\percent}'' claim is a benchmark-point claim and fails by
more than an order of magnitude at the cold corner. This is the same failure
mode as quoting any other extremum without the set it is an extremum over, and
neither sentence may travel again without its regime attached.

What this calculation does not settle should be as clear as what it does. The
geometry is assumed, not derived: a homogeneous slab, uniform emission, an
isotropic radiation field, a Doppler line profile, evaluated at the plasma
centre. A $0$-D model cannot check any of those. One assumption inside it
\emph{was} checked and is not load-bearing: the temperature of the absorbing
atoms, which sets the Doppler width. Fixing it at \SI{3}{\electronvolt}, the
value a molecule dissociating by the Franck--Condon process would give, instead
of at $\Te$, changes the census at $D = \SI{5}{\centi\metre}$ from $170/680$ and
a worst case of $0.2550$ to $173/680$ and $0.2653$. Continuum lowering is
untouched here and remains open. Opacity in the Balmer lines themselves, which
would act on the observed ratio directly rather than through the level
populations, has been computed: the largest $H_\alpha$ optical depth above
\SI{2}{\electronvolt} is $0.010$, and the worst cold cell, $[0,7]$, reaches
$0.214$ (Section~\ref{sec:balmer_equivalence},
\texttt{validation/balmer\_optical\_depth/}).

%-----------------------------------------------------------------------------
\section{Transport: the atoms do not stay where the model puts them}
\label{sec:transport}
% QUESTION: the slow clock is the ionisation time of a stationary atom -- what
% happens to the result when the atom leaves before it ionises?

\subsection{The question}

This is the objection that cannot be answered inside this model, and it is a
more serious threat to the cold-edge number than anything in
Section~\ref{sec:opacity}. It is also the simplest to state.

\subsection{The picture: a reactor with no walls}

The model of Chapter~\ref{ch:theory} is a batch reactor. A fixed parcel of gas
is given a temperature, and its composition evolves until it stops changing. The
slow eigenvalue of that system, $\tauslow$, is how long the ground-state
population takes to reach its new value, and it is long, because ionising an
atom out of the ground state is slow.

A divertor is not a batch reactor. It is closer to a continuous-flow reactor
with an enormous throughput, in which the residence time of a molecule is far
shorter than the time it needs to react. Neutral atoms are born at the target
surface when ions strike it and recombine there, which is \emph{recycling},
and they fly back into the plasma. They do not wait around to be ionised
locally; the ones that do not ionise leave, and fresh ones arrive continuously.

The comparison is not close. At the worst point of the map $\tauslow$ is
\SI{233}{\milli\second}: nearly a quarter of a second for a stationary atom to
be ionised. A deuterium atom at \SIrange{1}{3}{\electronvolt} would cross a
\SI{10}{\centi\metre} plasma ballistically in \SIrange{6}{10}{\micro\second}, but
it does not travel ballistically: resonant charge exchange with the surrounding
ions randomises its direction after about \SI{1.7}{\centi\metre} at this
density, so escape is diffusive and takes closer to \SI{72}{\micro\second}. That
correction runs in this model's favour. The threshold at which the worst point
would fall below \SI{10}{\percent} is a ground-state renewal time of about
\SI{26}{\micro\second}, so the result survives by a factor of $2.8$ rather than
failing outright. It survives only because charge exchange traps the neutrals,
which is a thin margin to rest on and is stated as one.
That is one grid point, and it is one this thesis does not defend:
$[0,4]$ lies below the \SI{2}{\electronvolt} scope boundary.
Section~\ref{sec:transport_selection} asks the same question at the points that
do carry the result, and gets a worse answer.

\begin{quote}
\textbf{In a divertor, the ground-state neutral density at a given point is not
set by local ionisation balance. It is set by recycling.}
\end{quote}

\subsection{Why this bears on the result specifically}

The error of Chapter~\ref{ch:results} exists for one reason: after a temperature
step, the ground-state reservoir is \emph{stale}. Its population has not moved
while the temperature has, the excited levels equilibrate against a reservoir
that belongs to the old conditions, and the lookup table, which assumes the
reservoir finished settling, reads the wrong answer. Everything hangs on the
reservoir not being refreshed.

Neutral transport refreshes it. Any process that replaces the ground-state
population on a timescale shorter than $\tauslow$ resets the staleness, and this
model contains no such process; the \SI{26}{\micro\second} renewal threshold
quoted above is the map's own measure of how fast that process would have to be.

\subsection{The assumption fails hardest where the result is generated}
\label{sec:transport_selection}

The worst point was the wrong place to test transport: $[0,4]$ lies below the
scope boundary of Section~\ref{sec:scope}. The pairs that carry the defended
claim are the $45$ above \SI{2}{\electronvolt} whose \SI{100}{\micro\second}
estimate exceeds \SI{10}{\percent} (Table~\ref{tab:elm_census}), and there the
answer is different.

Comparing the neutral escape time with $\tauslow$ pair by pair over the grid
gives Table~\ref{tab:transport_selection}. The escape model is the one used
above, stated as a formula: the atom moves at $v=\sqrt{2k\Tn/m_{\mathrm D}}$
with $\Tn=\Te$, resonant charge exchange with the ambient ions randomises its
direction after $\lambda = v/(\nel\langle\sigma v\rangle_{\mathrm{CX}})$ with
$\langle\sigma v\rangle_{\mathrm{CX}} = \SI{1.1e-8}{\cubic\centi\metre\per\second}$,
and escape from a boundary a distance $L$ away takes $L/v$ when
$\lambda \ge L$ and $4L^{2}/\pi^{2}D$ with $D = v\lambda/3$ otherwise. At
$[0,4]$ it returns $\lambda = \SI{1.72}{\centi\metre}$ and
$\tauesc = \SI{72.3}{\micro\second}$, the values used above
(\texttt{verify\_transport\_selection.py}).

\begin{table}[htbp]
\centering
\caption{Ratio of neutral escape time to $\tauslow$, over the $45$ defended
breakdown pairs and over the other $635$ analysed pairs. $L$ is the distance
from the emitting parcel to the plasma boundary. ``Satisfies'' means
$\tauesc > \tauslow$, that is, the atom is ionised before it leaves, which is
what the closed parcel assumes.}
\label{tab:transport_selection}
\begin{tabular}{llrrr}
  \toprule
  $L$ & set & median $\tauesc/\tauslow$ & largest & satisfies \\
  \midrule
  \SI{10}{\centi\metre} & the 45 breakdown pairs & $0.0079$ & $0.159$ & $0/45$ \\
  \SI{10}{\centi\metre} & the other 635 & $0.0471$ & $384$ & $168/635$ \\
  \SI{20}{\centi\metre} & the 45 breakdown pairs & $0.0192$ & $0.636$ & $0/45$ \\
  \SI{20}{\centi\metre} & the other 635 & $0.129$ & $1535$ & $227/635$ \\
  \bottomrule
\end{tabular}
\end{table}

\textbf{Not one of the $45$ pairs that carry the defended result satisfies the
assumption used to compute it, at either scale length, while a quarter to a
third of the pairs that do not break down satisfy it comfortably.}

A count of zero is worth little if the $45$ sit just under the line, so the
distance matters. The closest of them is $[18,4]$ heating,
$\Te = \SI{2.33}{\electronvolt}$, $\nel = \SI{5.18e13}{\per\cubic\centi\metre}$,
at $\tauesc/\tauslow = 0.636$ for $L = \SI{20}{\centi\metre}$: its escape time
would have to rise by a factor $1.57$ for the parcel to close. The median of the
$45$ would have to rise by a factor $52$. One pair is marginal and the rest are
not.

Three things were varied to see whether the conclusion depends on a choice.
Tripling $\langle\sigma v\rangle_{\mathrm{CX}}$, which lengthens every diffusive
escape time in proportion, moves the count from $0$ to $3$ of $45$; dividing it
by three leaves it at $0$. Replacing $\sqrt{2k\Tn/m}$ by the Maxwellian mean
speed $\sqrt{8k\Tn/\pi m}$ shortens every escape time and leaves the count at
$0$. Doubling $L$ from $10$ to \SI{20}{\centi\metre} leaves it at $0$. The
finding is not an artifact of any of the three.

\paragraph{The selection is structural, not incidental.} A pair enters the
census because $\tauslow$ is long, since the single-slow-mode estimate of
Eq.~\eqref{eq:elm_bound} rises with $\tauslow$ at fixed $\taudrive$. A long
$\tauslow$ means slow ionisation out of the ground state. Slow ionisation is
exactly the condition under which neutral transport, whose timescale does not
depend on $\tauslow$ at all, sets the ground-state density instead of the local
atomic physics. The census therefore selects for the regime in which its own
assumption is weakest, and it does so by construction rather than by accident.
That is why testing the assumption at the cold corner gave the reassuring answer
and testing it inside the defended scope does not.

\subsection{What this reaches, and what it does not}
\label{sec:transport_partition}

The finding above is about magnitudes; it does not touch the algebra.

\textbf{Untouched by transport.} The two-channel split of
Eq.~\eqref{eq:two_supplies}, exact to a relative residual of
$3.075\times10^{-14}$; the unit-width logistic form of the ground-fed fraction;
the bound $\max|f_3-f_4| = \tanh(|\Delta|/4)$ and its corollary
$|\mathrm{d}\ln\Robs/\mathrm{d}\ln\bdep| < 1$; the sensitivity map
$\overline{S}$; and the accuracy of the excited-state closure,
$6.73\times10^{-9}$ at the worst point. These are properties of the operator
$\Lmat(\Te,\nel)$ at a fixed operating point, or theorems about its structure.
$\overline{S}$ says how the line ratio responds to whatever the reservoir does;
it is silent on what made the reservoir do it. Transport changes the reservoir's
history and cannot reach any of them.

\textbf{Conditional on the closed parcel.} The reservoir gain $G$, which is
a secant along the CRE curve and therefore assumes the ground
state responds through collisional--radiative balance alone; $\epsplat$, which
is built from it; the breakdown census; and every time-averaged magnitude in
Chapter~\ref{ch:results}. If recycling renews $\ngs$ on a timescale shorter than
$\tauslow$, $G$ is not the gain the plasma has. These numbers are not upper
estimates. A fast-exchanging reservoir with a steady source freezes the
closed-parcel error rather than removing it, and a rising source enlarges it: over
the defended set the true error exceeds the closed-parcel plateau at $165$ of the
$180$ (pair, $L$, source-multiplier) combinations tested
(\texttt{verify\_open\_reservoir.py}). The closed-parcel value is a reference
magnitude conditional on the reservoir's history, not a bound on it.

The thesis therefore leads with the structure and the bound and reports the
magnitudes as conditional; Section~\ref{sec:transport_selection} adds that the
condition is measured to fail at the points that matter, by one to two orders
of magnitude.

\subsection{The claim, restated as the conditional it is}

The honest form is therefore not a caveat but a conditional, and it should be
written this way wherever the number appears:

\begin{quote}
\emph{If} the ground-state neutral density is frozen for the duration of the
event, the error is as computed.
\end{quote}

That conditional is close to exactly true for the excited manifold, which
equilibrates in nanoseconds, and it is the standard implicit assumption of every
$0$-D collisional--radiative table in use, the object this thesis set out to
test. But it is an assumption about the plasma, not about the atoms, and no
eigenvalue analysis inside a $0$-D matrix can validate it: the two timescales
are measured to differ by four orders of magnitude and assumed not to interact.
What would settle it is a two-region or transport-coupled calculation, an edge
code such as SOLPS (Scrape-Off Layer Plasma Simulation) supplying $\ngs$ from a
recycling influx; that is outside the scope of this thesis and is named in
Chapter~\ref{ch:conclusions} among the further work.

%-----------------------------------------------------------------------------
\section{No molecular channels, in a gas that is mostly molecules}
\label{sec:molecules}
% QUESTION: the model contains only atomic hydrogen -- at the conditions where
% the error is largest, is that a defensible description of the gas?

\subsection{The question, and why a referee asks it first}

Every state in Chapter~\ref{ch:atoms}'s state space is a level of a single
hydrogen \emph{atom}, plus a reservoir of bare protons. There is no
$\mathrm{H}_2$ (or $\mathrm{D}_2$), no molecular ion $\mathrm{D}_2^+$, no
negative ion $\mathrm{D}^-$, and no reactions between them. A search of the rate
assembly for any of these returns nothing at all.

\subsection{The picture: a third feed stream}

Chapter~\ref{ch:theory} fed every excited level from two directions, up from
the ground state by electron collisions and down from the ion continuum by
recombination, and the error mechanism is the difference in how the $n=3$ and
$n=4$ shells split their supply between those two feed streams.

Molecules open a third stream. In a cold, dense, mostly neutral gas, hydrogen
atoms form $\mathrm{H}_2$; the molecules are ionised to $\mathrm{H}_2^+$, or
exchange charge, and the resulting molecular ions recombine with electrons and
break apart, leaving one of the fragments in an \emph{excited} atomic level. The
whole family of such routes is called MAR, molecular activated recombination.
The important feature for this thesis is not that MAR exists but where it
delivers: it populates $n=3$ directly, without the atom ever having passed
through the ground-state reservoir whose staleness is the entire mechanism here.
A supply channel that bypasses the reservoir does not inherit the reservoir's
lag.

\subsection{How much gas is involved}

At the worst point of the map the model's own equilibrium solution is
\SI{96.6}{\percent} neutral: $\ngs = 1.47\times10^{15}$ against
$\nel = 5.18\times10^{13}\,\si{\per\cubic\centi\metre}$. A hydrogen gas at
\SI{1}{\electronvolt} that is \SI{96.6}{\percent} neutral is a gas in which
molecules form readily. The atomic description is not being applied at the edge
of its validity there; it is being applied outside it.

The size of what is being left out is documented experimentally.
\emph{Detachment} is the operating regime in which the plasma cools enough that
the ion flux to the target collapses, and the regime this diagnostic is chiefly
used to identify. In spectroscopic inversions of divertor Balmer emission at its
onset, molecularly-induced emission is found to account for up to
\SIrange{60}{70}{\percent} of the $D_\alpha$ light and
\SIrange{10}{20}{\percent} of $D_\gamma$~\cite{Wijkamp2023,Verhaegh2021}. An
$H_\alpha/H_\beta$ ratio computed from atomic processes alone is not the ratio
such an instrument records.

\subsection{What survives it, and what does not}

Two consequences follow, and they point in opposite directions, so both must be
stated.

The \emph{numerical} cold-edge results are atomic-model results and nothing
more. Adding a third supply channel that feeds $n=3$ preferentially would change
$\ffrac{3}-\ffrac{4}$, and could plausibly change it a great deal; this work
offers no direct calculation of how much. A bound can nonetheless be
constructed, from the emissivity fractions cited in this section. For a given
upper level the molecular fraction of the emissivity equals the molecular
fraction of the population, because both channels emit through the same
transition probability, and molecular-assisted recombination bypasses the ground
state entirely, so it dilutes the ground-fed fraction as
$\ffrac{p} \to \ffrac{p}(1-\phi_p)$. Taking $\phi_3 = 0.60$ and $\phi_4 = 0.30$
from the reported \SIrange{60}{70}{\percent} of $\mathrm{D}\alpha$ and
\SIrange{10}{20}{\percent} of $\mathrm{D}\gamma$, the difference of ground-fed
fractions at the cold corner falls from $0.4866$ to $0.1088$ and at the density
crest from $0.4580$ to $0.0998$, and the secant sensitivity $\overline{S}$
falls with it, from $0.4822$ to $0.0903$ and from $0.4552$ to $0.0903$.

The plateau error does not fall in the same proportion, because
$\epsplat = |\mathrm{e}^{\overline{S}\Lambda}-1|$ is not linear in
$\overline{S}$; treating it as linear understates the reduction. Imposing the
same fractions on the model's own operator, with a non-negative source
depositing into $n=3$ and $n=4$ so that $(\phi_3,\phi_4)$ is reached exactly at
the plateau, the cold-corner plateau error falls from \SI{38.7}{\percent} to
\SI{6.3}{\percent}, a \textbf{factor of $6.1$}; at the upper end of the
reported ranges, $(\phi_3,\phi_4) = (0.70,0.45)$, it falls to \SI{4.1}{\percent},
a factor of $9.4$. At the crest the same two targets give factors of $5.4$ and
$8.1$, and at the benchmark $3.2$ and $4.4$
(\texttt{validation/molecular\_channel/}).

Two features of that calculation are worth stating, because neither is visible
in the dilution formula. The two fractions cannot be chosen independently: a
source large enough to make $\phi_3 = 0.60$ already drives more than $0.15$ of
$n=4$ through collisional $3\to4$ transfer, so $(\phi_3,\phi_4) = (0.60,0.15)$
requires a negative $n=4$ source and is not realisable on this operator. And
the reduction is not monotone in $\phi_3$. The difference
$\ffrac{3}(1-\phi_3) - \ffrac{4}(1-\phi_4)$ passes through zero near
$\phi_3 = 0.74$ at $\phi_4 = 0.30$, beyond which $n=4$ is the more ground-fed
of the two and the error grows again with $\overline{S}$ carrying the opposite
sign. At the crest, where that zero is reachable with non-negative sources, a
two-source dilution at $\phi_4 = 0.30$ brings the plateau error down to
\SI{0.35}{\percent} at $\phi_3 = 0.75$, past the zero at $\phi_3 = 0.7355$; a
single $n=3$ source there gives \SI{0.30}{\percent} at $\phi_3 = 0.70$ and
\SI{1.1}{\percent} again at $\phi_3 = 0.80$, with $\overline{S}$ negative. A
molecular channel strong enough cancels the diagnostic's sensitivity to the
reservoir rather than merely reducing it, and past the cancellation the error
returns.

The bound is crude, it inherits every uncertainty in the quoted fractions, and
it applies only where those fractions do.

The \emph{structural} results, the list of Section~\ref{sec:transport_partition},
are theorems about whatever matrix $\Lmat$ is supplied and hold for one
containing molecular states; what such a matrix changes is the numerical value
of the ingredients, not the form of the result.

What is \emph{not} available is a calculation with molecules in it: the factors
above are the model's own operator driven by an imposed source of the published
strength, not a solution of a matrix containing molecular states. This is an
atomic model, its cold-edge numbers are atomic-model numbers, and the bound
above is an estimate of how wrong they are rather than a correction to them.

%-----------------------------------------------------------------------------
\section{Where the ridge sits depends on what the model assumed about neutrals}
\label{sec:ridge_conditional}
% QUESTION: is the density at which the error peaks a property of hydrogen, or
% a property of this model's own ionisation balance?

\subsection{The question}

Section~\ref{sec:mechanism} found an interior maximum, a ridge, in the
error map at $\nel \approx 2\times10^{13}\,\si{\per\cubic\centi\metre}$, and
showed that it moves with the shell pair as Griem's criterion for local
thermodynamic equilibrium (LTE) predicts. The mechanism is robust; three
independent tests agreed. But a mechanism and a location are different claims,
and only the first of them has been established.

\subsection{The picture: the ridge is where the two feeds are balanced}

The ridge exists because the $n=3$ and $n=4$ shells switch from
recombination-fed to ground-fed at different reservoir strengths, so their
difference vanishes when either feed dominates and peaks in between. The
controlling variable is therefore not the electron density directly; it is
\emph{how much neutral gas there is to feed from}. The electron density enters
only because, in this model, it determines the neutral density through the CR
ionisation balance.

That is the conditionality. Change what the model assumes about neutrals and the
ridge moves, whatever the electron density does.

\subsection{The measurement}

Within the two-channel split the ground-fed fraction is exactly
$\ffrac{p} = a_1 n_g/(a_0 + a_1 n_g)$, where $n_g$ is the ground-state density
feeding the excitation channel and $a_0$, $a_1$ are the network responses of
Section~\ref{sec:R_depends_on_b}. Scaling $n_g$ is therefore an exact operation
on the ridge, not an approximation to one. Doing it, at
$\Te = \SI{2.02}{\electronvolt}$:

\begin{center}
\begin{tabular}{lr}
  \toprule
  $\ngs$ scaling & density at which $|\ffrac{3}-\ffrac{4}|$ peaks \\
  \midrule
  $\times 0.01$ & $\le 10^{12}\,\si{\per\cubic\centi\metre}$ (off grid) \\
  $\times 0.1$  & $2.28\times10^{12}\,\si{\per\cubic\centi\metre}$ \\
  $\times 1$ (the model's own CRE value) & $1.68\times10^{13}\,\si{\per\cubic\centi\metre}$ \\
  $\times 10$   & $2.48\times10^{14}\,\si{\per\cubic\centi\metre}$ \\
  $\times 100$  & $\ge 10^{15}\,\si{\per\cubic\centi\metre}$ (off grid) \\
  \bottomrule
\end{tabular}
\end{center}

\noindent \textbf{One decade in the assumed neutral density moves the ridge by
roughly one decade in electron density.} The response is close to
proportional, and the whole three-decade span of the density grid is traversed
by a two-decade change in $\ngs$.

The agreement with Griem's criterion at unit scaling, $1.93\times10^{13}$
measured against $2.05\times10^{13}\,\si{\per\cubic\centi\metre}$ predicted, is
therefore an agreement evaluated at \emph{this model's} collisional--radiative
ionisation balance, which contains no transport and no recycling. In a real
divertor the neutral density is set by recycling, which is
Section~\ref{sec:transport}'s objection arriving in a different place. The
sentence ``the ridge sits at
$\nel \approx 2\times10^{13}\,\si{\per\cubic\centi\metre}$'' is a statement about
an ionisation balance, not about a divertor, and should be quoted with that
condition attached or not quoted.

The ridge is also one line pair's and not the Balmer series': the $(3,5)$ pair
of $H_\alpha/H_\gamma$ has its worst density at
$7.2\times10^{12}\,\si{\per\cubic\centi\metre}$, a factor $2.68$ lower, and the
form of words this imposes is stated in Section~\ref{sec:error_maps}.

%-----------------------------------------------------------------------------
\section{The model does not identify the crest with detachment}
\label{sec:not_detachment}
% QUESTION: can the density at which the ridge sits be identified with a named
% operating regime of a tokamak divertor?

\subsection{The question, and the temptation}

The identification is tempting. The ridge sits where ionising and recombining
supplies compete; detachment, the regime ITER must operate in to survive its
exhaust power and the transition this diagnostic is chiefly used to identify
(Section~\ref{sec:molecules}), is where ionising and recombining plasma compete.
The words match.

\subsection{The arithmetic}

The numbers do not. Section~\ref{sec:withdrawn} withdrew the identification:
the ridge, at $1.93\times10^{19}\,\si{\per\cubic\metre}$, sits a factor $5.2$
below the detached-divertor band of published edge
modelling~\cite{Guillemaut2011} and a factor $1.6$ below the modelled upstream
separatrix density~\cite{Stangeby2023,Pitts2019}, in a range neither source
characterises.

The comparison band deserves a word, because a factor of $1.6$ is not a large
margin to be arguing over. Read from the primary source, the band is
\SIrange{3e19}{6e19}{\per\cubic\metre} across the $53$ ITER SOLPS-4.3 cases
\cite[Sec.~3.1.1]{Pitts2019}, and it is the density \emph{at the outer
midplane}, which that paper states explicitly where it plots the dependence
\cite[Sec.~3.2.2]{Pitts2019}. So the quantity is settled: outer midplane, not
X-point. What is not settled by a single factor is where in the band to
measure from. The ridge sits $1.6$ times below the \emph{bottom} of it and
$3.1$ times below the top, and the honest statement is the range rather than
either end. On the same flux tube the outboard-midplane and X-point densities
differ by roughly \SI{50}{\percent}, which is comparable to the width of the
band itself. The conclusion does not turn on any of this, since the ridge is
below the band on every reading, but the factor should not be quoted more
precisely than the band it is measured against.
The repository recorded \cite{Stangeby2023} as ``Nucl.\ Fusion 62 (2022)'' in
\texttt{src/validation/grid.py}; the DOI was right, the volume and year wrong.
It is Nuclear Fusion \textbf{63}(1) 016017 (2023), companion part A at 63(1)
016016 \cite{Stangeby2023a}; \cite{Lore2022} in the same block of that file
genuinely is volume 62.

\subsection{What that leaves}

The honest statement is uncomfortable and should be made anyway. \textbf{The
structure this work found is real and reproducible, and it sits at a density
this thesis cannot identify with any named region of a tokamak.} It is not
detachment; it is not the upstream separatrix; and calling it either would be a
claim about edge transport that a $0$-D atomic model cannot support. The ridge
is a property of how two shells of hydrogen divide their supply. It needs no
divertor geometry to produce it, and it acquires none by being described in one.

The ridge and the detached band do not overlap: restricted to
$\Te \ge \SI{2}{\electronvolt}$ and $\nel \ge 10^{14}\,\si{\per\cubic\centi\metre}$
the census is $0$ of $108$ with a worst estimate of \SI{7.2}{\percent}
(Table~\ref{tab:trapping_census}), so where the diagnostic is most used the
transient error is real but modest.

%-----------------------------------------------------------------------------
\section{The open boundary: an attack that was expected to be fatal}
\label{sec:open_system}
% QUESTION: does holding the ion density fixed manufacture the slow clock that
% the whole result rests on?

\subsection{The question}

Section~\ref{sec:rate_equation} was explicit that the ion density $\nion$ is
held fixed: the recombination feed $\Svec\nion$ enters as a constant source term
a feed stream into a CSTR in the reactor picture, rather than as a
matrix element, and the $43$-state system has no $44$th state for ionised atoms
to accumulate in. Atoms ionise and vanish from the accounting. The slow
eigenvalue $\lambda_0$ is therefore the CR ionisation rate of ground-state
hydrogen against a prescribed reservoir, not the equilibration rate of a closed
system.

\subsection{The falsifier, and what happened}

If the long $\tauslow$ values were an artifact of the open boundary, so would be
the claim that the error outlasts an ELM. The falsifier named in
Section~\ref{sec:attacks} before the test was run is a closure factor of $2300$
at $[0,4]$, the value of $\tauslow/\taudrive$ there. Closing the manifold with
a $44$th ion state gives a closure factor with minimum $1.00$, median $1.00$
and maximum $41.4$ over the grid, moves the ELM census from $202$ to $200$ of
$680$ and the worst case by one part in $325$, and leaves the count above
\SI{2}{\electronvolt} at $45$ before and $45$ after (Table~\ref{tab:closure}).
The factor is large only where $\tauslow$ already exceeds the event duration by
three orders of magnitude and is unity wherever $\tauslow$ approaches
$\taudrive$, so where the boundary condition matters the outcome does not
depend on it. The measured factor at $[0,4]$ is $15$ against the $2300$
required, a cold-corner margin; inside the defended range, at the ridge, it is
$1.01$ against $20$, and at the benchmark $\tauslow$ moves by one part in
$10^{3}$. The attack was survived everywhere it was pressed.

\subsection{The caveat that must not be softened}

The closure adds an ion reservoir. It does \emph{not} close the electron
balance, since $\nel$ is held fixed while the ion population moves, and it adds
no transport whatever. It is a test of the boundary condition on the atomic
manifold, and nothing more. Section~\ref{sec:transport} is the objection this
test does not answer, and no amount of closing the atomic system will answer it.

The two index labels for this point, $[1,4]$ in Table~\ref{tab:closure} and
$[0,4]$ in the census, are the post- and pre-step operators of one heating
step: the census rows carry the pre-step index, every quoted timescale belongs
to the post-step operator that governs the relaxation (Chapter~\ref{ch:results}),
and the two must not be interchanged. Diagonalising each gives
$\tauslow = \SI{0.421012}{\second}$ at $\Lmat[0,4]$ and
\SI{0.233239}{\second} at $\Lmat[1,4]$, so the quoted \SI{0.23324}{\second}
belongs to $\Lmat[1,4]$.

%-----------------------------------------------------------------------------
\section{An external benchmark this model fails, and what the failure measures}
\label{sec:r1_deficit}
% QUESTION: where does the comparison against published atomic data break down,
% and does the failure touch the mechanism?

\subsection{The question}

Chapter~\ref{ch:validation} presented the comparison against Fujimoto's
published population coefficients as the strongest external check this model
passes. It is, for one of the two coefficients. The other disagrees at low
principal quantum number, and this section reports what that disagreement
measures. The short answer is that it is unexplained, that it is not the
$\ell$~closure, and that it is direct evidence about the coefficients
Chapter~\ref{ch:results} uses.

\subsection{The picture: two coefficients, one for each feed}

Chapter~\ref{ch:theory} split every excited population into two pieces
corresponding to the two supply channels. The population coefficient $r_0(p)$
measures how strongly level $p$ is fed from the ion continuum by recombination;
$r_1(p)$ measures how strongly it is fed from the ground state by excitation.
The published tabulation gives both, computed independently, decades ago, by a
different method. Agreement is a real test; disagreement is a real failure.

\subsection{The disagreement, and what it measures}

The recombination-fed coefficients agree. At
$\nel = 10^{12}\,\si{\per\cubic\centi\metre}$, $r_0$ matches to
\SI{1.3}{\percent} at $p=2$ and to \SIrange{12}{14}{\percent} at $p=3$ and
$p=4$. The ground-fed coefficients at low principal quantum number do not: this
model's $r_1$ sits a factor $8.3$ low at $p=3$ and a factor $4.4$ low at $p=4$.

Those are the two shells whose difference is the mechanism of
Chapter~\ref{ch:results}, so the disagreement had to be understood rather than
noted.

Two candidate explanations were eliminated before the source was consulted. If
the excitation rates feeding the ground-fed channel were wrong, the deficit
would follow, but those rates were benchmarked against
R-matrix-with-pseudostates calculations and the dominant $\Delta\ell = 1$
channels agree at a median of $-\SI{5.0}{\percent}$ over $120$ comparisons. And
a normalisation error is excluded by $r_0$ itself: $r_0$ and $r_1$ carry the
same $p$-dependent factor $Z(p)$, so an error large enough to move $r_1$ by a
factor of $8$ would have moved $r_0$ with it.

\subsection{The source's closure, imposed and tested}

Table~4.1(b) is indexed by $p = 2, 3, 4, 5, 7, 10, 15$. Those are principal
quantum numbers, and the table carries no $\ell$ label anywhere. The published
coefficients are shell quantities computed with the $\ell$~substructure
bundled.

This model resolves $\ell$ below $n=9$, so the closures differ. An earlier
draft of this section treated that difference as the explanation, on the
evidence that recomputing the model with proton $\ell$-mixing removed moves
$r_1(2)$ by a factor $118$ and $r_1(3)$ by a factor $3.35$ at
$\nel = 10^{12}\,\si{\per\cubic\centi\metre}$. That argument does not survive
its own test. Removing $\ell$-mixing is the limit in which the process is
absent; the tabulation assumes the limit in which it is infinitely fast, and
those are opposite ends rather than a bracket. Imposing the source's own closure
directly, by bundling this model's operator under statistical $\ell$~weights,
moves $r_1(3)$ by \SI{0.42}{\percent} and makes the $p=2$ agreement worse
(\texttt{verify\_fujimoto\_bundle.py}, stamped in
\texttt{validation/fujimoto\_bundle/}).

The disagreement therefore does \emph{not} measure the $\ell$~closure, and this
thesis does not explain it. An earlier draft of this chapter said the deficit
did not reach $\ffrac{3}-\ffrac{4}$ because that difference never passes
through a bundled quantity; that was wrong, and bundling is not the protection.
Section~\ref{sec:fujimoto} makes $a_p = r_1(p)\,\Zsb{p}/\Zsb{1}$, so a deficit
in $r_1(3)$ and $r_1(4)$ is a deficit in $a_3$ and $a_4$ themselves, and the
only protection is that $\Delta$ and the cap depend on the ratio $a_3/a_4$, in
which a deficit common to both shells cancels. It is not common at
\SI{e12}{\per\cubic\centi\metre}, $8.3$ against $4.4$, and nearly so at
\SI{e15}{\per\cubic\centi\metre}, $1.8$ against $1.9$. Dividing $a_3$ and
$a_4$ by those ratios at every grid point
(\texttt{validation/fujimoto\_r1\_consequence/}) moves $\Delta$ by $+0.64$ and
$-0.08$; over the $280$ points above \SI{2}{\electronvolt} the cap moves by
\SI{+29}{\percent} and \SI{-4}{\percent} at the median, and the sensitivity at
the operating point by \SI{+30}{\percent} and \SI{+11}{\percent} at the median,
with ranges of $-58$ to \SI{+524}{\percent} and $-38$ to \SI{+68}{\percent}.
The disagreement is direct evidence about the coefficients
Chapter~\ref{ch:results} uses; what is open is whether the deficit lies in this
model or in the tabulation, and the stake is stated in
Section~\ref{sec:fujimoto}. The open status is carried in
Section~\ref{sec:scope_open}.

\subsection{Two limitations of the benchmark, stated precisely}

Table~4.1(b) is computed at $T_\mathrm{e} = \SI{1.28e5}{\kelvin}$, which is
\SI{11.03}{\electronvolt}. The grid used here reaches \SI{10}{\electronvolt},
so the comparison is made at the grid edge and \SI{9.3}{\percent} below the
tabulated temperature.

There is no second usable case. The companion Table~4.1(a) is at
$T_\mathrm{e} = \SI{e3}{\kelvin}$, which is \SI{0.0862}{\electronvolt}, two
orders of magnitude below this grid. The external check is therefore one
temperature at one edge, and it should not be described as more than that.

\subsection{What the benchmark establishes, and what it does not}

What the benchmark does establish is the recombination-fed channel, which
agrees, and that is worth stating as a pass rather than leaving implicit beside
a disagreement. What it does not establish is the absolute normalisation of the
ground-fed channel against an independent calculation, because no tabulation
resolved in $\ell$ was available to compare against.

The sensitivity of the density maximum to the ground-fed supply remains a
separate and open question, and Section~\ref{sec:ridge_conditional} treats it
on its own terms: the location responds roughly decade-for-decade to the assumed
neutral density, which is a stronger reason to quote it as a range than
anything the benchmark shows.

%-----------------------------------------------------------------------------
\section{The top of the level ladder, and what the bundling costs}
\label{sec:nmax}
% QUESTION: the model stops at n=15 and bundles the levels above n=8 -- what
% rests on the part that is approximated?

\subsection{The question}

A hydrogen atom has infinitely many bound levels, crowding together towards the
ionisation limit. A computation cannot have infinitely many states, so a
boundary was chosen: this model tracks levels up to $n_{\max} = 15$ and
discards the rest. Above $n = 8$ it also stops resolving orbital angular
momentum $\ell$, treating each shell as a single bundled state populated
statistically. Both are choices, and this section states what they cost.

\subsection{Three consequences, in decreasing order of severity}

\paragraph{The terminal shell is over-populated.} Truncation removes from the
terminal shell's balance the collisional excitation out of it to $n > n_{\max}$,
a way out, and the de-excitation and cascade into it from those levels; an
earlier draft said the shell had ``nothing above it to cascade out to'', which
names the wrong process, since cascade is downward. Against published values
this model's ground-fed coefficient at $n = 15$ is high by a factor of $4.9$ to
$6.2$, while at $n = 10$ it agrees to \SIrange{4}{20}{\percent}; the budget is
in Section~\ref{sec:convergence}. \emph{No result in this thesis rests on the
top shell}: the observable is built from $n = 3$ and $n = 4$, both fully
$\ell$-resolved and far from the truncation.

\paragraph{The bundling nonetheless enters the result indirectly.} The
recombination flux cascading down from the bundled shells feeds the
recombination-fed channel of both $n=3$ and $n=4$, one of the two channels the
mechanism compares, which is why Section~\ref{sec:ridge_alternatives} lists
truncation among the alternative explanations for the ridge that this work does
not exclude.

\paragraph{There is no $\ell$-mixing inside the bundled block.} The
proton-impact $\ell$-mixing that Section~\ref{sec:lmixing} showed to be
essential for the $n = 2$ manifold is present only for $n = 2$--$8$; the
$\ell$-mixing rate array is identically zero above that. Within the bundled
shells the $\ell$-distribution is assumed statistical rather than computed. At
high $n$ that assumption is far more defensible than it would be at low $n$,
because the sublevels are nearly degenerate and every collisional process mixes
them, but it is assumed, not measured, and this thesis has not measured it.

\subsection{The check, run}
\label{sec:bundling_check}

The script that quantifies this was written and, until now, never executed.
Running it produced a result and also produced three defects in the script
itself, which are reported here because two of them would have made its verdict
worthless.

\paragraph{What the script could not do.} It promised to test the bundled
shells against the pipeline's own $\ell$-mixing array, and that is impossible:
the array is identically zero for states $36$ to $42$, in both directions,
because $\ell$-mixing was computed for the resolved block alone. The file
contains nothing to test. The comparison has to be made against the rate
\emph{function} rather than the rate table, and the function is the one
\texttt{compute\_lmix.py} used to build the resolved block, so it is the model's
own rate rather than a second one. Calling it at $n = 2$ to $8$ reproduces the
stored array to $1.0000$, which is a wiring check and not a physical
confirmation; it establishes only that the rate carried up to the bundled
shells is the same function with the same arguments.

\paragraph{The two defects that would have inverted the answer.} The script
carried its own Pengelly--Seaton expression, independent of the pipeline's,
which disagrees with the rate the model actually uses by a factor of $95$ at
$n = 2$ rising to $5.6\times10^{3}$ at $n = 8$, growing roughly as $n^{2}$. Used
as written, it would have reported the mixing rate two to three orders of
magnitude below its true value and could have returned a verdict of marginal
where the answer is not close. It is deleted rather than corrected: a
verification script has no business maintaining a second implementation of a
rate the pipeline already computes. The script also fell back on a synthesised
logarithmic temperature grid of its own when the pipeline's grid files were
absent,
printing a warning and continuing, which is the failure mode described in
Chapter~\ref{ch:validation}: every ratio would have carried a correct-looking
label attached to conditions the model was never evaluated at. It now raises.

\paragraph{The result.} The criterion is that the $\ell$-mixing rate out of the
$np$ substate, the substate the Lyman channel empties fastest, must exceed that
substate's total radiative rate. Over the seven bundled shells and all $400$
grid points, the ratio is above the well-mixed threshold of $10$ at every one of
the $2800$ combinations. The worst case is $n = 9$ at
$\Te = \SI{1}{\electronvolt}$, $\nel = \SI{e12}{\per\cubic\centi\metre}$, where
the mixing rate is $\SI{1.155e10}{\per\second}$ against
$A(9p) = \SI{7.51e6}{\per\second}$, a margin of $1538$, that is, a factor $154$
clear of the threshold (\texttt{verify\_bundling\_psm20.py},
\texttt{outputs/bundling/bundling\_report.txt}).

Two inputs to that comparison were checked rather than assumed. The Einstein
coefficients $A(np\to1s)$ for $n = 9$ to $15$ are hardcoded in the script,
because the model's own radiative file resolves $\ell$ only to $n = 8$ and has
nothing to read above it. Since $A(np\to1s)\,n^{3}$ tends to a constant for
hydrogen, the model's resolved values and the hardcoded table must lie on one
curve: they do, with a step of \SI{0.43}{\percent} across the boundary at
$n = 8$ to $9$ and a monotone approach to $4.13\times10^{9}$ at $n = 15$. And
the $\ell$-mixing coefficient carries a Debye cutoff frozen at
$\nel = \SI{e14}{\per\cubic\centi\metre}$, because the array has no density
axis. That freezing is not small: re-evaluating the cutoff at the local density
changes the $n = 15$ rate by a factor $20$ upward at
\SI{e12}{\per\cubic\centi\metre} and a factor $30$ downward at
\SI{e15}{\per\cubic\centi\metre}. Against the radiative criterion it does not
move the verdict, because the worst point sits at the lowest density, where the
local cutoff raises the rate: the worst ratio is $1538$ with the freezing and
$3.26\times10^{4}$ without it. At the highest densities the same correction
lowers the rate, and that matters for the criterion that follows.

\paragraph{What this closes and what it does not.} The criterion above compares
mixing with radiative decay, and radiative decay is not how a high shell
empties at these densities. The total loss out of a bundled state, $-L_{kk}$,
is \SIrange{89.5}{99.1}{\percent} collisional transfer to other shells,
\SIrange{0.9}{10.5}{\percent} ionisation and at most \SI{0.19}{\percent}
radiative (\texttt{verify\_bundling\_total\_loss.py},
\texttt{validation/bundling\_total\_loss/}). Against that total the margin is
not three orders of magnitude. With the frozen cutoff the model is built with,
mixing exceeds total escape at every one of the $2800$ combinations, by $26$
at $n=9$ falling to $2.85$ at $n=14$, the minimum, all at
\SI{1}{\electronvolt} and \SI{e12}{\per\cubic\centi\metre}, and $288$
combinations sit below $10$. With the cutoff evaluated at the local density
the ratio falls below unity at $121$ of the $2800$, in every bundled shell, at
the two highest densities and $\Te \le \SI{2.95}{\electronvolt}$; at
\SI{1}{\electronvolt} and \SI{e15}{\per\cubic\centi\metre} it runs from $0.86$
at $n=9$ to $0.09$ at $n=14$. The comment in \texttt{compute\_lmix.py} that the
cutoff factor varies by about \SI{30}{\percent} across the grid understates the
variation already at $n=2$, where the local-to-frozen ratio of $F$ runs from
$1.96$ at \SI{e12}{\per\cubic\centi\metre} to $0.53$ at
\SI{e15}{\per\cubic\centi\metre} at \SI{1}{\electronvolt}
(\texttt{validation/bundling\_total\_loss/}, RESULT~4), and fails outright
above $n=9$, where $U_m > 1$ and $F$ falls as $\nel^{-3/2}$; it is reported
here and left in place. What survives is narrower than an
earlier draft claimed. Where mixing wins, the $\ell$-distribution is
statistical whatever the feeds do. Where collisional transfer wins, at high
density, the distribution is set by the feeds, and $n$-changing collisions
between adjacent shells, which carry \SIrange{71}{95}{\percent} of the
collisional loss, are close to $\ell$-blind at high $n$, so the statistical
assumption is plausible there for a different reason, one this thesis has not
tested. It does not close the terminal-shell over-population above, which is a
different defect with a different cause, and it does not make the bundled
block's $\ell$-distribution a computed quantity. It remains assumed, now with
two numbers attached and a stated region where the first of them does not
license it.

%-----------------------------------------------------------------------------
\section{Remaining assumptions, and the sensitivities that are known}
\label{sec:other_assumptions}
% QUESTION: what else was assumed, and which of those assumptions has actually
% been tested?

Four assumptions remain. Each is stated with what is known about it, rather than
listed as a disclaimer, and two of the four have been tested.

\paragraph{Maxwellian electrons: untested.} Every rate coefficient in
Chapter~\ref{ch:atoms} is an average of a cross section over a Maxwell--Boltzmann
electron energy distribution. Ionisation and excitation are driven by the
high-energy tail of that distribution, which holds a small minority of the
electrons. A \SI{1}{\electronvolt} plasma next to a recycling surface has no
guarantee of a Maxwellian tail, and a depleted or enhanced tail would change the
rates that set both clocks. Not tested here.

\paragraph{$T_i = \Te$ in the $\ell$-mixing rates: untested.} The
proton-impact $\ell$-mixing rates of Section~\ref{sec:lmixing} are evaluated at
an ion temperature equal to the electron temperature. In a detaching divertor
the ions are frequently colder than the electrons. Not tested here.

\paragraph{The step is instantaneous: tested, and safe for a reason that is
easy to get wrong.} The error maps impose a sudden temperature step, whereas a
real ELM has a finite rise time. The tempting ratio
$\taurel/\taudrive \approx 10^{-5}$ is about the wrong quantity, the fast
excited-manifold step error (Section~\ref{sec:persistence}); the plateau error
is governed by the slow ground-state response, for which the relevant ratio is
$\tauslow/\taudrive$. At the worst point of the grid that
ratio is $2332$, but that point sits at \SI{1}{\electronvolt} and the figure is
not the margin this result runs on. Over the $448$ defended pairs the ratio has
median $0.815$, and over the $45$ that break down it runs from $1.95$ to $288$,
median $11.2$ (\texttt{validation/divertor\_map/}). The step idealisation is
therefore safe by one to two orders of magnitude where a plateau is reported,
not by three. The governing quantity is in any case
$D_{\rm e} = \tauslow/t_{\rm ramp}$ rather than $\tauslow/\taudrive$: a linear
ramp reaches $D_{\rm e}(1-\mathrm{e}^{-1/D_{\rm e}})$ of the step's plateau
(\texttt{validation/ramp\_plateau/}), so any ramp with
$\taurel \ll t_{\rm ramp} \ll \tauslow$ reaches the same plateau and the
idealisation \emph{is} safe. It is safe because of the slow ratio, and the fast
ratio must never be quoted for or against this result.

\paragraph{The step is in temperature only: tested, and the correction goes the
wrong way.} A real ELM raises $\nel$ as well as $\Te$. The joint map over $4968$
cases (Section~\ref{sec:joint_step}) leaves the exact decomposition unchanged,
with the gain generalising to an additive gradient accurate to
\SI{0.2}{\percent} at the median; over $\Te\ge\SI{2}{\electronvolt}$ the
breakdown fraction moves from $10.0$ to \SI{11.3}{\percent} when the density
steps up one index and the worst case from $0.1748$ to $0.2206$. The earlier
benchmark spot check, a bound falling by a factor $4$, had the right number and
the wrong sign for the grid, because the benchmark and the worst point lie on
opposite sides of the density crest; the fixed-density map understates the
worst case.

\medskip
Two sensitivities of the reported census belong in the same place, because both
are properties of a reporting choice rather than of the physics, and both would
otherwise be mistaken for measurements.

The step size is one interval of the temperature grid, which displaces the
reservoir by $|\ln\xscale|$ between $0.124$ and $0.682$ across the grid. A
four-interval step gives $0.516$ to $2.563$ and raises $\epsplat$ at the ridge
from \SI{18.1}{\percent} to \SI{81.1}{\percent}. On step size alone, therefore,
the reported numbers \emph{understate} the error rather than overstate it.

The \SI{10}{\percent} breakdown threshold is doing real work: at \SI{5}{\percent},
\SI{10}{\percent} and \SI{20}{\percent} the count over the same $680$ points is
$348$, $202$ and $61$ (Section~\ref{sec:persistence}), so the count is not robust
to the threshold, the worst case is, and ``$N$ points break down'' must carry both.

%-----------------------------------------------------------------------------
\section{Scope: five questions this thesis leaves open}
\label{sec:scope_open}
% QUESTION: which questions are open, why can none of them be closed from
% inside a 0-D atomic model, and what would close each?

The objections computed in this chapter were all answerable from inside a
$0$-D collisional--radiative matrix. Five are not, and they are gathered here
rather than left scattered so that a reader can see at once what this thesis
does not determine. Each is stated with the reason it is out of reach and with
the calculation that would settle it. None is a defect of the implementation;
each is a boundary of the model class.

\paragraph{1. The $r_1$ deficit is unexplained.} This model's ground-fed
population coefficient sits a factor $8.3$ below Fujimoto's Table~4.1(b) at
$p=3$, the shell the mechanism runs on, and a factor $1.8$ to $2.0$ below it at
$\log_{10}(\nel/\si{\per\cubic\metre}) = 21$; truncation (\SI{1.2}{\percent}),
Lyman trapping and the source's own statistical-$\ell$ closure
(\SI{0.42}{\percent}) have each been eliminated (Section~\ref{sec:r1_deficit}).
What would settle it is an $\ell$-resolved published tabulation or the atomic
data underlying the tabulated values, neither of which was available. The
deficit is reported as open, with its consequence for $\ffrac{3}-\ffrac{4}$
computed in Section~\ref{sec:fujimoto}.

\paragraph{2. There is no convergence study above $n_{\max} = 15$.} The state
space is cut at $n = 15$ because the atomic data stop there: the quantum
cross sections extend to $n = 10$, shells $11$ to $15$ are already semiclassical
Vriens--Smeets values, and no radiative or recombination coefficients exist
above $15$ in any form. Truncation has been tested downward, and it moves
$\ffrac{3}-\ffrac{4}$ by about \SI{1}{\percent} with the sign changing across
the grid (Section~\ref{sec:convergence}). What is untested is the sensitivity of
the ridge \emph{position} to $n_{\max}$, which would require rebuilding the rate
pipeline at $n_{\max} = 20$ rather than running a validation script: new
excitation, ionisation, radiative and recombination coefficients for five
further shells, and a re-indexing of every array built on the $43$-state layout.
That is a rate-generation exercise, not an analysis one.

\paragraph{3. The ground-state population is solved for, not supplied.} Every
magnitude in Chapter~\ref{ch:results} is conditional on $\ngs$ coming from
$\Lmat\nvec + \Svec\nion = 0$; in a divertor it is set by recycling, on a
transport timescale four orders of magnitude shorter than the local ionisation
balance (Section~\ref{sec:transport}). The open-boundary test of
Section~\ref{sec:attacks} bounds the sensitivity without removing the
conditional, and what would settle it is stated in Section~\ref{sec:transport}.

\paragraph{4. There are no molecular channels in $\Lmat$.} At the onset of
detachment a large fraction of the Balmer light is molecular in origin. The
dilution estimate of Section~\ref{sec:molecules} says how far the cold-corner
number would move under one specific dilution picture, not what a molecular
model would give. What would settle it is a state space containing
$\mathrm{H}_2$, $\mathrm{H}_2^+$ and $\mathrm{H}^-$ with their rate set, a
different model rather than an extension of this one.

\paragraph{5. The electrons are assumed Maxwellian.} The reasons are given in
Section~\ref{sec:other_assumptions}. The rate integrals as built do not accept
an arbitrary electron energy distribution, so this is not a parameter that can
be varied within the present code, and the $T_i = \Te$ assumption in the
$\ell$-mixing rates is untested for the same reason.

\medskip

\noindent Items 3 and 4 are named in Chapter~\ref{ch:conclusions} among the
further work. Items 1 and 2 are bounded in their consequences rather than in
their causes: the thesis states what each would change if it went the wrong
way, and for item 1 that stake is computed in Section~\ref{sec:fujimoto}.

%-----------------------------------------------------------------------------
\section{What is left standing}
\label{sec:limits_summary}
% QUESTION: after all of that, what may this thesis actually claim?

\subsection{What this model cannot say}

It cannot say anything quantitative about the cold edge of its own grid. Below
\SI{2}{\electronvolt} the optically thin plateau error is inflated by a factor
of $2.5$ to $3.3$ by neglected Lyman trapping; the size of the correction
depends on a slab thickness the model does not contain; and below
\SI{1.15}{\electronvolt} even the location of the ridge is undetermined.

It cannot say what a real divertor does, because recycling renews the
ground-state population there on a transport timescale the model does not
contain: \SI{72}{\micro\second} against \SI{233}{\milli\second} at the worst
point, a factor $3.2\times10^{3}$, below the scope boundary; inside the defended
range one to three orders rather than four, with a median $\tauesc/\tauslow$ of
$0.0079$ over the $45$ defended pairs (Table~\ref{tab:transport_selection}); and
$3.6$ at the benchmark, where transport is the \emph{slower} of the two.

It cannot say what an instrument records at the onset of detachment, because a
large fraction of that light is molecular in origin and there are no molecules
in $\Lmat$.

It cannot name the density at which its own ridge sits as a region of a tokamak.
The ridge falls between the modelled upstream separatrix and the detached
divertor band, and belongs to neither.

And it cannot establish the absolute normalisation of its ground-fed
population coefficient against an independent calculation: $r_1$ sits a factor
$8.3$ below the standard tabulation at $p = 3$, the disagreement is
unexplained, and its stake for $\ffrac{3}-\ffrac{4}$ is computed in
Section~\ref{sec:fujimoto} (Sections~\ref{sec:r1_deficit}
and~\ref{sec:scope_open}).

\subsection{What this model can say, with the objections computed}

Above \SI{2}{\electronvolt} the breakdown census is invariant under Lyman
trapping across a twentyfold range of assumed slab thickness, $45$ of $448$ in
every run with the worst case moving by \SI{0.5}{\percent}
(Section~\ref{sec:opacity}). Closing the open boundary moves the census by two
points out of $680$ and the worst case by one part in $325$; the factor $2300$
that would have refuted the result at the cold corner measures $15$, and at the
ridge the factor required is $20$ against a measured $1.01$
(Section~\ref{sec:attacks}). The ridge mechanism was verified over
$\Te \ge \SI{2}{\electronvolt}$, where trapping is now shown to be irrelevant,
so the agreement with Griem's boundary criterion is not an artifact of neglected
opacity. At citable divertor densities above \SI{2}{\electronvolt} the
\SI{100}{\micro\second} time-averaged error never reaches \SI{10}{\percent} over
$108$ points, worst case \SI{7.2}{\percent} (Section~\ref{sec:not_detachment}).

\subsection{The bridge}

The pattern is consistent, and it is why this chapter is not an apology. Every
objection that could be computed inside a $0$-D atomic model was computed, and
the defensible claim came through all of them with its numbers unchanged. The
objections that remain, neutral transport and molecular chemistry, are the
ones that cannot be computed inside a $0$-D atomic model at all, and they have
been stated as conditionals rather than dissolved into caveats.

That is the boundary of the claim. What is inside it is a quantitative statement
about when a spectroscopist may trust a lookup table and when she may not, and
Chapter~\ref{ch:conclusions} puts it into the form she can use.
 -- those labels are reused below, so every existing
%  cross-reference from Chapters 3 and 5 continues to resolve.
%
%  Sources for every number, in order of appearance:
%    outputs/findings_10_four_agent_review.md   ss 3.1, 4.1-4.4, 7.3, 8, 10,
%                                               and the ADDENDUM (A.1-A.6)
%    validation/lyman_trapping/lyman_trapping.txt    (run 9-10 Sep 2026,
%                                               src/validation/verify_lyman_trapping.py)
%    thesis_tex/chapter5_C5D.tex:284    (the "It is not detachment" wording)
%    outputs/findings_09_central_quantity_misnamed.md   ss 8, 9
%=============================================================================

\chapter{What This Model Cannot Say}
\label{ch:limits}

Chapter~\ref{ch:results} handed the reader two things: a map of how wrong a
tabulated $H_\alpha/H_\beta$ inversion can be while the ground state is still
catching up, and a mechanism that explains the map's shape. The map has a ridge
in it, the ridge has a closed-form explanation, and the explanation was tested
three separate ways. At this point in the argument the reader believes a
result.

This chapter's job is to bound that belief without destroying it. Every number
in Chapter~\ref{ch:results} came out of one $43$-state matrix describing one
kind of atom under one set of assumptions, and some of those assumptions are
badly violated in exactly the corner of the map where the numbers are largest.
An examiner will find that corner. It is better to arrive there first, with
arithmetic.

\medskip
\noindent\textbf{The standard this chapter holds itself to.} A limitation that
has been \emph{computed} is worth more than a limitation that has been hedged.
Where an objection could be turned into a calculation, it was, and the
calculation is reported here with the number that would have refuted the
result. Two objections were expected to be fatal and were not; they are given
the same space as the ones that bite. Where an objection cannot be computed
inside a model with no spatial dimension, it is stated as a \emph{conditional},
``\emph{if} this holds, the error is as computed'', and not softened into a
caveat sentence. Where a check has simply not been run, this chapter says so.

\medskip
\noindent\textbf{Three abbreviations recur below} and are expanded here for
this chapter: \emph{CR} for collisional--radiative, the class of model built in
Chapter~\ref{ch:theory}; \emph{QSS} for quasi-steady-state, the Bodenstein-style
approximation that eliminates the fast excited levels; and \emph{CRE} for
collisional--radiative equilibrium, the fully settled steady state that the
diagnostic lookup table assumes and that this thesis measures the cost of
assuming. An \emph{ELM} (edge-localised mode) is the repetitive burst of
heat and particles that a tokamak edge throws at the divertor every few
milliseconds, and it is the transient every ``event duration'' below refers to;
$\taudrive = \SI{100}{\micro\second}$ throughout, the duration used in
Chapter~\ref{ch:results} and the aggressive end of the range; the
ITER-pedestal extrapolation is \SI{506}{\micro\second}
\cite[Sec.~5]{Loarte2003} (Section~\ref{sec:persistence}).

%-----------------------------------------------------------------------------
\section{Optical thickness: when the plasma cannot let its own light out}
\label{sec:opacity}
% QUESTION: the model lets every emitted photon escape the plasma -- where is
% that false, what does it do to the error map, and how much of the map
% survives it?

\subsection{The question, and the picture behind it}

Chapter~\ref{ch:atoms} gave every downward transition an Einstein $A$
coefficient: the probability per second that an excited atom drops to a lower
level and emits a photon. Those coefficients entered the CR matrix $\Lmat$
bare. Entering bare means one specific physical assumption: that every photon
emitted anywhere in the plasma leaves it and never comes back. A plasma for
which this is true is called \emph{optically thin}.

A chemical engineer has met the failure of this assumption in another costume.
Consider a reaction whose product is removed by desorption from a catalyst
surface. If the product leaves the reactor, the desorption step runs at its
intrinsic rate. If the product cannot leave, if it re-adsorbs somewhere else
in the bed and is put back onto a site, then the \emph{effective} desorption
rate is smaller than the intrinsic one, and the whole network downstream of that
step shifts. Radiation trapping is that, with photons.

Here is the specific chain. The strongest transition in hydrogen is
Lyman-$\alpha$: an atom in the $2p$ level drops to the ground state $1s$ and
emits an ultraviolet photon. If the gas around it contains many ground-state
atoms, that photon has a good chance of being absorbed by one of them, which
puts \emph{that} atom back into $2p$. Nothing has left the plasma; a
$2p$ population has simply been moved from one place to another. Averaged over
the volume, the $2p$ level decays more slowly than its $A$ coefficient says. The
standard way to write this down is to multiply the coefficient by an
\emph{escape factor} $\Theta_P$ between $0$ and $1$: $\Theta_P = 1$ is a
transparent plasma, $\Theta_P = 10^{-5}$ is one in which the transition is
effectively switched off. The formulation used here is Holstein's population
escape factor~\cite{Holstein1947}.

\textbf{Why this reaches the populations, and therefore the result.} If $2p$
lives longer, atoms pile up in $n=2$. Atoms in $n=2$ are only
\SI{10.2}{\electronvolt} from the ground state but only
\SI{3.4}{\electronvolt} from the continuum, so they are far easier to ionise
than ground-state atoms. A longer $n=2$ lifetime therefore means more
\emph{stepwise} ionisation, ionisation via an excited level rather than
directly out of the ground state, which means fewer neutrals and a faster
turnover of the ground-state reservoir. The ground-state reservoir is the whole
mechanism of this thesis. Trapping is not a cosmetic correction to a line
intensity; it acts on the slow clock.

\subsection{Where the assumption fails, in numbers}

The formalism needed is one line. A photon travelling a distance $D$ through a
gas of absorbers is attenuated by $e^{-\tau}$, where the \emph{optical depth}
$\tau = \sigma_0\,\ngs\,D$ is dimensionless: $\ngs$ is the ground-state density
in \si{\per\cubic\centi\metre}, $D$ the path length in \si{\centi\metre}, and
$\sigma_0$ the absorption cross section at line centre in
\si{\square\centi\metre}. For a Doppler-broadened line,
$\sigma_0 = 2.654\times10^{-2}\,f/(\sqrt{\pi}\,\Delta\nu_D)$, with $f$ the
oscillator strength (a dimensionless measure of how strongly the transition
couples to light) and $\Delta\nu_D$ the Doppler width set by the temperature of
the absorbing atoms. $\tau \ll 1$ is transparent; $\tau \gg 1$ is opaque.

The number that matters is $\ngs$, not $\nel$, and Chapter~\ref{ch:theory}
already established the awkward fact that these two run in opposite directions:
the plasma is most neutral where it is coldest. The worst point of the error map
is $\Te = \SI{1.0}{\electronvolt}$,
$\nel = 5.18\times10^{13}\,\si{\per\cubic\centi\metre}$, the maximum of
$\epsplat$ over all $680$ window-resolvable points of the map, cold edge
included. There the model's own CR solution gives $\ngs/\nion = 28.3$, so
quasi-neutrality fixes
$\ngs = 1.47\times10^{15}\,\si{\per\cubic\centi\metre}$. At
\SI{1}{\electronvolt} that gives
$\sigma_0 = 5.47\times10^{-14}\,\si{\square\centi\metre}$, and the Lyman-$\alpha$
optical depth per centimetre of path is

\begin{center}
\begin{tabular}{lr}
  \toprule
  $\Te$ (eV) & $\tau_{\rm Ly\alpha}$ per \si{\centi\metre} at
              $\nel = 5.18\times10^{13}\,\si{\per\cubic\centi\metre}$ \\
  \midrule
  1.00 & 80.3 \\
  1.15 & 10.6 \\
  1.60 & 0.224 \\
  2.02 & 0.0255 \\
  2.95 & 0.00163 \\
  \bottomrule
\end{tabular}
\end{center}

\noindent An optical depth of $80$ per centimetre means a Lyman-$\alpha$ photon
travels about a tenth of a millimetre before it is reabsorbed. Over
\SI{5}{\centi\metre} of plasma the line-centre depth
over the half slab is $201$ and the population escape factor is
$1.2\times10^{-3}$; roughly one photon in eight hundred gets out. Chapter~%
\ref{ch:theory} quotes $3\times10^{-5}$ at a different grid point, a different
slab thickness and a different density, so the two are not the same quantity and
the agreement in order of magnitude should not be read as a reproduction of
one by the other.
The blunt reading is that the effective $A_{2p\to1s}$ belonging in $\Lmat$ at
that point is wrong, self-consistently, by about $1.8$ orders of magnitude, or
$2.9$ orders on the thin one-shot estimate
(\texttt{verify\_lyman\_optical\_depth.py}).

Two rows down the same table the problem has evaporated. At
\SI{2.02}{\electronvolt} the depth per centimetre is $0.0255$, and at the
benchmark point (\benchTe, \benchne) it is $0.0039$; a photon there crosses
metres of plasma before anything happens to it. \textbf{The measured quantity
is not ``how opaque is this plasma'' but ``how fast does opacity switch on as
the plasma cools'',} and the answer is a factor of $5\times10^{4}$ between
\SI{1}{\electronvolt} and \SI{2.95}{\electronvolt}.

That is why the results of Chapter~\ref{ch:results} were declared for
$\Te \gtrsim \SI{2}{\electronvolt}$ in the first place
(Section~\ref{sec:two_references}). It is also, uncomfortably, where the
headline was living. Recounting the \SI{100}{\micro\second} breakdown census against optical depth
rather than against temperature: of the $202$ points at which the
\SI{100}{\micro\second} time-averaged estimate of the error exceeds \SI{10}{\percent}, an extremum count over
the $680$ window-resolvable points of the full grid, $157$ lie below
\SI{2}{\electronvolt}, $85$ have
$\tau_{\rm Ly\alpha}(\SI{5}{\centi\metre}) > 1$, and $15$ have it above $100$,
on the half-slab depth used above.

An earlier draft gave $105$ and $26$, a convention error rather than a
recomputation: those counts are reproduced only by combining the full-path
depth $\kappa_0 D$ with a Doppler width $\sqrt{kT_n/m}$, short of the correct
$\sqrt{2kT_n/m}$ by $\sqrt{2}$. With the correct width the census is $85$ and
$15$ on the half-slab depth quoted here and $100$ and $23$ on the full-path
depth. A deuterium cross-section with the full-path depth also returns $105$
and $26$, since the isotope mass enters the width as the same $\sqrt{2}$, so
the two errors are indistinguishable from the counts alone. All seven
conventions are tabulated in \texttt{validation/lyman\_optical\_depth/}.
Nothing in the defended range depends on the choice: every crossing lies below
\SI{2}{\electronvolt} except two, both at $D = \SI{20}{\centi\metre}$.

\subsection{What we expected trapping to do, and what it did}

The expectation, stated before the calculation was run, was that trapping would
demolish the map: an effective $A$ coefficient wrong by about three orders of
magnitude at the worst point ought to invalidate everything computed there and
nearby, and the project's own records warned that under strong trapping the
fast relaxation time $\taurel$ rises to \SI{2.05}{\micro\second} and the
timescale separation $\Msep$ falls by a factor $158$. That expectation was
wrong.

So the calculation was done. The Lyman series was given a self-consistent
Holstein escape factor and the entire error map rebuilt. Self-consistent matters
here: $\Theta_P$ depends on $\ngs$, and $\ngs$ depends on $\Theta_P$ through the
CR balance, so the two were iterated to a fixed point rather than evaluated once.
The fixed point converged at all $400$ grid points for every slab thickness
tested, and non-convergence was made to raise an error rather than silently
return the last iterate. All $14$ Lyman channels were trapped, the seven
resolved $n\mathrm{P}\to1\mathrm{S}$ transitions and the seven bundled ones,
not Lyman-$\alpha$ alone.

Three gates were passed before any result was looked at. The oscillator
strengths were derived from this repository's own $A$ coefficients by inverting
the Einstein relation rather than imported, and come out as $f = 0.4162$,
$0.0791$, $0.0290$, $0.0139$ for Lyman-$\alpha$ through Lyman-$\delta$, the
accepted literature values to four figures. The derived Lyman-$\alpha$
$\sigma_0$ agrees with an independent implementation to \SI{0.059}{\percent},
after a first run that failed at \SI{1156}{\percent} on a $4\pi$ factor; and
rebuilding the matrix with trapping switched off reproduces the canonical
$\Lmat$ exactly, so every trapped number below is a difference against the real
matrix (Section~\ref{sec:severe_checks}).

\textbf{One new quantity had to be invented to do this at all, and it is not in
the model.} The escape factor depends on how far a photon must travel to get
out, so it depends on the size of the plasma, the slab thickness $D$. A model
with no spatial dimension has no $D$. It was therefore swept, from
\SI{1}{\centi\metre} to \SI{20}{\centi\metre}, and never fixed; every trapped
number below carries the $D$ it was computed at. Reporting a single trapped
number without its $D$ would be reporting an assumption as a measurement.

\begin{table}[tbp]
  \centering
  \caption[Breakdown census under Lyman trapping]{Effect of radiation
  trapping on the \SI{100}{\micro\second} breakdown census. Each entry counts the grid points at
  which the time-averaged estimate of the CR-equilibrium-reference error,
  averaged over a chosen duration $\taudrive = \SI{100}{\micro\second}$, exceeds
  \SI{10}{\percent}, out of the points at which a
  quasi-steady-state window is resolvable; ``worst'' is the largest estimate
  within that scope. Rows differ only in the assumed slab thickness $D$ used to
  compute the self-consistent Holstein escape factor applied to all $14$ Lyman
  transitions. The first row is the optically thin matrix used throughout
  Chapter~\ref{ch:results}. Above \SI{2}{\electronvolt} the count does not move
  at all across a twentyfold range of $D$. The $D = \SI{20}{\centi\metre}$ row
  has a slightly smaller denominator because two points lose their resolvable
  window under trapping. Source:
  \texttt{validation/lyman\_trapping/lyman\_trapping.txt}.}
  \label{tab:trapping_census}
  \footnotesize
  \begin{tabular}{lccc}
    \toprule
    matrix & all window-resolvable & $\Te \ge \SI{2}{\electronvolt}$ &
          $\Te \ge \SI{2}{\electronvolt}$, $\nel \ge 10^{14}$ \\
    \midrule
    optically thin (canonical)          & $202/680$, worst $0.3868$ & $45/448$, worst $0.1748$ & $0/108$, worst $0.0717$ \\
    trapped, $D = \SI{1}{\centi\metre}$  & $176/680$, worst $0.2925$ & $45/448$, worst $0.1748$ & $0/108$, worst $0.0712$ \\
    trapped, $D = \SI{5}{\centi\metre}$  & $170/680$, worst $0.2550$ & $45/448$, worst $0.1746$ & $0/108$, worst $0.0694$ \\
    trapped, $D = \SI{20}{\centi\metre}$ & $160/678$, worst $0.2237$ & $45/446$, worst $0.1740$ & $0/106$, worst $0.0627$ \\
    \bottomrule
  \end{tabular}
\end{table}

\subsection{Above two electron-volts, nothing moves}

Table~\ref{tab:trapping_census} contains the strongest robustness statement in
this thesis, and it is the opposite of what was expected.

\textbf{Above \SI{2}{\electronvolt} the breakdown count does not change at all,
$45$ of $448$ points in every run, and the worst case within that set
moves by \SI{0.5}{\percent}, from $0.1748$ to $0.1740$, across a twentyfold
range in assumed slab thickness.} The escape factor at the benchmark point
(\benchTe, \benchne) runs from $0.9986$ at $D = \SI{1}{\centi\metre}$ to
$0.973$ at $D = \SI{20}{\centi\metre}$. Even at the thickest slab tested,
\SI{97}{\percent} of Lyman-$\alpha$ photons still leave: the plasma is optically
thin there in the sense that actually matters, and assuming so costs less than
the resolution of the map.

This converts a disclaimer into a measurement. The
$\Te \gtrsim \SI{2}{\electronvolt}$ restriction was chosen on an escape-factor
argument before any of this was recomputed, and it sits almost exactly at the
temperature above which trapping stops mattering: a guess that has now been
checked, and survived. The defensible claim of Chapter~\ref{ch:results}, $45$
of $448$ points above \SI{2}{\electronvolt} at \SI{100}{\micro\second} and $24$
of $448$ at \SI{506}{\micro\second} (Section~\ref{sec:persistence}), is
invariant under the one physical process most likely to have invalidated it.

\subsection{Below two electron-volts, the cold-edge number is inflated
threefold}

The same calculation is much less kind to the cold edge, and it should be,
because that is where the optical depth was $80$ per centimetre.
Table~\ref{tab:trapping_worstpoint} follows the worst point of the optically
thin map through the sweep.

\begin{table}[tbp]
  \centering
  \caption[Lyman trapping at the cold-edge worst point]{The cold-edge worst
  point of the optically thin error map, at $\Te = \SI{1.0}{\electronvolt}$ and
  $\nel = 5.18\times10^{13}\,\si{\per\cubic\centi\metre}$ under a heating step,
  is the maximum of $\epsplat$ over all $680$ window-resolvable grid points; it
  is traced here as a function of the assumed slab thickness $D$. $\epsplat$ is the plateau
  error, the fractional error a lookup table makes while the ground state is
  stale; the ELM bound is its average over an event of
  $\taudrive = \SI{100}{\micro\second}$. The plateau error falls by a factor
  $2.5$ to $3.3$, but the diagnostic's intrinsic sensitivity
  $|\ffrac{3}-\ffrac{4}|$ barely moves: what collapses is how far the ground
  state is displaced, not how strongly the line ratio responds to it. Source:
  \texttt{validation/lyman\_trapping/lyman\_trapping.txt}.}
  \label{tab:trapping_worstpoint}
  \small
  \begin{tabular}{lrrrrrrr}
    \toprule
    $D$ (cm) & $\epsplat$ & ELM bound & $\tauslow$ (s) & $\Msep$ &
    $\ffrac{3}$ & $\ffrac{4}$ & $|\ffrac{3}-\ffrac{4}|$ \\
    \midrule
    $0$ (thin) & $0.38690$ & $0.38682$ & $2.3324\times10^{-1}$ & $3.74\times10^{7}$ & $0.7728$ & $0.2862$ & $0.4866$ \\
    $1$        & $0.15748$ & $0.15739$ & $9.0651\times10^{-2}$ & $2.75\times10^{6}$ & $0.6643$ & $0.2055$ & $0.4589$ \\
    $5$        & $0.13274$ & $0.13260$ & $4.6940\times10^{-2}$ & $6.13\times10^{5}$ & $0.5967$ & $0.1793$ & $0.4174$ \\
    $20$       & $0.11570$ & $0.11550$ & $2.8500\times10^{-2}$ & $2.14\times10^{5}$ & $0.5139$ & $0.1501$ & $0.3639$ \\
    \bottomrule
  \end{tabular}
\end{table}

\textbf{The \SI{38.7}{\percent} becomes \SIrange{11.6}{15.7}{\percent}}, a
factor of $2.5$ to $3.3$, and which value inside that range it takes is set
by a slab thickness the model does not contain. The point still exceeds the
\SI{10}{\percent} threshold at every thickness tested, so the qualitative
statement that the cold corner breaks down survives intact; the \emph{number}
does not, and it must never appear without its trapped range and without the
optical depth that produced it.

The route by which the number falls is itself a result, and it is not the
obvious route. The obvious guess is that trapping desensitises the diagnostic,
that a trapped plasma simply responds less. It does not. The sensitivity
$|\ffrac{3}-\ffrac{4}|$, which is the factor in the exact decomposition of
Section~\ref{sec:mechanism} measuring how strongly the line ratio answers to its
reservoir, moves by only \SI{6}{\percent} at $D = \SI{1}{\centi\metre}$ and
\SI{25}{\percent} at $D = \SI{20}{\centi\metre}$. What collapses is the other
factor: the displacement.

The chain is the one set out at the start of this section, now with numbers on
it. Trapping lengthens the $n=2$ lifetime; the longer lifetime raises stepwise
ionisation; the ground-state population falls; and the slow clock $\tauslow$
shortens by a factor $8$ at the worst point. At $\Te = \SI{1.0}{\electronvolt}$,
$\nel = 1.93\times10^{13}\,\si{\per\cubic\centi\metre}$, $\ngs$ drops from
$2.99\times10^{14}$ to $7.26\times10^{13}\,\si{\per\cubic\centi\metre}$ between
$D = 1$ and \SI{20}{\centi\metre}. A reservoir that turns over faster is a less
stale reservoir, so $|\Delta\ln\xscale|$ roughly halves. This is exactly the
two-axis attribution of Section~\ref{sec:temperature_mechanism} arriving from a
third direction: trapping acts on the displacement, not on the sensitivity.

\subsection{Below 1.15 electron-volts the ridge location is not determined}

One more thing moves at the cold edge, and it is a location rather than a
magnitude. Taking the density at which $\epsplat$ is maximal, row by row: at
$\Te \ge \SI{1.6}{\electronvolt}$ that maximum falls on
$\nel = 1.93\times10^{13}\,\si{\per\cubic\centi\metre}$ in every run, optically
thin and trapped alike, and the values themselves barely register the change:
at the benchmark temperature $\epsplat = 0.12166$, $0.12166$, $0.12165$,
$0.12163$ for $D = 0, 1, 5$ and \SI{20}{\centi\metre}, with
$|\ffrac{3}-\ffrac{4}|$ changing in the fifth decimal place. Below
\SI{1.15}{\electronvolt} the maximum wanders across three density columns, a
factor of $7$, as a function of an assumed parameter. \textbf{The two coldest
rows of the grid cannot support a claim about where the ridge is at all.}

That result cuts in the mechanism's favour, which is worth saying explicitly
because it is easy to miss. The ridge attribution of
Section~\ref{sec:mechanism} was verified over $\Te \ge \SI{2}{\electronvolt}$,
and $\Te \ge \SI{2}{\electronvolt}$ is precisely the range in which trapping has
now been shown to be irrelevant. The agreement with Griem's boundary criterion
is therefore not an artifact of neglected opacity; it was established where
opacity does not act.

\subsection{A contradiction resolved, and what is still open}

Two statements about trapping had been circulating in this project's records:
that $\taurel$ is insensitive to Lyman-$\alpha$ trapping to within
\SI{1}{\percent}, and that under strong trapping $\taurel$ rises to
\SI{2.05}{\micro\second} with $\Msep$ falling $158$-fold. They look
contradictory. Measured directly, both are true and neither had stated its
scope: at the benchmark point (\benchTe, \benchne) $\taurel$ moves from
$2.2769\times10^{-9}$ to $2.2947\times10^{-9}\,\si{\second}$ at
$D = \SI{20}{\centi\metre}$, a change of \SI{0.78}{\percent}; at
$\Te = \SI{1.0}{\electronvolt}$,
$\nel = 1.93\times10^{13}\,\si{\per\cubic\centi\metre}$ it moves from
$8.8842\times10^{-9}$ to $2.4289\times10^{-7}\,\si{\second}$, a factor of $27.3$.
The ``within \SI{1}{\percent}'' claim is a benchmark-point claim and fails by
more than an order of magnitude at the cold corner. This is the same failure
mode as quoting any other extremum without the set it is an extremum over, and
neither sentence may travel again without its regime attached.

What this calculation does not settle should be as clear as what it does. The
geometry is assumed, not derived: a homogeneous slab, uniform emission, an
isotropic radiation field, a Doppler line profile, evaluated at the plasma
centre. A $0$-D model cannot check any of those. One assumption inside it
\emph{was} checked and is not load-bearing: the temperature of the absorbing
atoms, which sets the Doppler width. Fixing it at \SI{3}{\electronvolt}, the
value a molecule dissociating by the Franck--Condon process would give, instead
of at $\Te$, changes the census at $D = \SI{5}{\centi\metre}$ from $170/680$ and
a worst case of $0.2550$ to $173/680$ and $0.2653$. Continuum lowering is
untouched here and remains open. Opacity in the Balmer lines themselves, which
would act on the observed ratio directly rather than through the level
populations, has been computed: the largest $H_\alpha$ optical depth above
\SI{2}{\electronvolt} is $0.010$, and the worst cold cell, $[0,7]$, reaches
$0.214$ (Section~\ref{sec:balmer_equivalence},
\texttt{validation/balmer\_optical\_depth/}).

%-----------------------------------------------------------------------------
\section{Transport: the atoms do not stay where the model puts them}
\label{sec:transport}
% QUESTION: the slow clock is the ionisation time of a stationary atom -- what
% happens to the result when the atom leaves before it ionises?

\subsection{The question}

This is the objection that cannot be answered inside this model, and it is a
more serious threat to the cold-edge number than anything in
Section~\ref{sec:opacity}. It is also the simplest to state.

\subsection{The picture: a reactor with no walls}

The model of Chapter~\ref{ch:theory} is a batch reactor. A fixed parcel of gas
is given a temperature, and its composition evolves until it stops changing. The
slow eigenvalue of that system, $\tauslow$, is how long the ground-state
population takes to reach its new value, and it is long, because ionising an
atom out of the ground state is slow.

A divertor is not a batch reactor. It is closer to a continuous-flow reactor
with an enormous throughput, in which the residence time of a molecule is far
shorter than the time it needs to react. Neutral atoms are born at the target
surface when ions strike it and recombine there, which is \emph{recycling},
and they fly back into the plasma. They do not wait around to be ionised
locally; the ones that do not ionise leave, and fresh ones arrive continuously.

The comparison is not close. At the worst point of the map $\tauslow$ is
\SI{233}{\milli\second}: nearly a quarter of a second for a stationary atom to
be ionised. A deuterium atom at \SIrange{1}{3}{\electronvolt} would cross a
\SI{10}{\centi\metre} plasma ballistically in \SIrange{6}{10}{\micro\second}, but
it does not travel ballistically: resonant charge exchange with the surrounding
ions randomises its direction after about \SI{1.7}{\centi\metre} at this
density, so escape is diffusive and takes closer to \SI{72}{\micro\second}. That
correction runs in this model's favour. The threshold at which the worst point
would fall below \SI{10}{\percent} is a ground-state renewal time of about
\SI{26}{\micro\second}, so the result survives by a factor of $2.8$ rather than
failing outright. It survives only because charge exchange traps the neutrals,
which is a thin margin to rest on and is stated as one.
That is one grid point, and it is one this thesis does not defend:
$[0,4]$ lies below the \SI{2}{\electronvolt} scope boundary.
Section~\ref{sec:transport_selection} asks the same question at the points that
do carry the result, and gets a worse answer.

\begin{quote}
\textbf{In a divertor, the ground-state neutral density at a given point is not
set by local ionisation balance. It is set by recycling.}
\end{quote}

\subsection{Why this bears on the result specifically}

The error of Chapter~\ref{ch:results} exists for one reason: after a temperature
step, the ground-state reservoir is \emph{stale}. Its population has not moved
while the temperature has, the excited levels equilibrate against a reservoir
that belongs to the old conditions, and the lookup table, which assumes the
reservoir finished settling, reads the wrong answer. Everything hangs on the
reservoir not being refreshed.

Neutral transport refreshes it. Any process that replaces the ground-state
population on a timescale shorter than $\tauslow$ resets the staleness, and this
model contains no such process; the \SI{26}{\micro\second} renewal threshold
quoted above is the map's own measure of how fast that process would have to be.

\subsection{The assumption fails hardest where the result is generated}
\label{sec:transport_selection}

The worst point was the wrong place to test transport: $[0,4]$ lies below the
scope boundary of Section~\ref{sec:scope}. The pairs that carry the defended
claim are the $45$ above \SI{2}{\electronvolt} whose \SI{100}{\micro\second}
estimate exceeds \SI{10}{\percent} (Table~\ref{tab:elm_census}), and there the
answer is different.

Comparing the neutral escape time with $\tauslow$ pair by pair over the grid
gives Table~\ref{tab:transport_selection}. The escape model is the one used
above, stated as a formula: the atom moves at $v=\sqrt{2k\Tn/m_{\mathrm D}}$
with $\Tn=\Te$, resonant charge exchange with the ambient ions randomises its
direction after $\lambda = v/(\nel\langle\sigma v\rangle_{\mathrm{CX}})$ with
$\langle\sigma v\rangle_{\mathrm{CX}} = \SI{1.1e-8}{\cubic\centi\metre\per\second}$,
and escape from a boundary a distance $L$ away takes $L/v$ when
$\lambda \ge L$ and $4L^{2}/\pi^{2}D$ with $D = v\lambda/3$ otherwise. At
$[0,4]$ it returns $\lambda = \SI{1.72}{\centi\metre}$ and
$\tauesc = \SI{72.3}{\micro\second}$, the values used above
(\texttt{verify\_transport\_selection.py}).

\begin{table}[htbp]
\centering
\caption{Ratio of neutral escape time to $\tauslow$, over the $45$ defended
breakdown pairs and over the other $635$ analysed pairs. $L$ is the distance
from the emitting parcel to the plasma boundary. ``Satisfies'' means
$\tauesc > \tauslow$, that is, the atom is ionised before it leaves, which is
what the closed parcel assumes.}
\label{tab:transport_selection}
\begin{tabular}{llrrr}
  \toprule
  $L$ & set & median $\tauesc/\tauslow$ & largest & satisfies \\
  \midrule
  \SI{10}{\centi\metre} & the 45 breakdown pairs & $0.0079$ & $0.159$ & $0/45$ \\
  \SI{10}{\centi\metre} & the other 635 & $0.0471$ & $384$ & $168/635$ \\
  \SI{20}{\centi\metre} & the 45 breakdown pairs & $0.0192$ & $0.636$ & $0/45$ \\
  \SI{20}{\centi\metre} & the other 635 & $0.129$ & $1535$ & $227/635$ \\
  \bottomrule
\end{tabular}
\end{table}

\textbf{Not one of the $45$ pairs that carry the defended result satisfies the
assumption used to compute it, at either scale length, while a quarter to a
third of the pairs that do not break down satisfy it comfortably.}

A count of zero is worth little if the $45$ sit just under the line, so the
distance matters. The closest of them is $[18,4]$ heating,
$\Te = \SI{2.33}{\electronvolt}$, $\nel = \SI{5.18e13}{\per\cubic\centi\metre}$,
at $\tauesc/\tauslow = 0.636$ for $L = \SI{20}{\centi\metre}$: its escape time
would have to rise by a factor $1.57$ for the parcel to close. The median of the
$45$ would have to rise by a factor $52$. One pair is marginal and the rest are
not.

Three things were varied to see whether the conclusion depends on a choice.
Tripling $\langle\sigma v\rangle_{\mathrm{CX}}$, which lengthens every diffusive
escape time in proportion, moves the count from $0$ to $3$ of $45$; dividing it
by three leaves it at $0$. Replacing $\sqrt{2k\Tn/m}$ by the Maxwellian mean
speed $\sqrt{8k\Tn/\pi m}$ shortens every escape time and leaves the count at
$0$. Doubling $L$ from $10$ to \SI{20}{\centi\metre} leaves it at $0$. The
finding is not an artifact of any of the three.

\paragraph{The selection is structural, not incidental.} A pair enters the
census because $\tauslow$ is long, since the single-slow-mode estimate of
Eq.~\eqref{eq:elm_bound} rises with $\tauslow$ at fixed $\taudrive$. A long
$\tauslow$ means slow ionisation out of the ground state. Slow ionisation is
exactly the condition under which neutral transport, whose timescale does not
depend on $\tauslow$ at all, sets the ground-state density instead of the local
atomic physics. The census therefore selects for the regime in which its own
assumption is weakest, and it does so by construction rather than by accident.
That is why testing the assumption at the cold corner gave the reassuring answer
and testing it inside the defended scope does not.

\subsection{What this reaches, and what it does not}
\label{sec:transport_partition}

The finding above is about magnitudes; it does not touch the algebra.

\textbf{Untouched by transport.} The two-channel split of
Eq.~\eqref{eq:two_supplies}, exact to a relative residual of
$3.075\times10^{-14}$; the unit-width logistic form of the ground-fed fraction;
the bound $\max|f_3-f_4| = \tanh(|\Delta|/4)$ and its corollary
$|\mathrm{d}\ln\Robs/\mathrm{d}\ln\bdep| < 1$; the sensitivity map
$\overline{S}$; and the accuracy of the excited-state closure,
$6.73\times10^{-9}$ at the worst point. These are properties of the operator
$\Lmat(\Te,\nel)$ at a fixed operating point, or theorems about its structure.
$\overline{S}$ says how the line ratio responds to whatever the reservoir does;
it is silent on what made the reservoir do it. Transport changes the reservoir's
history and cannot reach any of them.

\textbf{Conditional on the closed parcel.} The reservoir gain $G$, which is
a secant along the CRE curve and therefore assumes the ground
state responds through collisional--radiative balance alone; $\epsplat$, which
is built from it; the breakdown census; and every time-averaged magnitude in
Chapter~\ref{ch:results}. If recycling renews $\ngs$ on a timescale shorter than
$\tauslow$, $G$ is not the gain the plasma has. These numbers are not upper
estimates. A fast-exchanging reservoir with a steady source freezes the
closed-parcel error rather than removing it, and a rising source enlarges it: over
the defended set the true error exceeds the closed-parcel plateau at $165$ of the
$180$ (pair, $L$, source-multiplier) combinations tested
(\texttt{verify\_open\_reservoir.py}). The closed-parcel value is a reference
magnitude conditional on the reservoir's history, not a bound on it.

The thesis therefore leads with the structure and the bound and reports the
magnitudes as conditional; Section~\ref{sec:transport_selection} adds that the
condition is measured to fail at the points that matter, by one to two orders
of magnitude.

\subsection{The claim, restated as the conditional it is}

The honest form is therefore not a caveat but a conditional, and it should be
written this way wherever the number appears:

\begin{quote}
\emph{If} the ground-state neutral density is frozen for the duration of the
event, the error is as computed.
\end{quote}

That conditional is close to exactly true for the excited manifold, which
equilibrates in nanoseconds, and it is the standard implicit assumption of every
$0$-D collisional--radiative table in use, the object this thesis set out to
test. But it is an assumption about the plasma, not about the atoms, and no
eigenvalue analysis inside a $0$-D matrix can validate it: the two timescales
are measured to differ by four orders of magnitude and assumed not to interact.
What would settle it is a two-region or transport-coupled calculation, an edge
code such as SOLPS (Scrape-Off Layer Plasma Simulation) supplying $\ngs$ from a
recycling influx; that is outside the scope of this thesis and is named in
Chapter~\ref{ch:conclusions} among the further work.

%-----------------------------------------------------------------------------
\section{No molecular channels, in a gas that is mostly molecules}
\label{sec:molecules}
% QUESTION: the model contains only atomic hydrogen -- at the conditions where
% the error is largest, is that a defensible description of the gas?

\subsection{The question, and why a referee asks it first}

Every state in Chapter~\ref{ch:atoms}'s state space is a level of a single
hydrogen \emph{atom}, plus a reservoir of bare protons. There is no
$\mathrm{H}_2$ (or $\mathrm{D}_2$), no molecular ion $\mathrm{D}_2^+$, no
negative ion $\mathrm{D}^-$, and no reactions between them. A search of the rate
assembly for any of these returns nothing at all.

\subsection{The picture: a third feed stream}

Chapter~\ref{ch:theory} fed every excited level from two directions, up from
the ground state by electron collisions and down from the ion continuum by
recombination, and the error mechanism is the difference in how the $n=3$ and
$n=4$ shells split their supply between those two feed streams.

Molecules open a third stream. In a cold, dense, mostly neutral gas, hydrogen
atoms form $\mathrm{H}_2$; the molecules are ionised to $\mathrm{H}_2^+$, or
exchange charge, and the resulting molecular ions recombine with electrons and
break apart, leaving one of the fragments in an \emph{excited} atomic level. The
whole family of such routes is called MAR, molecular activated recombination.
The important feature for this thesis is not that MAR exists but where it
delivers: it populates $n=3$ directly, without the atom ever having passed
through the ground-state reservoir whose staleness is the entire mechanism here.
A supply channel that bypasses the reservoir does not inherit the reservoir's
lag.

\subsection{How much gas is involved}

At the worst point of the map the model's own equilibrium solution is
\SI{96.6}{\percent} neutral: $\ngs = 1.47\times10^{15}$ against
$\nel = 5.18\times10^{13}\,\si{\per\cubic\centi\metre}$. A hydrogen gas at
\SI{1}{\electronvolt} that is \SI{96.6}{\percent} neutral is a gas in which
molecules form readily. The atomic description is not being applied at the edge
of its validity there; it is being applied outside it.

The size of what is being left out is documented experimentally.
\emph{Detachment} is the operating regime in which the plasma cools enough that
the ion flux to the target collapses, and the regime this diagnostic is chiefly
used to identify. In spectroscopic inversions of divertor Balmer emission at its
onset, molecularly-induced emission is found to account for up to
\SIrange{60}{70}{\percent} of the $D_\alpha$ light and
\SIrange{10}{20}{\percent} of $D_\gamma$~\cite{Wijkamp2023,Verhaegh2021}. An
$H_\alpha/H_\beta$ ratio computed from atomic processes alone is not the ratio
such an instrument records.

\subsection{What survives it, and what does not}

Two consequences follow, and they point in opposite directions, so both must be
stated.

The \emph{numerical} cold-edge results are atomic-model results and nothing
more. Adding a third supply channel that feeds $n=3$ preferentially would change
$\ffrac{3}-\ffrac{4}$, and could plausibly change it a great deal; this work
offers no direct calculation of how much. A bound can nonetheless be
constructed, from the emissivity fractions cited in this section. For a given
upper level the molecular fraction of the emissivity equals the molecular
fraction of the population, because both channels emit through the same
transition probability, and molecular-assisted recombination bypasses the ground
state entirely, so it dilutes the ground-fed fraction as
$\ffrac{p} \to \ffrac{p}(1-\phi_p)$. Taking $\phi_3 = 0.60$ and $\phi_4 = 0.30$
from the reported \SIrange{60}{70}{\percent} of $\mathrm{D}\alpha$ and
\SIrange{10}{20}{\percent} of $\mathrm{D}\gamma$, the difference of ground-fed
fractions at the cold corner falls from $0.4866$ to $0.1088$ and at the density
crest from $0.4580$ to $0.0998$, and the secant sensitivity $\overline{S}$
falls with it, from $0.4822$ to $0.0903$ and from $0.4552$ to $0.0903$.

The plateau error does not fall in the same proportion, because
$\epsplat = |\mathrm{e}^{\overline{S}\Lambda}-1|$ is not linear in
$\overline{S}$; treating it as linear understates the reduction. Imposing the
same fractions on the model's own operator, with a non-negative source
depositing into $n=3$ and $n=4$ so that $(\phi_3,\phi_4)$ is reached exactly at
the plateau, the cold-corner plateau error falls from \SI{38.7}{\percent} to
\SI{6.3}{\percent}, a \textbf{factor of $6.1$}; at the upper end of the
reported ranges, $(\phi_3,\phi_4) = (0.70,0.45)$, it falls to \SI{4.1}{\percent},
a factor of $9.4$. At the crest the same two targets give factors of $5.4$ and
$8.1$, and at the benchmark $3.2$ and $4.4$
(\texttt{validation/molecular\_channel/}).

Two features of that calculation are worth stating, because neither is visible
in the dilution formula. The two fractions cannot be chosen independently: a
source large enough to make $\phi_3 = 0.60$ already drives more than $0.15$ of
$n=4$ through collisional $3\to4$ transfer, so $(\phi_3,\phi_4) = (0.60,0.15)$
requires a negative $n=4$ source and is not realisable on this operator. And
the reduction is not monotone in $\phi_3$. The difference
$\ffrac{3}(1-\phi_3) - \ffrac{4}(1-\phi_4)$ passes through zero near
$\phi_3 = 0.74$ at $\phi_4 = 0.30$, beyond which $n=4$ is the more ground-fed
of the two and the error grows again with $\overline{S}$ carrying the opposite
sign. At the crest, where that zero is reachable with non-negative sources, a
two-source dilution at $\phi_4 = 0.30$ brings the plateau error down to
\SI{0.35}{\percent} at $\phi_3 = 0.75$, past the zero at $\phi_3 = 0.7355$; a
single $n=3$ source there gives \SI{0.30}{\percent} at $\phi_3 = 0.70$ and
\SI{1.1}{\percent} again at $\phi_3 = 0.80$, with $\overline{S}$ negative. A
molecular channel strong enough cancels the diagnostic's sensitivity to the
reservoir rather than merely reducing it, and past the cancellation the error
returns.

The bound is crude, it inherits every uncertainty in the quoted fractions, and
it applies only where those fractions do.

The \emph{structural} results, the list of Section~\ref{sec:transport_partition},
are theorems about whatever matrix $\Lmat$ is supplied and hold for one
containing molecular states; what such a matrix changes is the numerical value
of the ingredients, not the form of the result.

What is \emph{not} available is a calculation with molecules in it: the factors
above are the model's own operator driven by an imposed source of the published
strength, not a solution of a matrix containing molecular states. This is an
atomic model, its cold-edge numbers are atomic-model numbers, and the bound
above is an estimate of how wrong they are rather than a correction to them.

%-----------------------------------------------------------------------------
\section{Where the ridge sits depends on what the model assumed about neutrals}
\label{sec:ridge_conditional}
% QUESTION: is the density at which the error peaks a property of hydrogen, or
% a property of this model's own ionisation balance?

\subsection{The question}

Section~\ref{sec:mechanism} found an interior maximum, a ridge, in the
error map at $\nel \approx 2\times10^{13}\,\si{\per\cubic\centi\metre}$, and
showed that it moves with the shell pair as Griem's criterion for local
thermodynamic equilibrium (LTE) predicts. The mechanism is robust; three
independent tests agreed. But a mechanism and a location are different claims,
and only the first of them has been established.

\subsection{The picture: the ridge is where the two feeds are balanced}

The ridge exists because the $n=3$ and $n=4$ shells switch from
recombination-fed to ground-fed at different reservoir strengths, so their
difference vanishes when either feed dominates and peaks in between. The
controlling variable is therefore not the electron density directly; it is
\emph{how much neutral gas there is to feed from}. The electron density enters
only because, in this model, it determines the neutral density through the CR
ionisation balance.

That is the conditionality. Change what the model assumes about neutrals and the
ridge moves, whatever the electron density does.

\subsection{The measurement}

Within the two-channel split the ground-fed fraction is exactly
$\ffrac{p} = a_1 n_g/(a_0 + a_1 n_g)$, where $n_g$ is the ground-state density
feeding the excitation channel and $a_0$, $a_1$ are the network responses of
Section~\ref{sec:R_depends_on_b}. Scaling $n_g$ is therefore an exact operation
on the ridge, not an approximation to one. Doing it, at
$\Te = \SI{2.02}{\electronvolt}$:

\begin{center}
\begin{tabular}{lr}
  \toprule
  $\ngs$ scaling & density at which $|\ffrac{3}-\ffrac{4}|$ peaks \\
  \midrule
  $\times 0.01$ & $\le 10^{12}\,\si{\per\cubic\centi\metre}$ (off grid) \\
  $\times 0.1$  & $2.28\times10^{12}\,\si{\per\cubic\centi\metre}$ \\
  $\times 1$ (the model's own CRE value) & $1.68\times10^{13}\,\si{\per\cubic\centi\metre}$ \\
  $\times 10$   & $2.48\times10^{14}\,\si{\per\cubic\centi\metre}$ \\
  $\times 100$  & $\ge 10^{15}\,\si{\per\cubic\centi\metre}$ (off grid) \\
  \bottomrule
\end{tabular}
\end{center}

\noindent \textbf{One decade in the assumed neutral density moves the ridge by
roughly one decade in electron density.} The response is close to
proportional, and the whole three-decade span of the density grid is traversed
by a two-decade change in $\ngs$.

The agreement with Griem's criterion at unit scaling, $1.93\times10^{13}$
measured against $2.05\times10^{13}\,\si{\per\cubic\centi\metre}$ predicted, is
therefore an agreement evaluated at \emph{this model's} collisional--radiative
ionisation balance, which contains no transport and no recycling. In a real
divertor the neutral density is set by recycling, which is
Section~\ref{sec:transport}'s objection arriving in a different place. The
sentence ``the ridge sits at
$\nel \approx 2\times10^{13}\,\si{\per\cubic\centi\metre}$'' is a statement about
an ionisation balance, not about a divertor, and should be quoted with that
condition attached or not quoted.

The ridge is also one line pair's and not the Balmer series': the $(3,5)$ pair
of $H_\alpha/H_\gamma$ has its worst density at
$7.2\times10^{12}\,\si{\per\cubic\centi\metre}$, a factor $2.68$ lower, and the
form of words this imposes is stated in Section~\ref{sec:error_maps}.

%-----------------------------------------------------------------------------
\section{The model does not identify the crest with detachment}
\label{sec:not_detachment}
% QUESTION: can the density at which the ridge sits be identified with a named
% operating regime of a tokamak divertor?

\subsection{The question, and the temptation}

The identification is tempting. The ridge sits where ionising and recombining
supplies compete; detachment, the regime ITER must operate in to survive its
exhaust power and the transition this diagnostic is chiefly used to identify
(Section~\ref{sec:molecules}), is where ionising and recombining plasma compete.
The words match.

\subsection{The arithmetic}

The numbers do not. Section~\ref{sec:withdrawn} withdrew the identification:
the ridge, at $1.93\times10^{19}\,\si{\per\cubic\metre}$, sits a factor $5.2$
below the detached-divertor band of published edge
modelling~\cite{Guillemaut2011} and a factor $1.6$ below the modelled upstream
separatrix density~\cite{Stangeby2023,Pitts2019}, in a range neither source
characterises.

The comparison band deserves a word, because a factor of $1.6$ is not a large
margin to be arguing over. Read from the primary source, the band is
\SIrange{3e19}{6e19}{\per\cubic\metre} across the $53$ ITER SOLPS-4.3 cases
\cite[Sec.~3.1.1]{Pitts2019}, and it is the density \emph{at the outer
midplane}, which that paper states explicitly where it plots the dependence
\cite[Sec.~3.2.2]{Pitts2019}. So the quantity is settled: outer midplane, not
X-point. What is not settled by a single factor is where in the band to
measure from. The ridge sits $1.6$ times below the \emph{bottom} of it and
$3.1$ times below the top, and the honest statement is the range rather than
either end. On the same flux tube the outboard-midplane and X-point densities
differ by roughly \SI{50}{\percent}, which is comparable to the width of the
band itself. The conclusion does not turn on any of this, since the ridge is
below the band on every reading, but the factor should not be quoted more
precisely than the band it is measured against.
The repository recorded \cite{Stangeby2023} as ``Nucl.\ Fusion 62 (2022)'' in
\texttt{src/validation/grid.py}; the DOI was right, the volume and year wrong.
It is Nuclear Fusion \textbf{63}(1) 016017 (2023), companion part A at 63(1)
016016 \cite{Stangeby2023a}; \cite{Lore2022} in the same block of that file
genuinely is volume 62.

\subsection{What that leaves}

The honest statement is uncomfortable and should be made anyway. \textbf{The
structure this work found is real and reproducible, and it sits at a density
this thesis cannot identify with any named region of a tokamak.} It is not
detachment; it is not the upstream separatrix; and calling it either would be a
claim about edge transport that a $0$-D atomic model cannot support. The ridge
is a property of how two shells of hydrogen divide their supply. It needs no
divertor geometry to produce it, and it acquires none by being described in one.

The ridge and the detached band do not overlap: restricted to
$\Te \ge \SI{2}{\electronvolt}$ and $\nel \ge 10^{14}\,\si{\per\cubic\centi\metre}$
the census is $0$ of $108$ with a worst estimate of \SI{7.2}{\percent}
(Table~\ref{tab:trapping_census}), so where the diagnostic is most used the
transient error is real but modest.

%-----------------------------------------------------------------------------
\section{The open boundary: an attack that was expected to be fatal}
\label{sec:open_system}
% QUESTION: does holding the ion density fixed manufacture the slow clock that
% the whole result rests on?

\subsection{The question}

Section~\ref{sec:rate_equation} was explicit that the ion density $\nion$ is
held fixed: the recombination feed $\Svec\nion$ enters as a constant source term
a feed stream into a CSTR in the reactor picture, rather than as a
matrix element, and the $43$-state system has no $44$th state for ionised atoms
to accumulate in. Atoms ionise and vanish from the accounting. The slow
eigenvalue $\lambda_0$ is therefore the CR ionisation rate of ground-state
hydrogen against a prescribed reservoir, not the equilibration rate of a closed
system.

\subsection{The falsifier, and what happened}

If the long $\tauslow$ values were an artifact of the open boundary, so would be
the claim that the error outlasts an ELM. The falsifier named in
Section~\ref{sec:attacks} before the test was run is a closure factor of $2300$
at $[0,4]$, the value of $\tauslow/\taudrive$ there. Closing the manifold with
a $44$th ion state gives a closure factor with minimum $1.00$, median $1.00$
and maximum $41.4$ over the grid, moves the ELM census from $202$ to $200$ of
$680$ and the worst case by one part in $325$, and leaves the count above
\SI{2}{\electronvolt} at $45$ before and $45$ after (Table~\ref{tab:closure}).
The factor is large only where $\tauslow$ already exceeds the event duration by
three orders of magnitude and is unity wherever $\tauslow$ approaches
$\taudrive$, so where the boundary condition matters the outcome does not
depend on it. The measured factor at $[0,4]$ is $15$ against the $2300$
required, a cold-corner margin; inside the defended range, at the ridge, it is
$1.01$ against $20$, and at the benchmark $\tauslow$ moves by one part in
$10^{3}$. The attack was survived everywhere it was pressed.

\subsection{The caveat that must not be softened}

The closure adds an ion reservoir. It does \emph{not} close the electron
balance, since $\nel$ is held fixed while the ion population moves, and it adds
no transport whatever. It is a test of the boundary condition on the atomic
manifold, and nothing more. Section~\ref{sec:transport} is the objection this
test does not answer, and no amount of closing the atomic system will answer it.

The two index labels for this point, $[1,4]$ in Table~\ref{tab:closure} and
$[0,4]$ in the census, are the post- and pre-step operators of one heating
step: the census rows carry the pre-step index, every quoted timescale belongs
to the post-step operator that governs the relaxation (Chapter~\ref{ch:results}),
and the two must not be interchanged. Diagonalising each gives
$\tauslow = \SI{0.421012}{\second}$ at $\Lmat[0,4]$ and
\SI{0.233239}{\second} at $\Lmat[1,4]$, so the quoted \SI{0.23324}{\second}
belongs to $\Lmat[1,4]$.

%-----------------------------------------------------------------------------
\section{An external benchmark this model fails, and what the failure measures}
\label{sec:r1_deficit}
% QUESTION: where does the comparison against published atomic data break down,
% and does the failure touch the mechanism?

\subsection{The question}

Chapter~\ref{ch:validation} presented the comparison against Fujimoto's
published population coefficients as the strongest external check this model
passes. It is, for one of the two coefficients. The other disagrees at low
principal quantum number, and this section reports what that disagreement
measures. The short answer is that it is unexplained, that it is not the
$\ell$~closure, and that it is direct evidence about the coefficients
Chapter~\ref{ch:results} uses.

\subsection{The picture: two coefficients, one for each feed}

Chapter~\ref{ch:theory} split every excited population into two pieces
corresponding to the two supply channels. The population coefficient $r_0(p)$
measures how strongly level $p$ is fed from the ion continuum by recombination;
$r_1(p)$ measures how strongly it is fed from the ground state by excitation.
The published tabulation gives both, computed independently, decades ago, by a
different method. Agreement is a real test; disagreement is a real failure.

\subsection{The disagreement, and what it measures}

The recombination-fed coefficients agree. At
$\nel = 10^{12}\,\si{\per\cubic\centi\metre}$, $r_0$ matches to
\SI{1.3}{\percent} at $p=2$ and to \SIrange{12}{14}{\percent} at $p=3$ and
$p=4$. The ground-fed coefficients at low principal quantum number do not: this
model's $r_1$ sits a factor $8.3$ low at $p=3$ and a factor $4.4$ low at $p=4$.

Those are the two shells whose difference is the mechanism of
Chapter~\ref{ch:results}, so the disagreement had to be understood rather than
noted.

Two candidate explanations were eliminated before the source was consulted. If
the excitation rates feeding the ground-fed channel were wrong, the deficit
would follow, but those rates were benchmarked against
R-matrix-with-pseudostates calculations and the dominant $\Delta\ell = 1$
channels agree at a median of $-\SI{5.0}{\percent}$ over $120$ comparisons. And
a normalisation error is excluded by $r_0$ itself: $r_0$ and $r_1$ carry the
same $p$-dependent factor $Z(p)$, so an error large enough to move $r_1$ by a
factor of $8$ would have moved $r_0$ with it.

\subsection{The source's closure, imposed and tested}

Table~4.1(b) is indexed by $p = 2, 3, 4, 5, 7, 10, 15$. Those are principal
quantum numbers, and the table carries no $\ell$ label anywhere. The published
coefficients are shell quantities computed with the $\ell$~substructure
bundled.

This model resolves $\ell$ below $n=9$, so the closures differ. An earlier
draft of this section treated that difference as the explanation, on the
evidence that recomputing the model with proton $\ell$-mixing removed moves
$r_1(2)$ by a factor $118$ and $r_1(3)$ by a factor $3.35$ at
$\nel = 10^{12}\,\si{\per\cubic\centi\metre}$. That argument does not survive
its own test. Removing $\ell$-mixing is the limit in which the process is
absent; the tabulation assumes the limit in which it is infinitely fast, and
those are opposite ends rather than a bracket. Imposing the source's own closure
directly, by bundling this model's operator under statistical $\ell$~weights,
moves $r_1(3)$ by \SI{0.42}{\percent} and makes the $p=2$ agreement worse
(\texttt{verify\_fujimoto\_bundle.py}, stamped in
\texttt{validation/fujimoto\_bundle/}).

The disagreement therefore does \emph{not} measure the $\ell$~closure, and this
thesis does not explain it. An earlier draft of this chapter said the deficit
did not reach $\ffrac{3}-\ffrac{4}$ because that difference never passes
through a bundled quantity; that was wrong, and bundling is not the protection.
Section~\ref{sec:fujimoto} makes $a_p = r_1(p)\,\Zsb{p}/\Zsb{1}$, so a deficit
in $r_1(3)$ and $r_1(4)$ is a deficit in $a_3$ and $a_4$ themselves, and the
only protection is that $\Delta$ and the cap depend on the ratio $a_3/a_4$, in
which a deficit common to both shells cancels. It is not common at
\SI{e12}{\per\cubic\centi\metre}, $8.3$ against $4.4$, and nearly so at
\SI{e15}{\per\cubic\centi\metre}, $1.8$ against $1.9$. Dividing $a_3$ and
$a_4$ by those ratios at every grid point
(\texttt{validation/fujimoto\_r1\_consequence/}) moves $\Delta$ by $+0.64$ and
$-0.08$; over the $280$ points above \SI{2}{\electronvolt} the cap moves by
\SI{+29}{\percent} and \SI{-4}{\percent} at the median, and the sensitivity at
the operating point by \SI{+30}{\percent} and \SI{+11}{\percent} at the median,
with ranges of $-58$ to \SI{+524}{\percent} and $-38$ to \SI{+68}{\percent}.
The disagreement is direct evidence about the coefficients
Chapter~\ref{ch:results} uses; what is open is whether the deficit lies in this
model or in the tabulation, and the stake is stated in
Section~\ref{sec:fujimoto}. The open status is carried in
Section~\ref{sec:scope_open}.

\subsection{Two limitations of the benchmark, stated precisely}

Table~4.1(b) is computed at $T_\mathrm{e} = \SI{1.28e5}{\kelvin}$, which is
\SI{11.03}{\electronvolt}. The grid used here reaches \SI{10}{\electronvolt},
so the comparison is made at the grid edge and \SI{9.3}{\percent} below the
tabulated temperature.

There is no second usable case. The companion Table~4.1(a) is at
$T_\mathrm{e} = \SI{e3}{\kelvin}$, which is \SI{0.0862}{\electronvolt}, two
orders of magnitude below this grid. The external check is therefore one
temperature at one edge, and it should not be described as more than that.

\subsection{What the benchmark establishes, and what it does not}

What the benchmark does establish is the recombination-fed channel, which
agrees, and that is worth stating as a pass rather than leaving implicit beside
a disagreement. What it does not establish is the absolute normalisation of the
ground-fed channel against an independent calculation, because no tabulation
resolved in $\ell$ was available to compare against.

The sensitivity of the density maximum to the ground-fed supply remains a
separate and open question, and Section~\ref{sec:ridge_conditional} treats it
on its own terms: the location responds roughly decade-for-decade to the assumed
neutral density, which is a stronger reason to quote it as a range than
anything the benchmark shows.

%-----------------------------------------------------------------------------
\section{The top of the level ladder, and what the bundling costs}
\label{sec:nmax}
% QUESTION: the model stops at n=15 and bundles the levels above n=8 -- what
% rests on the part that is approximated?

\subsection{The question}

A hydrogen atom has infinitely many bound levels, crowding together towards the
ionisation limit. A computation cannot have infinitely many states, so a
boundary was chosen: this model tracks levels up to $n_{\max} = 15$ and
discards the rest. Above $n = 8$ it also stops resolving orbital angular
momentum $\ell$, treating each shell as a single bundled state populated
statistically. Both are choices, and this section states what they cost.

\subsection{Three consequences, in decreasing order of severity}

\paragraph{The terminal shell is over-populated.} Truncation removes from the
terminal shell's balance the collisional excitation out of it to $n > n_{\max}$,
a way out, and the de-excitation and cascade into it from those levels; an
earlier draft said the shell had ``nothing above it to cascade out to'', which
names the wrong process, since cascade is downward. Against published values
this model's ground-fed coefficient at $n = 15$ is high by a factor of $4.9$ to
$6.2$, while at $n = 10$ it agrees to \SIrange{4}{20}{\percent}; the budget is
in Section~\ref{sec:convergence}. \emph{No result in this thesis rests on the
top shell}: the observable is built from $n = 3$ and $n = 4$, both fully
$\ell$-resolved and far from the truncation.

\paragraph{The bundling nonetheless enters the result indirectly.} The
recombination flux cascading down from the bundled shells feeds the
recombination-fed channel of both $n=3$ and $n=4$, one of the two channels the
mechanism compares, which is why Section~\ref{sec:ridge_alternatives} lists
truncation among the alternative explanations for the ridge that this work does
not exclude.

\paragraph{There is no $\ell$-mixing inside the bundled block.} The
proton-impact $\ell$-mixing that Section~\ref{sec:lmixing} showed to be
essential for the $n = 2$ manifold is present only for $n = 2$--$8$; the
$\ell$-mixing rate array is identically zero above that. Within the bundled
shells the $\ell$-distribution is assumed statistical rather than computed. At
high $n$ that assumption is far more defensible than it would be at low $n$,
because the sublevels are nearly degenerate and every collisional process mixes
them, but it is assumed, not measured, and this thesis has not measured it.

\subsection{The check, run}
\label{sec:bundling_check}

The script that quantifies this was written and, until now, never executed.
Running it produced a result and also produced three defects in the script
itself, which are reported here because two of them would have made its verdict
worthless.

\paragraph{What the script could not do.} It promised to test the bundled
shells against the pipeline's own $\ell$-mixing array, and that is impossible:
the array is identically zero for states $36$ to $42$, in both directions,
because $\ell$-mixing was computed for the resolved block alone. The file
contains nothing to test. The comparison has to be made against the rate
\emph{function} rather than the rate table, and the function is the one
\texttt{compute\_lmix.py} used to build the resolved block, so it is the model's
own rate rather than a second one. Calling it at $n = 2$ to $8$ reproduces the
stored array to $1.0000$, which is a wiring check and not a physical
confirmation; it establishes only that the rate carried up to the bundled
shells is the same function with the same arguments.

\paragraph{The two defects that would have inverted the answer.} The script
carried its own Pengelly--Seaton expression, independent of the pipeline's,
which disagrees with the rate the model actually uses by a factor of $95$ at
$n = 2$ rising to $5.6\times10^{3}$ at $n = 8$, growing roughly as $n^{2}$. Used
as written, it would have reported the mixing rate two to three orders of
magnitude below its true value and could have returned a verdict of marginal
where the answer is not close. It is deleted rather than corrected: a
verification script has no business maintaining a second implementation of a
rate the pipeline already computes. The script also fell back on a synthesised
logarithmic temperature grid of its own when the pipeline's grid files were
absent,
printing a warning and continuing, which is the failure mode described in
Chapter~\ref{ch:validation}: every ratio would have carried a correct-looking
label attached to conditions the model was never evaluated at. It now raises.

\paragraph{The result.} The criterion is that the $\ell$-mixing rate out of the
$np$ substate, the substate the Lyman channel empties fastest, must exceed that
substate's total radiative rate. Over the seven bundled shells and all $400$
grid points, the ratio is above the well-mixed threshold of $10$ at every one of
the $2800$ combinations. The worst case is $n = 9$ at
$\Te = \SI{1}{\electronvolt}$, $\nel = \SI{e12}{\per\cubic\centi\metre}$, where
the mixing rate is $\SI{1.155e10}{\per\second}$ against
$A(9p) = \SI{7.51e6}{\per\second}$, a margin of $1538$, that is, a factor $154$
clear of the threshold (\texttt{verify\_bundling\_psm20.py},
\texttt{outputs/bundling/bundling\_report.txt}).

Two inputs to that comparison were checked rather than assumed. The Einstein
coefficients $A(np\to1s)$ for $n = 9$ to $15$ are hardcoded in the script,
because the model's own radiative file resolves $\ell$ only to $n = 8$ and has
nothing to read above it. Since $A(np\to1s)\,n^{3}$ tends to a constant for
hydrogen, the model's resolved values and the hardcoded table must lie on one
curve: they do, with a step of \SI{0.43}{\percent} across the boundary at
$n = 8$ to $9$ and a monotone approach to $4.13\times10^{9}$ at $n = 15$. And
the $\ell$-mixing coefficient carries a Debye cutoff frozen at
$\nel = \SI{e14}{\per\cubic\centi\metre}$, because the array has no density
axis. That freezing is not small: re-evaluating the cutoff at the local density
changes the $n = 15$ rate by a factor $20$ upward at
\SI{e12}{\per\cubic\centi\metre} and a factor $30$ downward at
\SI{e15}{\per\cubic\centi\metre}. Against the radiative criterion it does not
move the verdict, because the worst point sits at the lowest density, where the
local cutoff raises the rate: the worst ratio is $1538$ with the freezing and
$3.26\times10^{4}$ without it. At the highest densities the same correction
lowers the rate, and that matters for the criterion that follows.

\paragraph{What this closes and what it does not.} The criterion above compares
mixing with radiative decay, and radiative decay is not how a high shell
empties at these densities. The total loss out of a bundled state, $-L_{kk}$,
is \SIrange{89.5}{99.1}{\percent} collisional transfer to other shells,
\SIrange{0.9}{10.5}{\percent} ionisation and at most \SI{0.19}{\percent}
radiative (\texttt{verify\_bundling\_total\_loss.py},
\texttt{validation/bundling\_total\_loss/}). Against that total the margin is
not three orders of magnitude. With the frozen cutoff the model is built with,
mixing exceeds total escape at every one of the $2800$ combinations, by $26$
at $n=9$ falling to $2.85$ at $n=14$, the minimum, all at
\SI{1}{\electronvolt} and \SI{e12}{\per\cubic\centi\metre}, and $288$
combinations sit below $10$. With the cutoff evaluated at the local density
the ratio falls below unity at $121$ of the $2800$, in every bundled shell, at
the two highest densities and $\Te \le \SI{2.95}{\electronvolt}$; at
\SI{1}{\electronvolt} and \SI{e15}{\per\cubic\centi\metre} it runs from $0.86$
at $n=9$ to $0.09$ at $n=14$. The comment in \texttt{compute\_lmix.py} that the
cutoff factor varies by about \SI{30}{\percent} across the grid understates the
variation already at $n=2$, where the local-to-frozen ratio of $F$ runs from
$1.96$ at \SI{e12}{\per\cubic\centi\metre} to $0.53$ at
\SI{e15}{\per\cubic\centi\metre} at \SI{1}{\electronvolt}
(\texttt{validation/bundling\_total\_loss/}, RESULT~4), and fails outright
above $n=9$, where $U_m > 1$ and $F$ falls as $\nel^{-3/2}$; it is reported
here and left in place. What survives is narrower than an
earlier draft claimed. Where mixing wins, the $\ell$-distribution is
statistical whatever the feeds do. Where collisional transfer wins, at high
density, the distribution is set by the feeds, and $n$-changing collisions
between adjacent shells, which carry \SIrange{71}{95}{\percent} of the
collisional loss, are close to $\ell$-blind at high $n$, so the statistical
assumption is plausible there for a different reason, one this thesis has not
tested. It does not close the terminal-shell over-population above, which is a
different defect with a different cause, and it does not make the bundled
block's $\ell$-distribution a computed quantity. It remains assumed, now with
two numbers attached and a stated region where the first of them does not
license it.

%-----------------------------------------------------------------------------
\section{Remaining assumptions, and the sensitivities that are known}
\label{sec:other_assumptions}
% QUESTION: what else was assumed, and which of those assumptions has actually
% been tested?

Four assumptions remain. Each is stated with what is known about it, rather than
listed as a disclaimer, and two of the four have been tested.

\paragraph{Maxwellian electrons: untested.} Every rate coefficient in
Chapter~\ref{ch:atoms} is an average of a cross section over a Maxwell--Boltzmann
electron energy distribution. Ionisation and excitation are driven by the
high-energy tail of that distribution, which holds a small minority of the
electrons. A \SI{1}{\electronvolt} plasma next to a recycling surface has no
guarantee of a Maxwellian tail, and a depleted or enhanced tail would change the
rates that set both clocks. Not tested here.

\paragraph{$T_i = \Te$ in the $\ell$-mixing rates: untested.} The
proton-impact $\ell$-mixing rates of Section~\ref{sec:lmixing} are evaluated at
an ion temperature equal to the electron temperature. In a detaching divertor
the ions are frequently colder than the electrons. Not tested here.

\paragraph{The step is instantaneous: tested, and safe for a reason that is
easy to get wrong.} The error maps impose a sudden temperature step, whereas a
real ELM has a finite rise time. The tempting ratio
$\taurel/\taudrive \approx 10^{-5}$ is about the wrong quantity, the fast
excited-manifold step error (Section~\ref{sec:persistence}); the plateau error
is governed by the slow ground-state response, for which the relevant ratio is
$\tauslow/\taudrive$. At the worst point of the grid that
ratio is $2332$, but that point sits at \SI{1}{\electronvolt} and the figure is
not the margin this result runs on. Over the $448$ defended pairs the ratio has
median $0.815$, and over the $45$ that break down it runs from $1.95$ to $288$,
median $11.2$ (\texttt{validation/divertor\_map/}). The step idealisation is
therefore safe by one to two orders of magnitude where a plateau is reported,
not by three. The governing quantity is in any case
$D_{\rm e} = \tauslow/t_{\rm ramp}$ rather than $\tauslow/\taudrive$: a linear
ramp reaches $D_{\rm e}(1-\mathrm{e}^{-1/D_{\rm e}})$ of the step's plateau
(\texttt{validation/ramp\_plateau/}), so any ramp with
$\taurel \ll t_{\rm ramp} \ll \tauslow$ reaches the same plateau and the
idealisation \emph{is} safe. It is safe because of the slow ratio, and the fast
ratio must never be quoted for or against this result.

\paragraph{The step is in temperature only: tested, and the correction goes the
wrong way.} A real ELM raises $\nel$ as well as $\Te$. The joint map over $4968$
cases (Section~\ref{sec:joint_step}) leaves the exact decomposition unchanged,
with the gain generalising to an additive gradient accurate to
\SI{0.2}{\percent} at the median; over $\Te\ge\SI{2}{\electronvolt}$ the
breakdown fraction moves from $10.0$ to \SI{11.3}{\percent} when the density
steps up one index and the worst case from $0.1748$ to $0.2206$. The earlier
benchmark spot check, a bound falling by a factor $4$, had the right number and
the wrong sign for the grid, because the benchmark and the worst point lie on
opposite sides of the density crest; the fixed-density map understates the
worst case.

\medskip
Two sensitivities of the reported census belong in the same place, because both
are properties of a reporting choice rather than of the physics, and both would
otherwise be mistaken for measurements.

The step size is one interval of the temperature grid, which displaces the
reservoir by $|\ln\xscale|$ between $0.124$ and $0.682$ across the grid. A
four-interval step gives $0.516$ to $2.563$ and raises $\epsplat$ at the ridge
from \SI{18.1}{\percent} to \SI{81.1}{\percent}. On step size alone, therefore,
the reported numbers \emph{understate} the error rather than overstate it.

The \SI{10}{\percent} breakdown threshold is doing real work: at \SI{5}{\percent},
\SI{10}{\percent} and \SI{20}{\percent} the count over the same $680$ points is
$348$, $202$ and $61$ (Section~\ref{sec:persistence}), so the count is not robust
to the threshold, the worst case is, and ``$N$ points break down'' must carry both.

%-----------------------------------------------------------------------------
\section{Scope: five questions this thesis leaves open}
\label{sec:scope_open}
% QUESTION: which questions are open, why can none of them be closed from
% inside a 0-D atomic model, and what would close each?

The objections computed in this chapter were all answerable from inside a
$0$-D collisional--radiative matrix. Five are not, and they are gathered here
rather than left scattered so that a reader can see at once what this thesis
does not determine. Each is stated with the reason it is out of reach and with
the calculation that would settle it. None is a defect of the implementation;
each is a boundary of the model class.

\paragraph{1. The $r_1$ deficit is unexplained.} This model's ground-fed
population coefficient sits a factor $8.3$ below Fujimoto's Table~4.1(b) at
$p=3$, the shell the mechanism runs on, and a factor $1.8$ to $2.0$ below it at
$\log_{10}(\nel/\si{\per\cubic\metre}) = 21$; truncation (\SI{1.2}{\percent}),
Lyman trapping and the source's own statistical-$\ell$ closure
(\SI{0.42}{\percent}) have each been eliminated (Section~\ref{sec:r1_deficit}).
What would settle it is an $\ell$-resolved published tabulation or the atomic
data underlying the tabulated values, neither of which was available. The
deficit is reported as open, with its consequence for $\ffrac{3}-\ffrac{4}$
computed in Section~\ref{sec:fujimoto}.

\paragraph{2. There is no convergence study above $n_{\max} = 15$.} The state
space is cut at $n = 15$ because the atomic data stop there: the quantum
cross sections extend to $n = 10$, shells $11$ to $15$ are already semiclassical
Vriens--Smeets values, and no radiative or recombination coefficients exist
above $15$ in any form. Truncation has been tested downward, and it moves
$\ffrac{3}-\ffrac{4}$ by about \SI{1}{\percent} with the sign changing across
the grid (Section~\ref{sec:convergence}). What is untested is the sensitivity of
the ridge \emph{position} to $n_{\max}$, which would require rebuilding the rate
pipeline at $n_{\max} = 20$ rather than running a validation script: new
excitation, ionisation, radiative and recombination coefficients for five
further shells, and a re-indexing of every array built on the $43$-state layout.
That is a rate-generation exercise, not an analysis one.

\paragraph{3. The ground-state population is solved for, not supplied.} Every
magnitude in Chapter~\ref{ch:results} is conditional on $\ngs$ coming from
$\Lmat\nvec + \Svec\nion = 0$; in a divertor it is set by recycling, on a
transport timescale four orders of magnitude shorter than the local ionisation
balance (Section~\ref{sec:transport}). The open-boundary test of
Section~\ref{sec:attacks} bounds the sensitivity without removing the
conditional, and what would settle it is stated in Section~\ref{sec:transport}.

\paragraph{4. There are no molecular channels in $\Lmat$.} At the onset of
detachment a large fraction of the Balmer light is molecular in origin. The
dilution estimate of Section~\ref{sec:molecules} says how far the cold-corner
number would move under one specific dilution picture, not what a molecular
model would give. What would settle it is a state space containing
$\mathrm{H}_2$, $\mathrm{H}_2^+$ and $\mathrm{H}^-$ with their rate set, a
different model rather than an extension of this one.

\paragraph{5. The electrons are assumed Maxwellian.} The reasons are given in
Section~\ref{sec:other_assumptions}. The rate integrals as built do not accept
an arbitrary electron energy distribution, so this is not a parameter that can
be varied within the present code, and the $T_i = \Te$ assumption in the
$\ell$-mixing rates is untested for the same reason.

\medskip

\noindent Items 3 and 4 are named in Chapter~\ref{ch:conclusions} among the
further work. Items 1 and 2 are bounded in their consequences rather than in
their causes: the thesis states what each would change if it went the wrong
way, and for item 1 that stake is computed in Section~\ref{sec:fujimoto}.

%-----------------------------------------------------------------------------
\section{What is left standing}
\label{sec:limits_summary}
% QUESTION: after all of that, what may this thesis actually claim?

\subsection{What this model cannot say}

It cannot say anything quantitative about the cold edge of its own grid. Below
\SI{2}{\electronvolt} the optically thin plateau error is inflated by a factor
of $2.5$ to $3.3$ by neglected Lyman trapping; the size of the correction
depends on a slab thickness the model does not contain; and below
\SI{1.15}{\electronvolt} even the location of the ridge is undetermined.

It cannot say what a real divertor does, because recycling renews the
ground-state population there on a transport timescale the model does not
contain: \SI{72}{\micro\second} against \SI{233}{\milli\second} at the worst
point, a factor $3.2\times10^{3}$, below the scope boundary; inside the defended
range one to three orders rather than four, with a median $\tauesc/\tauslow$ of
$0.0079$ over the $45$ defended pairs (Table~\ref{tab:transport_selection}); and
$3.6$ at the benchmark, where transport is the \emph{slower} of the two.

It cannot say what an instrument records at the onset of detachment, because a
large fraction of that light is molecular in origin and there are no molecules
in $\Lmat$.

It cannot name the density at which its own ridge sits as a region of a tokamak.
The ridge falls between the modelled upstream separatrix and the detached
divertor band, and belongs to neither.

And it cannot establish the absolute normalisation of its ground-fed
population coefficient against an independent calculation: $r_1$ sits a factor
$8.3$ below the standard tabulation at $p = 3$, the disagreement is
unexplained, and its stake for $\ffrac{3}-\ffrac{4}$ is computed in
Section~\ref{sec:fujimoto} (Sections~\ref{sec:r1_deficit}
and~\ref{sec:scope_open}).

\subsection{What this model can say, with the objections computed}

Above \SI{2}{\electronvolt} the breakdown census is invariant under Lyman
trapping across a twentyfold range of assumed slab thickness, $45$ of $448$ in
every run with the worst case moving by \SI{0.5}{\percent}
(Section~\ref{sec:opacity}). Closing the open boundary moves the census by two
points out of $680$ and the worst case by one part in $325$; the factor $2300$
that would have refuted the result at the cold corner measures $15$, and at the
ridge the factor required is $20$ against a measured $1.01$
(Section~\ref{sec:attacks}). The ridge mechanism was verified over
$\Te \ge \SI{2}{\electronvolt}$, where trapping is now shown to be irrelevant,
so the agreement with Griem's boundary criterion is not an artifact of neglected
opacity. At citable divertor densities above \SI{2}{\electronvolt} the
\SI{100}{\micro\second} time-averaged error never reaches \SI{10}{\percent} over
$108$ points, worst case \SI{7.2}{\percent} (Section~\ref{sec:not_detachment}).

\subsection{The bridge}

The pattern is consistent, and it is why this chapter is not an apology. Every
objection that could be computed inside a $0$-D atomic model was computed, and
the defensible claim came through all of them with its numbers unchanged. The
objections that remain, neutral transport and molecular chemistry, are the
ones that cannot be computed inside a $0$-D atomic model at all, and they have
been stated as conditionals rather than dissolved into caveats.

That is the boundary of the claim. What is inside it is a quantitative statement
about when a spectroscopist may trust a lookup table and when she may not, and
Chapter~\ref{ch:conclusions} puts it into the form she can use.
 -- those labels are reused below, so every existing
%  cross-reference from Chapters 3 and 5 continues to resolve.
%
%  Sources for every number, in order of appearance:
%    outputs/findings_10_four_agent_review.md   ss 3.1, 4.1-4.4, 7.3, 8, 10,
%                                               and the ADDENDUM (A.1-A.6)
%    validation/lyman_trapping/lyman_trapping.txt    (run 9-10 Sep 2026,
%                                               src/validation/verify_lyman_trapping.py)
%    thesis_tex/chapter5_C5D.tex:284    (the "It is not detachment" wording)
%    outputs/findings_09_central_quantity_misnamed.md   ss 8, 9
%=============================================================================

\chapter{What This Model Cannot Say}
\label{ch:limits}

Chapter~\ref{ch:results} handed the reader two things: a map of how wrong a
tabulated $H_\alpha/H_\beta$ inversion can be while the ground state is still
catching up, and a mechanism that explains the map's shape. The map has a ridge
in it, the ridge has a closed-form explanation, and the explanation was tested
three separate ways. At this point in the argument the reader believes a
result.

This chapter's job is to bound that belief without destroying it. Every number
in Chapter~\ref{ch:results} came out of one $43$-state matrix describing one
kind of atom under one set of assumptions, and some of those assumptions are
badly violated in exactly the corner of the map where the numbers are largest.
An examiner will find that corner. It is better to arrive there first, with
arithmetic.

\medskip
\noindent\textbf{The standard this chapter holds itself to.} A limitation that
has been \emph{computed} is worth more than a limitation that has been hedged.
Where an objection could be turned into a calculation, it was, and the
calculation is reported here with the number that would have refuted the
result. Two objections were expected to be fatal and were not; they are given
the same space as the ones that bite. Where an objection cannot be computed
inside a model with no spatial dimension, it is stated as a \emph{conditional},
``\emph{if} this holds, the error is as computed'', and not softened into a
caveat sentence. Where a check has simply not been run, this chapter says so.

\medskip
\noindent\textbf{Three abbreviations recur below} and are expanded here for
this chapter: \emph{CR} for collisional--radiative, the class of model built in
Chapter~\ref{ch:theory}; \emph{QSS} for quasi-steady-state, the Bodenstein-style
approximation that eliminates the fast excited levels; and \emph{CRE} for
collisional--radiative equilibrium, the fully settled steady state that the
diagnostic lookup table assumes and that this thesis measures the cost of
assuming. An \emph{ELM} (edge-localised mode) is the repetitive burst of
heat and particles that a tokamak edge throws at the divertor every few
milliseconds, and it is the transient every ``event duration'' below refers to;
$\taudrive = \SI{100}{\micro\second}$ throughout, the duration used in
Chapter~\ref{ch:results} and the aggressive end of the range; the
ITER-pedestal extrapolation is \SI{506}{\micro\second}
\cite[Sec.~5]{Loarte2003} (Section~\ref{sec:persistence}).

%-----------------------------------------------------------------------------
\section{Optical thickness: when the plasma cannot let its own light out}
\label{sec:opacity}
% QUESTION: the model lets every emitted photon escape the plasma -- where is
% that false, what does it do to the error map, and how much of the map
% survives it?

\subsection{The question, and the picture behind it}

Chapter~\ref{ch:atoms} gave every downward transition an Einstein $A$
coefficient: the probability per second that an excited atom drops to a lower
level and emits a photon. Those coefficients entered the CR matrix $\Lmat$
bare. Entering bare means one specific physical assumption: that every photon
emitted anywhere in the plasma leaves it and never comes back. A plasma for
which this is true is called \emph{optically thin}.

A chemical engineer has met the failure of this assumption in another costume.
Consider a reaction whose product is removed by desorption from a catalyst
surface. If the product leaves the reactor, the desorption step runs at its
intrinsic rate. If the product cannot leave, if it re-adsorbs somewhere else
in the bed and is put back onto a site, then the \emph{effective} desorption
rate is smaller than the intrinsic one, and the whole network downstream of that
step shifts. Radiation trapping is that, with photons.

Here is the specific chain. The strongest transition in hydrogen is
Lyman-$\alpha$: an atom in the $2p$ level drops to the ground state $1s$ and
emits an ultraviolet photon. If the gas around it contains many ground-state
atoms, that photon has a good chance of being absorbed by one of them, which
puts \emph{that} atom back into $2p$. Nothing has left the plasma; a
$2p$ population has simply been moved from one place to another. Averaged over
the volume, the $2p$ level decays more slowly than its $A$ coefficient says. The
standard way to write this down is to multiply the coefficient by an
\emph{escape factor} $\Theta_P$ between $0$ and $1$: $\Theta_P = 1$ is a
transparent plasma, $\Theta_P = 10^{-5}$ is one in which the transition is
effectively switched off. The formulation used here is Holstein's population
escape factor~\cite{Holstein1947}.

\textbf{Why this reaches the populations, and therefore the result.} If $2p$
lives longer, atoms pile up in $n=2$. Atoms in $n=2$ are only
\SI{10.2}{\electronvolt} from the ground state but only
\SI{3.4}{\electronvolt} from the continuum, so they are far easier to ionise
than ground-state atoms. A longer $n=2$ lifetime therefore means more
\emph{stepwise} ionisation, ionisation via an excited level rather than
directly out of the ground state, which means fewer neutrals and a faster
turnover of the ground-state reservoir. The ground-state reservoir is the whole
mechanism of this thesis. Trapping is not a cosmetic correction to a line
intensity; it acts on the slow clock.

\subsection{Where the assumption fails, in numbers}

The formalism needed is one line. A photon travelling a distance $D$ through a
gas of absorbers is attenuated by $e^{-\tau}$, where the \emph{optical depth}
$\tau = \sigma_0\,\ngs\,D$ is dimensionless: $\ngs$ is the ground-state density
in \si{\per\cubic\centi\metre}, $D$ the path length in \si{\centi\metre}, and
$\sigma_0$ the absorption cross section at line centre in
\si{\square\centi\metre}. For a Doppler-broadened line,
$\sigma_0 = 2.654\times10^{-2}\,f/(\sqrt{\pi}\,\Delta\nu_D)$, with $f$ the
oscillator strength (a dimensionless measure of how strongly the transition
couples to light) and $\Delta\nu_D$ the Doppler width set by the temperature of
the absorbing atoms. $\tau \ll 1$ is transparent; $\tau \gg 1$ is opaque.

The number that matters is $\ngs$, not $\nel$, and Chapter~\ref{ch:theory}
already established the awkward fact that these two run in opposite directions:
the plasma is most neutral where it is coldest. The worst point of the error map
is $\Te = \SI{1.0}{\electronvolt}$,
$\nel = 5.18\times10^{13}\,\si{\per\cubic\centi\metre}$, the maximum of
$\epsplat$ over all $680$ window-resolvable points of the map, cold edge
included. There the model's own CR solution gives $\ngs/\nion = 28.3$, so
quasi-neutrality fixes
$\ngs = 1.47\times10^{15}\,\si{\per\cubic\centi\metre}$. At
\SI{1}{\electronvolt} that gives
$\sigma_0 = 5.47\times10^{-14}\,\si{\square\centi\metre}$, and the Lyman-$\alpha$
optical depth per centimetre of path is

\begin{center}
\begin{tabular}{lr}
  \toprule
  $\Te$ (eV) & $\tau_{\rm Ly\alpha}$ per \si{\centi\metre} at
              $\nel = 5.18\times10^{13}\,\si{\per\cubic\centi\metre}$ \\
  \midrule
  1.00 & 80.3 \\
  1.15 & 10.6 \\
  1.60 & 0.224 \\
  2.02 & 0.0255 \\
  2.95 & 0.00163 \\
  \bottomrule
\end{tabular}
\end{center}

\noindent An optical depth of $80$ per centimetre means a Lyman-$\alpha$ photon
travels about a tenth of a millimetre before it is reabsorbed. Over
\SI{5}{\centi\metre} of plasma the line-centre depth
over the half slab is $201$ and the population escape factor is
$1.2\times10^{-3}$; roughly one photon in eight hundred gets out. Chapter~%
\ref{ch:theory} quotes $3\times10^{-5}$ at a different grid point, a different
slab thickness and a different density, so the two are not the same quantity and
the agreement in order of magnitude should not be read as a reproduction of
one by the other.
The blunt reading is that the effective $A_{2p\to1s}$ belonging in $\Lmat$ at
that point is wrong, self-consistently, by about $1.8$ orders of magnitude, or
$2.9$ orders on the thin one-shot estimate
(\texttt{verify\_lyman\_optical\_depth.py}).

Two rows down the same table the problem has evaporated. At
\SI{2.02}{\electronvolt} the depth per centimetre is $0.0255$, and at the
benchmark point (\benchTe, \benchne) it is $0.0039$; a photon there crosses
metres of plasma before anything happens to it. \textbf{The measured quantity
is not ``how opaque is this plasma'' but ``how fast does opacity switch on as
the plasma cools'',} and the answer is a factor of $5\times10^{4}$ between
\SI{1}{\electronvolt} and \SI{2.95}{\electronvolt}.

That is why the results of Chapter~\ref{ch:results} were declared for
$\Te \gtrsim \SI{2}{\electronvolt}$ in the first place
(Section~\ref{sec:two_references}). It is also, uncomfortably, where the
headline was living. Recounting the \SI{100}{\micro\second} breakdown census against optical depth
rather than against temperature: of the $202$ points at which the
\SI{100}{\micro\second} time-averaged estimate of the error exceeds \SI{10}{\percent}, an extremum count over
the $680$ window-resolvable points of the full grid, $157$ lie below
\SI{2}{\electronvolt}, $85$ have
$\tau_{\rm Ly\alpha}(\SI{5}{\centi\metre}) > 1$, and $15$ have it above $100$,
on the half-slab depth used above.

An earlier draft gave $105$ and $26$, a convention error rather than a
recomputation: those counts are reproduced only by combining the full-path
depth $\kappa_0 D$ with a Doppler width $\sqrt{kT_n/m}$, short of the correct
$\sqrt{2kT_n/m}$ by $\sqrt{2}$. With the correct width the census is $85$ and
$15$ on the half-slab depth quoted here and $100$ and $23$ on the full-path
depth. A deuterium cross-section with the full-path depth also returns $105$
and $26$, since the isotope mass enters the width as the same $\sqrt{2}$, so
the two errors are indistinguishable from the counts alone. All seven
conventions are tabulated in \texttt{validation/lyman\_optical\_depth/}.
Nothing in the defended range depends on the choice: every crossing lies below
\SI{2}{\electronvolt} except two, both at $D = \SI{20}{\centi\metre}$.

\subsection{What we expected trapping to do, and what it did}

The expectation, stated before the calculation was run, was that trapping would
demolish the map: an effective $A$ coefficient wrong by about three orders of
magnitude at the worst point ought to invalidate everything computed there and
nearby, and the project's own records warned that under strong trapping the
fast relaxation time $\taurel$ rises to \SI{2.05}{\micro\second} and the
timescale separation $\Msep$ falls by a factor $158$. That expectation was
wrong.

So the calculation was done. The Lyman series was given a self-consistent
Holstein escape factor and the entire error map rebuilt. Self-consistent matters
here: $\Theta_P$ depends on $\ngs$, and $\ngs$ depends on $\Theta_P$ through the
CR balance, so the two were iterated to a fixed point rather than evaluated once.
The fixed point converged at all $400$ grid points for every slab thickness
tested, and non-convergence was made to raise an error rather than silently
return the last iterate. All $14$ Lyman channels were trapped, the seven
resolved $n\mathrm{P}\to1\mathrm{S}$ transitions and the seven bundled ones,
not Lyman-$\alpha$ alone.

Three gates were passed before any result was looked at. The oscillator
strengths were derived from this repository's own $A$ coefficients by inverting
the Einstein relation rather than imported, and come out as $f = 0.4162$,
$0.0791$, $0.0290$, $0.0139$ for Lyman-$\alpha$ through Lyman-$\delta$, the
accepted literature values to four figures. The derived Lyman-$\alpha$
$\sigma_0$ agrees with an independent implementation to \SI{0.059}{\percent},
after a first run that failed at \SI{1156}{\percent} on a $4\pi$ factor; and
rebuilding the matrix with trapping switched off reproduces the canonical
$\Lmat$ exactly, so every trapped number below is a difference against the real
matrix (Section~\ref{sec:severe_checks}).

\textbf{One new quantity had to be invented to do this at all, and it is not in
the model.} The escape factor depends on how far a photon must travel to get
out, so it depends on the size of the plasma, the slab thickness $D$. A model
with no spatial dimension has no $D$. It was therefore swept, from
\SI{1}{\centi\metre} to \SI{20}{\centi\metre}, and never fixed; every trapped
number below carries the $D$ it was computed at. Reporting a single trapped
number without its $D$ would be reporting an assumption as a measurement.

\begin{table}[tbp]
  \centering
  \caption[Breakdown census under Lyman trapping]{Effect of radiation
  trapping on the \SI{100}{\micro\second} breakdown census. Each entry counts the grid points at
  which the time-averaged estimate of the CR-equilibrium-reference error,
  averaged over a chosen duration $\taudrive = \SI{100}{\micro\second}$, exceeds
  \SI{10}{\percent}, out of the points at which a
  quasi-steady-state window is resolvable; ``worst'' is the largest estimate
  within that scope. Rows differ only in the assumed slab thickness $D$ used to
  compute the self-consistent Holstein escape factor applied to all $14$ Lyman
  transitions. The first row is the optically thin matrix used throughout
  Chapter~\ref{ch:results}. Above \SI{2}{\electronvolt} the count does not move
  at all across a twentyfold range of $D$. The $D = \SI{20}{\centi\metre}$ row
  has a slightly smaller denominator because two points lose their resolvable
  window under trapping. Source:
  \texttt{validation/lyman\_trapping/lyman\_trapping.txt}.}
  \label{tab:trapping_census}
  \footnotesize
  \begin{tabular}{lccc}
    \toprule
    matrix & all window-resolvable & $\Te \ge \SI{2}{\electronvolt}$ &
          $\Te \ge \SI{2}{\electronvolt}$, $\nel \ge 10^{14}$ \\
    \midrule
    optically thin (canonical)          & $202/680$, worst $0.3868$ & $45/448$, worst $0.1748$ & $0/108$, worst $0.0717$ \\
    trapped, $D = \SI{1}{\centi\metre}$  & $176/680$, worst $0.2925$ & $45/448$, worst $0.1748$ & $0/108$, worst $0.0712$ \\
    trapped, $D = \SI{5}{\centi\metre}$  & $170/680$, worst $0.2550$ & $45/448$, worst $0.1746$ & $0/108$, worst $0.0694$ \\
    trapped, $D = \SI{20}{\centi\metre}$ & $160/678$, worst $0.2237$ & $45/446$, worst $0.1740$ & $0/106$, worst $0.0627$ \\
    \bottomrule
  \end{tabular}
\end{table}

\subsection{Above two electron-volts, nothing moves}

Table~\ref{tab:trapping_census} contains the strongest robustness statement in
this thesis, and it is the opposite of what was expected.

\textbf{Above \SI{2}{\electronvolt} the breakdown count does not change at all,
$45$ of $448$ points in every run, and the worst case within that set
moves by \SI{0.5}{\percent}, from $0.1748$ to $0.1740$, across a twentyfold
range in assumed slab thickness.} The escape factor at the benchmark point
(\benchTe, \benchne) runs from $0.9986$ at $D = \SI{1}{\centi\metre}$ to
$0.973$ at $D = \SI{20}{\centi\metre}$. Even at the thickest slab tested,
\SI{97}{\percent} of Lyman-$\alpha$ photons still leave: the plasma is optically
thin there in the sense that actually matters, and assuming so costs less than
the resolution of the map.

This converts a disclaimer into a measurement. The
$\Te \gtrsim \SI{2}{\electronvolt}$ restriction was chosen on an escape-factor
argument before any of this was recomputed, and it sits almost exactly at the
temperature above which trapping stops mattering: a guess that has now been
checked, and survived. The defensible claim of Chapter~\ref{ch:results}, $45$
of $448$ points above \SI{2}{\electronvolt} at \SI{100}{\micro\second} and $24$
of $448$ at \SI{506}{\micro\second} (Section~\ref{sec:persistence}), is
invariant under the one physical process most likely to have invalidated it.

\subsection{Below two electron-volts, the cold-edge number is inflated
threefold}

The same calculation is much less kind to the cold edge, and it should be,
because that is where the optical depth was $80$ per centimetre.
Table~\ref{tab:trapping_worstpoint} follows the worst point of the optically
thin map through the sweep.

\begin{table}[tbp]
  \centering
  \caption[Lyman trapping at the cold-edge worst point]{The cold-edge worst
  point of the optically thin error map, at $\Te = \SI{1.0}{\electronvolt}$ and
  $\nel = 5.18\times10^{13}\,\si{\per\cubic\centi\metre}$ under a heating step,
  is the maximum of $\epsplat$ over all $680$ window-resolvable grid points; it
  is traced here as a function of the assumed slab thickness $D$. $\epsplat$ is the plateau
  error, the fractional error a lookup table makes while the ground state is
  stale; the ELM bound is its average over an event of
  $\taudrive = \SI{100}{\micro\second}$. The plateau error falls by a factor
  $2.5$ to $3.3$, but the diagnostic's intrinsic sensitivity
  $|\ffrac{3}-\ffrac{4}|$ barely moves: what collapses is how far the ground
  state is displaced, not how strongly the line ratio responds to it. Source:
  \texttt{validation/lyman\_trapping/lyman\_trapping.txt}.}
  \label{tab:trapping_worstpoint}
  \small
  \begin{tabular}{lrrrrrrr}
    \toprule
    $D$ (cm) & $\epsplat$ & ELM bound & $\tauslow$ (s) & $\Msep$ &
    $\ffrac{3}$ & $\ffrac{4}$ & $|\ffrac{3}-\ffrac{4}|$ \\
    \midrule
    $0$ (thin) & $0.38690$ & $0.38682$ & $2.3324\times10^{-1}$ & $3.74\times10^{7}$ & $0.7728$ & $0.2862$ & $0.4866$ \\
    $1$        & $0.15748$ & $0.15739$ & $9.0651\times10^{-2}$ & $2.75\times10^{6}$ & $0.6643$ & $0.2055$ & $0.4589$ \\
    $5$        & $0.13274$ & $0.13260$ & $4.6940\times10^{-2}$ & $6.13\times10^{5}$ & $0.5967$ & $0.1793$ & $0.4174$ \\
    $20$       & $0.11570$ & $0.11550$ & $2.8500\times10^{-2}$ & $2.14\times10^{5}$ & $0.5139$ & $0.1501$ & $0.3639$ \\
    \bottomrule
  \end{tabular}
\end{table}

\textbf{The \SI{38.7}{\percent} becomes \SIrange{11.6}{15.7}{\percent}}, a
factor of $2.5$ to $3.3$, and which value inside that range it takes is set
by a slab thickness the model does not contain. The point still exceeds the
\SI{10}{\percent} threshold at every thickness tested, so the qualitative
statement that the cold corner breaks down survives intact; the \emph{number}
does not, and it must never appear without its trapped range and without the
optical depth that produced it.

The route by which the number falls is itself a result, and it is not the
obvious route. The obvious guess is that trapping desensitises the diagnostic,
that a trapped plasma simply responds less. It does not. The sensitivity
$|\ffrac{3}-\ffrac{4}|$, which is the factor in the exact decomposition of
Section~\ref{sec:mechanism} measuring how strongly the line ratio answers to its
reservoir, moves by only \SI{6}{\percent} at $D = \SI{1}{\centi\metre}$ and
\SI{25}{\percent} at $D = \SI{20}{\centi\metre}$. What collapses is the other
factor: the displacement.

The chain is the one set out at the start of this section, now with numbers on
it. Trapping lengthens the $n=2$ lifetime; the longer lifetime raises stepwise
ionisation; the ground-state population falls; and the slow clock $\tauslow$
shortens by a factor $8$ at the worst point. At $\Te = \SI{1.0}{\electronvolt}$,
$\nel = 1.93\times10^{13}\,\si{\per\cubic\centi\metre}$, $\ngs$ drops from
$2.99\times10^{14}$ to $7.26\times10^{13}\,\si{\per\cubic\centi\metre}$ between
$D = 1$ and \SI{20}{\centi\metre}. A reservoir that turns over faster is a less
stale reservoir, so $|\Delta\ln\xscale|$ roughly halves. This is exactly the
two-axis attribution of Section~\ref{sec:temperature_mechanism} arriving from a
third direction: trapping acts on the displacement, not on the sensitivity.

\subsection{Below 1.15 electron-volts the ridge location is not determined}

One more thing moves at the cold edge, and it is a location rather than a
magnitude. Taking the density at which $\epsplat$ is maximal, row by row: at
$\Te \ge \SI{1.6}{\electronvolt}$ that maximum falls on
$\nel = 1.93\times10^{13}\,\si{\per\cubic\centi\metre}$ in every run, optically
thin and trapped alike, and the values themselves barely register the change:
at the benchmark temperature $\epsplat = 0.12166$, $0.12166$, $0.12165$,
$0.12163$ for $D = 0, 1, 5$ and \SI{20}{\centi\metre}, with
$|\ffrac{3}-\ffrac{4}|$ changing in the fifth decimal place. Below
\SI{1.15}{\electronvolt} the maximum wanders across three density columns, a
factor of $7$, as a function of an assumed parameter. \textbf{The two coldest
rows of the grid cannot support a claim about where the ridge is at all.}

That result cuts in the mechanism's favour, which is worth saying explicitly
because it is easy to miss. The ridge attribution of
Section~\ref{sec:mechanism} was verified over $\Te \ge \SI{2}{\electronvolt}$,
and $\Te \ge \SI{2}{\electronvolt}$ is precisely the range in which trapping has
now been shown to be irrelevant. The agreement with Griem's boundary criterion
is therefore not an artifact of neglected opacity; it was established where
opacity does not act.

\subsection{A contradiction resolved, and what is still open}

Two statements about trapping had been circulating in this project's records:
that $\taurel$ is insensitive to Lyman-$\alpha$ trapping to within
\SI{1}{\percent}, and that under strong trapping $\taurel$ rises to
\SI{2.05}{\micro\second} with $\Msep$ falling $158$-fold. They look
contradictory. Measured directly, both are true and neither had stated its
scope: at the benchmark point (\benchTe, \benchne) $\taurel$ moves from
$2.2769\times10^{-9}$ to $2.2947\times10^{-9}\,\si{\second}$ at
$D = \SI{20}{\centi\metre}$, a change of \SI{0.78}{\percent}; at
$\Te = \SI{1.0}{\electronvolt}$,
$\nel = 1.93\times10^{13}\,\si{\per\cubic\centi\metre}$ it moves from
$8.8842\times10^{-9}$ to $2.4289\times10^{-7}\,\si{\second}$, a factor of $27.3$.
The ``within \SI{1}{\percent}'' claim is a benchmark-point claim and fails by
more than an order of magnitude at the cold corner. This is the same failure
mode as quoting any other extremum without the set it is an extremum over, and
neither sentence may travel again without its regime attached.

What this calculation does not settle should be as clear as what it does. The
geometry is assumed, not derived: a homogeneous slab, uniform emission, an
isotropic radiation field, a Doppler line profile, evaluated at the plasma
centre. A $0$-D model cannot check any of those. One assumption inside it
\emph{was} checked and is not load-bearing: the temperature of the absorbing
atoms, which sets the Doppler width. Fixing it at \SI{3}{\electronvolt}, the
value a molecule dissociating by the Franck--Condon process would give, instead
of at $\Te$, changes the census at $D = \SI{5}{\centi\metre}$ from $170/680$ and
a worst case of $0.2550$ to $173/680$ and $0.2653$. Continuum lowering is
untouched here and remains open. Opacity in the Balmer lines themselves, which
would act on the observed ratio directly rather than through the level
populations, has been computed: the largest $H_\alpha$ optical depth above
\SI{2}{\electronvolt} is $0.010$, and the worst cold cell, $[0,7]$, reaches
$0.214$ (Section~\ref{sec:balmer_equivalence},
\texttt{validation/balmer\_optical\_depth/}).

%-----------------------------------------------------------------------------
\section{Transport: the atoms do not stay where the model puts them}
\label{sec:transport}
% QUESTION: the slow clock is the ionisation time of a stationary atom -- what
% happens to the result when the atom leaves before it ionises?

\subsection{The question}

This is the objection that cannot be answered inside this model, and it is a
more serious threat to the cold-edge number than anything in
Section~\ref{sec:opacity}. It is also the simplest to state.

\subsection{The picture: a reactor with no walls}

The model of Chapter~\ref{ch:theory} is a batch reactor. A fixed parcel of gas
is given a temperature, and its composition evolves until it stops changing. The
slow eigenvalue of that system, $\tauslow$, is how long the ground-state
population takes to reach its new value, and it is long, because ionising an
atom out of the ground state is slow.

A divertor is not a batch reactor. It is closer to a continuous-flow reactor
with an enormous throughput, in which the residence time of a molecule is far
shorter than the time it needs to react. Neutral atoms are born at the target
surface when ions strike it and recombine there, which is \emph{recycling},
and they fly back into the plasma. They do not wait around to be ionised
locally; the ones that do not ionise leave, and fresh ones arrive continuously.

The comparison is not close. At the worst point of the map $\tauslow$ is
\SI{233}{\milli\second}: nearly a quarter of a second for a stationary atom to
be ionised. A deuterium atom at \SIrange{1}{3}{\electronvolt} would cross a
\SI{10}{\centi\metre} plasma ballistically in \SIrange{6}{10}{\micro\second}, but
it does not travel ballistically: resonant charge exchange with the surrounding
ions randomises its direction after about \SI{1.7}{\centi\metre} at this
density, so escape is diffusive and takes closer to \SI{72}{\micro\second}. That
correction runs in this model's favour. The threshold at which the worst point
would fall below \SI{10}{\percent} is a ground-state renewal time of about
\SI{26}{\micro\second}, so the result survives by a factor of $2.8$ rather than
failing outright. It survives only because charge exchange traps the neutrals,
which is a thin margin to rest on and is stated as one.
That is one grid point, and it is one this thesis does not defend:
$[0,4]$ lies below the \SI{2}{\electronvolt} scope boundary.
Section~\ref{sec:transport_selection} asks the same question at the points that
do carry the result, and gets a worse answer.

\begin{quote}
\textbf{In a divertor, the ground-state neutral density at a given point is not
set by local ionisation balance. It is set by recycling.}
\end{quote}

\subsection{Why this bears on the result specifically}

The error of Chapter~\ref{ch:results} exists for one reason: after a temperature
step, the ground-state reservoir is \emph{stale}. Its population has not moved
while the temperature has, the excited levels equilibrate against a reservoir
that belongs to the old conditions, and the lookup table, which assumes the
reservoir finished settling, reads the wrong answer. Everything hangs on the
reservoir not being refreshed.

Neutral transport refreshes it. Any process that replaces the ground-state
population on a timescale shorter than $\tauslow$ resets the staleness, and this
model contains no such process; the \SI{26}{\micro\second} renewal threshold
quoted above is the map's own measure of how fast that process would have to be.

\subsection{The assumption fails hardest where the result is generated}
\label{sec:transport_selection}

The worst point was the wrong place to test transport: $[0,4]$ lies below the
scope boundary of Section~\ref{sec:scope}. The pairs that carry the defended
claim are the $45$ above \SI{2}{\electronvolt} whose \SI{100}{\micro\second}
estimate exceeds \SI{10}{\percent} (Table~\ref{tab:elm_census}), and there the
answer is different.

Comparing the neutral escape time with $\tauslow$ pair by pair over the grid
gives Table~\ref{tab:transport_selection}. The escape model is the one used
above, stated as a formula: the atom moves at $v=\sqrt{2k\Tn/m_{\mathrm D}}$
with $\Tn=\Te$, resonant charge exchange with the ambient ions randomises its
direction after $\lambda = v/(\nel\langle\sigma v\rangle_{\mathrm{CX}})$ with
$\langle\sigma v\rangle_{\mathrm{CX}} = \SI{1.1e-8}{\cubic\centi\metre\per\second}$,
and escape from a boundary a distance $L$ away takes $L/v$ when
$\lambda \ge L$ and $4L^{2}/\pi^{2}D$ with $D = v\lambda/3$ otherwise. At
$[0,4]$ it returns $\lambda = \SI{1.72}{\centi\metre}$ and
$\tauesc = \SI{72.3}{\micro\second}$, the values used above
(\texttt{verify\_transport\_selection.py}).

\begin{table}[htbp]
\centering
\caption{Ratio of neutral escape time to $\tauslow$, over the $45$ defended
breakdown pairs and over the other $635$ analysed pairs. $L$ is the distance
from the emitting parcel to the plasma boundary. ``Satisfies'' means
$\tauesc > \tauslow$, that is, the atom is ionised before it leaves, which is
what the closed parcel assumes.}
\label{tab:transport_selection}
\begin{tabular}{llrrr}
  \toprule
  $L$ & set & median $\tauesc/\tauslow$ & largest & satisfies \\
  \midrule
  \SI{10}{\centi\metre} & the 45 breakdown pairs & $0.0079$ & $0.159$ & $0/45$ \\
  \SI{10}{\centi\metre} & the other 635 & $0.0471$ & $384$ & $168/635$ \\
  \SI{20}{\centi\metre} & the 45 breakdown pairs & $0.0192$ & $0.636$ & $0/45$ \\
  \SI{20}{\centi\metre} & the other 635 & $0.129$ & $1535$ & $227/635$ \\
  \bottomrule
\end{tabular}
\end{table}

\textbf{Not one of the $45$ pairs that carry the defended result satisfies the
assumption used to compute it, at either scale length, while a quarter to a
third of the pairs that do not break down satisfy it comfortably.}

A count of zero is worth little if the $45$ sit just under the line, so the
distance matters. The closest of them is $[18,4]$ heating,
$\Te = \SI{2.33}{\electronvolt}$, $\nel = \SI{5.18e13}{\per\cubic\centi\metre}$,
at $\tauesc/\tauslow = 0.636$ for $L = \SI{20}{\centi\metre}$: its escape time
would have to rise by a factor $1.57$ for the parcel to close. The median of the
$45$ would have to rise by a factor $52$. One pair is marginal and the rest are
not.

Three things were varied to see whether the conclusion depends on a choice.
Tripling $\langle\sigma v\rangle_{\mathrm{CX}}$, which lengthens every diffusive
escape time in proportion, moves the count from $0$ to $3$ of $45$; dividing it
by three leaves it at $0$. Replacing $\sqrt{2k\Tn/m}$ by the Maxwellian mean
speed $\sqrt{8k\Tn/\pi m}$ shortens every escape time and leaves the count at
$0$. Doubling $L$ from $10$ to \SI{20}{\centi\metre} leaves it at $0$. The
finding is not an artifact of any of the three.

\paragraph{The selection is structural, not incidental.} A pair enters the
census because $\tauslow$ is long, since the single-slow-mode estimate of
Eq.~\eqref{eq:elm_bound} rises with $\tauslow$ at fixed $\taudrive$. A long
$\tauslow$ means slow ionisation out of the ground state. Slow ionisation is
exactly the condition under which neutral transport, whose timescale does not
depend on $\tauslow$ at all, sets the ground-state density instead of the local
atomic physics. The census therefore selects for the regime in which its own
assumption is weakest, and it does so by construction rather than by accident.
That is why testing the assumption at the cold corner gave the reassuring answer
and testing it inside the defended scope does not.

\subsection{What this reaches, and what it does not}
\label{sec:transport_partition}

The finding above is about magnitudes; it does not touch the algebra.

\textbf{Untouched by transport.} The two-channel split of
Eq.~\eqref{eq:two_supplies}, exact to a relative residual of
$3.075\times10^{-14}$; the unit-width logistic form of the ground-fed fraction;
the bound $\max|f_3-f_4| = \tanh(|\Delta|/4)$ and its corollary
$|\mathrm{d}\ln\Robs/\mathrm{d}\ln\bdep| < 1$; the sensitivity map
$\overline{S}$; and the accuracy of the excited-state closure,
$6.73\times10^{-9}$ at the worst point. These are properties of the operator
$\Lmat(\Te,\nel)$ at a fixed operating point, or theorems about its structure.
$\overline{S}$ says how the line ratio responds to whatever the reservoir does;
it is silent on what made the reservoir do it. Transport changes the reservoir's
history and cannot reach any of them.

\textbf{Conditional on the closed parcel.} The reservoir gain $G$, which is
a secant along the CRE curve and therefore assumes the ground
state responds through collisional--radiative balance alone; $\epsplat$, which
is built from it; the breakdown census; and every time-averaged magnitude in
Chapter~\ref{ch:results}. If recycling renews $\ngs$ on a timescale shorter than
$\tauslow$, $G$ is not the gain the plasma has. These numbers are not upper
estimates. A fast-exchanging reservoir with a steady source freezes the
closed-parcel error rather than removing it, and a rising source enlarges it: over
the defended set the true error exceeds the closed-parcel plateau at $165$ of the
$180$ (pair, $L$, source-multiplier) combinations tested
(\texttt{verify\_open\_reservoir.py}). The closed-parcel value is a reference
magnitude conditional on the reservoir's history, not a bound on it.

The thesis therefore leads with the structure and the bound and reports the
magnitudes as conditional; Section~\ref{sec:transport_selection} adds that the
condition is measured to fail at the points that matter, by one to two orders
of magnitude.

\subsection{The claim, restated as the conditional it is}

The honest form is therefore not a caveat but a conditional, and it should be
written this way wherever the number appears:

\begin{quote}
\emph{If} the ground-state neutral density is frozen for the duration of the
event, the error is as computed.
\end{quote}

That conditional is close to exactly true for the excited manifold, which
equilibrates in nanoseconds, and it is the standard implicit assumption of every
$0$-D collisional--radiative table in use, the object this thesis set out to
test. But it is an assumption about the plasma, not about the atoms, and no
eigenvalue analysis inside a $0$-D matrix can validate it: the two timescales
are measured to differ by four orders of magnitude and assumed not to interact.
What would settle it is a two-region or transport-coupled calculation, an edge
code such as SOLPS (Scrape-Off Layer Plasma Simulation) supplying $\ngs$ from a
recycling influx; that is outside the scope of this thesis and is named in
Chapter~\ref{ch:conclusions} among the further work.

%-----------------------------------------------------------------------------
\section{No molecular channels, in a gas that is mostly molecules}
\label{sec:molecules}
% QUESTION: the model contains only atomic hydrogen -- at the conditions where
% the error is largest, is that a defensible description of the gas?

\subsection{The question, and why a referee asks it first}

Every state in Chapter~\ref{ch:atoms}'s state space is a level of a single
hydrogen \emph{atom}, plus a reservoir of bare protons. There is no
$\mathrm{H}_2$ (or $\mathrm{D}_2$), no molecular ion $\mathrm{D}_2^+$, no
negative ion $\mathrm{D}^-$, and no reactions between them. A search of the rate
assembly for any of these returns nothing at all.

\subsection{The picture: a third feed stream}

Chapter~\ref{ch:theory} fed every excited level from two directions, up from
the ground state by electron collisions and down from the ion continuum by
recombination, and the error mechanism is the difference in how the $n=3$ and
$n=4$ shells split their supply between those two feed streams.

Molecules open a third stream. In a cold, dense, mostly neutral gas, hydrogen
atoms form $\mathrm{H}_2$; the molecules are ionised to $\mathrm{H}_2^+$, or
exchange charge, and the resulting molecular ions recombine with electrons and
break apart, leaving one of the fragments in an \emph{excited} atomic level. The
whole family of such routes is called MAR, molecular activated recombination.
The important feature for this thesis is not that MAR exists but where it
delivers: it populates $n=3$ directly, without the atom ever having passed
through the ground-state reservoir whose staleness is the entire mechanism here.
A supply channel that bypasses the reservoir does not inherit the reservoir's
lag.

\subsection{How much gas is involved}

At the worst point of the map the model's own equilibrium solution is
\SI{96.6}{\percent} neutral: $\ngs = 1.47\times10^{15}$ against
$\nel = 5.18\times10^{13}\,\si{\per\cubic\centi\metre}$. A hydrogen gas at
\SI{1}{\electronvolt} that is \SI{96.6}{\percent} neutral is a gas in which
molecules form readily. The atomic description is not being applied at the edge
of its validity there; it is being applied outside it.

The size of what is being left out is documented experimentally.
\emph{Detachment} is the operating regime in which the plasma cools enough that
the ion flux to the target collapses, and the regime this diagnostic is chiefly
used to identify. In spectroscopic inversions of divertor Balmer emission at its
onset, molecularly-induced emission is found to account for up to
\SIrange{60}{70}{\percent} of the $D_\alpha$ light and
\SIrange{10}{20}{\percent} of $D_\gamma$~\cite{Wijkamp2023,Verhaegh2021}. An
$H_\alpha/H_\beta$ ratio computed from atomic processes alone is not the ratio
such an instrument records.

\subsection{What survives it, and what does not}

Two consequences follow, and they point in opposite directions, so both must be
stated.

The \emph{numerical} cold-edge results are atomic-model results and nothing
more. Adding a third supply channel that feeds $n=3$ preferentially would change
$\ffrac{3}-\ffrac{4}$, and could plausibly change it a great deal; this work
offers no direct calculation of how much. A bound can nonetheless be
constructed, from the emissivity fractions cited in this section. For a given
upper level the molecular fraction of the emissivity equals the molecular
fraction of the population, because both channels emit through the same
transition probability, and molecular-assisted recombination bypasses the ground
state entirely, so it dilutes the ground-fed fraction as
$\ffrac{p} \to \ffrac{p}(1-\phi_p)$. Taking $\phi_3 = 0.60$ and $\phi_4 = 0.30$
from the reported \SIrange{60}{70}{\percent} of $\mathrm{D}\alpha$ and
\SIrange{10}{20}{\percent} of $\mathrm{D}\gamma$, the difference of ground-fed
fractions at the cold corner falls from $0.4866$ to $0.1088$ and at the density
crest from $0.4580$ to $0.0998$, and the secant sensitivity $\overline{S}$
falls with it, from $0.4822$ to $0.0903$ and from $0.4552$ to $0.0903$.

The plateau error does not fall in the same proportion, because
$\epsplat = |\mathrm{e}^{\overline{S}\Lambda}-1|$ is not linear in
$\overline{S}$; treating it as linear understates the reduction. Imposing the
same fractions on the model's own operator, with a non-negative source
depositing into $n=3$ and $n=4$ so that $(\phi_3,\phi_4)$ is reached exactly at
the plateau, the cold-corner plateau error falls from \SI{38.7}{\percent} to
\SI{6.3}{\percent}, a \textbf{factor of $6.1$}; at the upper end of the
reported ranges, $(\phi_3,\phi_4) = (0.70,0.45)$, it falls to \SI{4.1}{\percent},
a factor of $9.4$. At the crest the same two targets give factors of $5.4$ and
$8.1$, and at the benchmark $3.2$ and $4.4$
(\texttt{validation/molecular\_channel/}).

Two features of that calculation are worth stating, because neither is visible
in the dilution formula. The two fractions cannot be chosen independently: a
source large enough to make $\phi_3 = 0.60$ already drives more than $0.15$ of
$n=4$ through collisional $3\to4$ transfer, so $(\phi_3,\phi_4) = (0.60,0.15)$
requires a negative $n=4$ source and is not realisable on this operator. And
the reduction is not monotone in $\phi_3$. The difference
$\ffrac{3}(1-\phi_3) - \ffrac{4}(1-\phi_4)$ passes through zero near
$\phi_3 = 0.74$ at $\phi_4 = 0.30$, beyond which $n=4$ is the more ground-fed
of the two and the error grows again with $\overline{S}$ carrying the opposite
sign. At the crest, where that zero is reachable with non-negative sources, a
two-source dilution at $\phi_4 = 0.30$ brings the plateau error down to
\SI{0.35}{\percent} at $\phi_3 = 0.75$, past the zero at $\phi_3 = 0.7355$; a
single $n=3$ source there gives \SI{0.30}{\percent} at $\phi_3 = 0.70$ and
\SI{1.1}{\percent} again at $\phi_3 = 0.80$, with $\overline{S}$ negative. A
molecular channel strong enough cancels the diagnostic's sensitivity to the
reservoir rather than merely reducing it, and past the cancellation the error
returns.

The bound is crude, it inherits every uncertainty in the quoted fractions, and
it applies only where those fractions do.

The \emph{structural} results, the list of Section~\ref{sec:transport_partition},
are theorems about whatever matrix $\Lmat$ is supplied and hold for one
containing molecular states; what such a matrix changes is the numerical value
of the ingredients, not the form of the result.

What is \emph{not} available is a calculation with molecules in it: the factors
above are the model's own operator driven by an imposed source of the published
strength, not a solution of a matrix containing molecular states. This is an
atomic model, its cold-edge numbers are atomic-model numbers, and the bound
above is an estimate of how wrong they are rather than a correction to them.

%-----------------------------------------------------------------------------
\section{Where the ridge sits depends on what the model assumed about neutrals}
\label{sec:ridge_conditional}
% QUESTION: is the density at which the error peaks a property of hydrogen, or
% a property of this model's own ionisation balance?

\subsection{The question}

Section~\ref{sec:mechanism} found an interior maximum, a ridge, in the
error map at $\nel \approx 2\times10^{13}\,\si{\per\cubic\centi\metre}$, and
showed that it moves with the shell pair as Griem's criterion for local
thermodynamic equilibrium (LTE) predicts. The mechanism is robust; three
independent tests agreed. But a mechanism and a location are different claims,
and only the first of them has been established.

\subsection{The picture: the ridge is where the two feeds are balanced}

The ridge exists because the $n=3$ and $n=4$ shells switch from
recombination-fed to ground-fed at different reservoir strengths, so their
difference vanishes when either feed dominates and peaks in between. The
controlling variable is therefore not the electron density directly; it is
\emph{how much neutral gas there is to feed from}. The electron density enters
only because, in this model, it determines the neutral density through the CR
ionisation balance.

That is the conditionality. Change what the model assumes about neutrals and the
ridge moves, whatever the electron density does.

\subsection{The measurement}

Within the two-channel split the ground-fed fraction is exactly
$\ffrac{p} = a_1 n_g/(a_0 + a_1 n_g)$, where $n_g$ is the ground-state density
feeding the excitation channel and $a_0$, $a_1$ are the network responses of
Section~\ref{sec:R_depends_on_b}. Scaling $n_g$ is therefore an exact operation
on the ridge, not an approximation to one. Doing it, at
$\Te = \SI{2.02}{\electronvolt}$:

\begin{center}
\begin{tabular}{lr}
  \toprule
  $\ngs$ scaling & density at which $|\ffrac{3}-\ffrac{4}|$ peaks \\
  \midrule
  $\times 0.01$ & $\le 10^{12}\,\si{\per\cubic\centi\metre}$ (off grid) \\
  $\times 0.1$  & $2.28\times10^{12}\,\si{\per\cubic\centi\metre}$ \\
  $\times 1$ (the model's own CRE value) & $1.68\times10^{13}\,\si{\per\cubic\centi\metre}$ \\
  $\times 10$   & $2.48\times10^{14}\,\si{\per\cubic\centi\metre}$ \\
  $\times 100$  & $\ge 10^{15}\,\si{\per\cubic\centi\metre}$ (off grid) \\
  \bottomrule
\end{tabular}
\end{center}

\noindent \textbf{One decade in the assumed neutral density moves the ridge by
roughly one decade in electron density.} The response is close to
proportional, and the whole three-decade span of the density grid is traversed
by a two-decade change in $\ngs$.

The agreement with Griem's criterion at unit scaling, $1.93\times10^{13}$
measured against $2.05\times10^{13}\,\si{\per\cubic\centi\metre}$ predicted, is
therefore an agreement evaluated at \emph{this model's} collisional--radiative
ionisation balance, which contains no transport and no recycling. In a real
divertor the neutral density is set by recycling, which is
Section~\ref{sec:transport}'s objection arriving in a different place. The
sentence ``the ridge sits at
$\nel \approx 2\times10^{13}\,\si{\per\cubic\centi\metre}$'' is a statement about
an ionisation balance, not about a divertor, and should be quoted with that
condition attached or not quoted.

The ridge is also one line pair's and not the Balmer series': the $(3,5)$ pair
of $H_\alpha/H_\gamma$ has its worst density at
$7.2\times10^{12}\,\si{\per\cubic\centi\metre}$, a factor $2.68$ lower, and the
form of words this imposes is stated in Section~\ref{sec:error_maps}.

%-----------------------------------------------------------------------------
\section{The model does not identify the crest with detachment}
\label{sec:not_detachment}
% QUESTION: can the density at which the ridge sits be identified with a named
% operating regime of a tokamak divertor?

\subsection{The question, and the temptation}

The identification is tempting. The ridge sits where ionising and recombining
supplies compete; detachment, the regime ITER must operate in to survive its
exhaust power and the transition this diagnostic is chiefly used to identify
(Section~\ref{sec:molecules}), is where ionising and recombining plasma compete.
The words match.

\subsection{The arithmetic}

The numbers do not. Section~\ref{sec:withdrawn} withdrew the identification:
the ridge, at $1.93\times10^{19}\,\si{\per\cubic\metre}$, sits a factor $5.2$
below the detached-divertor band of published edge
modelling~\cite{Guillemaut2011} and a factor $1.6$ below the modelled upstream
separatrix density~\cite{Stangeby2023,Pitts2019}, in a range neither source
characterises.

The comparison band deserves a word, because a factor of $1.6$ is not a large
margin to be arguing over. Read from the primary source, the band is
\SIrange{3e19}{6e19}{\per\cubic\metre} across the $53$ ITER SOLPS-4.3 cases
\cite[Sec.~3.1.1]{Pitts2019}, and it is the density \emph{at the outer
midplane}, which that paper states explicitly where it plots the dependence
\cite[Sec.~3.2.2]{Pitts2019}. So the quantity is settled: outer midplane, not
X-point. What is not settled by a single factor is where in the band to
measure from. The ridge sits $1.6$ times below the \emph{bottom} of it and
$3.1$ times below the top, and the honest statement is the range rather than
either end. On the same flux tube the outboard-midplane and X-point densities
differ by roughly \SI{50}{\percent}, which is comparable to the width of the
band itself. The conclusion does not turn on any of this, since the ridge is
below the band on every reading, but the factor should not be quoted more
precisely than the band it is measured against.
The repository recorded \cite{Stangeby2023} as ``Nucl.\ Fusion 62 (2022)'' in
\texttt{src/validation/grid.py}; the DOI was right, the volume and year wrong.
It is Nuclear Fusion \textbf{63}(1) 016017 (2023), companion part A at 63(1)
016016 \cite{Stangeby2023a}; \cite{Lore2022} in the same block of that file
genuinely is volume 62.

\subsection{What that leaves}

The honest statement is uncomfortable and should be made anyway. \textbf{The
structure this work found is real and reproducible, and it sits at a density
this thesis cannot identify with any named region of a tokamak.} It is not
detachment; it is not the upstream separatrix; and calling it either would be a
claim about edge transport that a $0$-D atomic model cannot support. The ridge
is a property of how two shells of hydrogen divide their supply. It needs no
divertor geometry to produce it, and it acquires none by being described in one.

The ridge and the detached band do not overlap: restricted to
$\Te \ge \SI{2}{\electronvolt}$ and $\nel \ge 10^{14}\,\si{\per\cubic\centi\metre}$
the census is $0$ of $108$ with a worst estimate of \SI{7.2}{\percent}
(Table~\ref{tab:trapping_census}), so where the diagnostic is most used the
transient error is real but modest.

%-----------------------------------------------------------------------------
\section{The open boundary: an attack that was expected to be fatal}
\label{sec:open_system}
% QUESTION: does holding the ion density fixed manufacture the slow clock that
% the whole result rests on?

\subsection{The question}

Section~\ref{sec:rate_equation} was explicit that the ion density $\nion$ is
held fixed: the recombination feed $\Svec\nion$ enters as a constant source term
a feed stream into a CSTR in the reactor picture, rather than as a
matrix element, and the $43$-state system has no $44$th state for ionised atoms
to accumulate in. Atoms ionise and vanish from the accounting. The slow
eigenvalue $\lambda_0$ is therefore the CR ionisation rate of ground-state
hydrogen against a prescribed reservoir, not the equilibration rate of a closed
system.

\subsection{The falsifier, and what happened}

If the long $\tauslow$ values were an artifact of the open boundary, so would be
the claim that the error outlasts an ELM. The falsifier named in
Section~\ref{sec:attacks} before the test was run is a closure factor of $2300$
at $[0,4]$, the value of $\tauslow/\taudrive$ there. Closing the manifold with
a $44$th ion state gives a closure factor with minimum $1.00$, median $1.00$
and maximum $41.4$ over the grid, moves the ELM census from $202$ to $200$ of
$680$ and the worst case by one part in $325$, and leaves the count above
\SI{2}{\electronvolt} at $45$ before and $45$ after (Table~\ref{tab:closure}).
The factor is large only where $\tauslow$ already exceeds the event duration by
three orders of magnitude and is unity wherever $\tauslow$ approaches
$\taudrive$, so where the boundary condition matters the outcome does not
depend on it. The measured factor at $[0,4]$ is $15$ against the $2300$
required, a cold-corner margin; inside the defended range, at the ridge, it is
$1.01$ against $20$, and at the benchmark $\tauslow$ moves by one part in
$10^{3}$. The attack was survived everywhere it was pressed.

\subsection{The caveat that must not be softened}

The closure adds an ion reservoir. It does \emph{not} close the electron
balance, since $\nel$ is held fixed while the ion population moves, and it adds
no transport whatever. It is a test of the boundary condition on the atomic
manifold, and nothing more. Section~\ref{sec:transport} is the objection this
test does not answer, and no amount of closing the atomic system will answer it.

The two index labels for this point, $[1,4]$ in Table~\ref{tab:closure} and
$[0,4]$ in the census, are the post- and pre-step operators of one heating
step: the census rows carry the pre-step index, every quoted timescale belongs
to the post-step operator that governs the relaxation (Chapter~\ref{ch:results}),
and the two must not be interchanged. Diagonalising each gives
$\tauslow = \SI{0.421012}{\second}$ at $\Lmat[0,4]$ and
\SI{0.233239}{\second} at $\Lmat[1,4]$, so the quoted \SI{0.23324}{\second}
belongs to $\Lmat[1,4]$.

%-----------------------------------------------------------------------------
\section{An external benchmark this model fails, and what the failure measures}
\label{sec:r1_deficit}
% QUESTION: where does the comparison against published atomic data break down,
% and does the failure touch the mechanism?

\subsection{The question}

Chapter~\ref{ch:validation} presented the comparison against Fujimoto's
published population coefficients as the strongest external check this model
passes. It is, for one of the two coefficients. The other disagrees at low
principal quantum number, and this section reports what that disagreement
measures. The short answer is that it is unexplained, that it is not the
$\ell$~closure, and that it is direct evidence about the coefficients
Chapter~\ref{ch:results} uses.

\subsection{The picture: two coefficients, one for each feed}

Chapter~\ref{ch:theory} split every excited population into two pieces
corresponding to the two supply channels. The population coefficient $r_0(p)$
measures how strongly level $p$ is fed from the ion continuum by recombination;
$r_1(p)$ measures how strongly it is fed from the ground state by excitation.
The published tabulation gives both, computed independently, decades ago, by a
different method. Agreement is a real test; disagreement is a real failure.

\subsection{The disagreement, and what it measures}

The recombination-fed coefficients agree. At
$\nel = 10^{12}\,\si{\per\cubic\centi\metre}$, $r_0$ matches to
\SI{1.3}{\percent} at $p=2$ and to \SIrange{12}{14}{\percent} at $p=3$ and
$p=4$. The ground-fed coefficients at low principal quantum number do not: this
model's $r_1$ sits a factor $8.3$ low at $p=3$ and a factor $4.4$ low at $p=4$.

Those are the two shells whose difference is the mechanism of
Chapter~\ref{ch:results}, so the disagreement had to be understood rather than
noted.

Two candidate explanations were eliminated before the source was consulted. If
the excitation rates feeding the ground-fed channel were wrong, the deficit
would follow, but those rates were benchmarked against
R-matrix-with-pseudostates calculations and the dominant $\Delta\ell = 1$
channels agree at a median of $-\SI{5.0}{\percent}$ over $120$ comparisons. And
a normalisation error is excluded by $r_0$ itself: $r_0$ and $r_1$ carry the
same $p$-dependent factor $Z(p)$, so an error large enough to move $r_1$ by a
factor of $8$ would have moved $r_0$ with it.

\subsection{The source's closure, imposed and tested}

Table~4.1(b) is indexed by $p = 2, 3, 4, 5, 7, 10, 15$. Those are principal
quantum numbers, and the table carries no $\ell$ label anywhere. The published
coefficients are shell quantities computed with the $\ell$~substructure
bundled.

This model resolves $\ell$ below $n=9$, so the closures differ. An earlier
draft of this section treated that difference as the explanation, on the
evidence that recomputing the model with proton $\ell$-mixing removed moves
$r_1(2)$ by a factor $118$ and $r_1(3)$ by a factor $3.35$ at
$\nel = 10^{12}\,\si{\per\cubic\centi\metre}$. That argument does not survive
its own test. Removing $\ell$-mixing is the limit in which the process is
absent; the tabulation assumes the limit in which it is infinitely fast, and
those are opposite ends rather than a bracket. Imposing the source's own closure
directly, by bundling this model's operator under statistical $\ell$~weights,
moves $r_1(3)$ by \SI{0.42}{\percent} and makes the $p=2$ agreement worse
(\texttt{verify\_fujimoto\_bundle.py}, stamped in
\texttt{validation/fujimoto\_bundle/}).

The disagreement therefore does \emph{not} measure the $\ell$~closure, and this
thesis does not explain it. An earlier draft of this chapter said the deficit
did not reach $\ffrac{3}-\ffrac{4}$ because that difference never passes
through a bundled quantity; that was wrong, and bundling is not the protection.
Section~\ref{sec:fujimoto} makes $a_p = r_1(p)\,\Zsb{p}/\Zsb{1}$, so a deficit
in $r_1(3)$ and $r_1(4)$ is a deficit in $a_3$ and $a_4$ themselves, and the
only protection is that $\Delta$ and the cap depend on the ratio $a_3/a_4$, in
which a deficit common to both shells cancels. It is not common at
\SI{e12}{\per\cubic\centi\metre}, $8.3$ against $4.4$, and nearly so at
\SI{e15}{\per\cubic\centi\metre}, $1.8$ against $1.9$. Dividing $a_3$ and
$a_4$ by those ratios at every grid point
(\texttt{validation/fujimoto\_r1\_consequence/}) moves $\Delta$ by $+0.64$ and
$-0.08$; over the $280$ points above \SI{2}{\electronvolt} the cap moves by
\SI{+29}{\percent} and \SI{-4}{\percent} at the median, and the sensitivity at
the operating point by \SI{+30}{\percent} and \SI{+11}{\percent} at the median,
with ranges of $-58$ to \SI{+524}{\percent} and $-38$ to \SI{+68}{\percent}.
The disagreement is direct evidence about the coefficients
Chapter~\ref{ch:results} uses; what is open is whether the deficit lies in this
model or in the tabulation, and the stake is stated in
Section~\ref{sec:fujimoto}. The open status is carried in
Section~\ref{sec:scope_open}.

\subsection{Two limitations of the benchmark, stated precisely}

Table~4.1(b) is computed at $T_\mathrm{e} = \SI{1.28e5}{\kelvin}$, which is
\SI{11.03}{\electronvolt}. The grid used here reaches \SI{10}{\electronvolt},
so the comparison is made at the grid edge and \SI{9.3}{\percent} below the
tabulated temperature.

There is no second usable case. The companion Table~4.1(a) is at
$T_\mathrm{e} = \SI{e3}{\kelvin}$, which is \SI{0.0862}{\electronvolt}, two
orders of magnitude below this grid. The external check is therefore one
temperature at one edge, and it should not be described as more than that.

\subsection{What the benchmark establishes, and what it does not}

What the benchmark does establish is the recombination-fed channel, which
agrees, and that is worth stating as a pass rather than leaving implicit beside
a disagreement. What it does not establish is the absolute normalisation of the
ground-fed channel against an independent calculation, because no tabulation
resolved in $\ell$ was available to compare against.

The sensitivity of the density maximum to the ground-fed supply remains a
separate and open question, and Section~\ref{sec:ridge_conditional} treats it
on its own terms: the location responds roughly decade-for-decade to the assumed
neutral density, which is a stronger reason to quote it as a range than
anything the benchmark shows.

%-----------------------------------------------------------------------------
\section{The top of the level ladder, and what the bundling costs}
\label{sec:nmax}
% QUESTION: the model stops at n=15 and bundles the levels above n=8 -- what
% rests on the part that is approximated?

\subsection{The question}

A hydrogen atom has infinitely many bound levels, crowding together towards the
ionisation limit. A computation cannot have infinitely many states, so a
boundary was chosen: this model tracks levels up to $n_{\max} = 15$ and
discards the rest. Above $n = 8$ it also stops resolving orbital angular
momentum $\ell$, treating each shell as a single bundled state populated
statistically. Both are choices, and this section states what they cost.

\subsection{Three consequences, in decreasing order of severity}

\paragraph{The terminal shell is over-populated.} Truncation removes from the
terminal shell's balance the collisional excitation out of it to $n > n_{\max}$,
a way out, and the de-excitation and cascade into it from those levels; an
earlier draft said the shell had ``nothing above it to cascade out to'', which
names the wrong process, since cascade is downward. Against published values
this model's ground-fed coefficient at $n = 15$ is high by a factor of $4.9$ to
$6.2$, while at $n = 10$ it agrees to \SIrange{4}{20}{\percent}; the budget is
in Section~\ref{sec:convergence}. \emph{No result in this thesis rests on the
top shell}: the observable is built from $n = 3$ and $n = 4$, both fully
$\ell$-resolved and far from the truncation.

\paragraph{The bundling nonetheless enters the result indirectly.} The
recombination flux cascading down from the bundled shells feeds the
recombination-fed channel of both $n=3$ and $n=4$, one of the two channels the
mechanism compares, which is why Section~\ref{sec:ridge_alternatives} lists
truncation among the alternative explanations for the ridge that this work does
not exclude.

\paragraph{There is no $\ell$-mixing inside the bundled block.} The
proton-impact $\ell$-mixing that Section~\ref{sec:lmixing} showed to be
essential for the $n = 2$ manifold is present only for $n = 2$--$8$; the
$\ell$-mixing rate array is identically zero above that. Within the bundled
shells the $\ell$-distribution is assumed statistical rather than computed. At
high $n$ that assumption is far more defensible than it would be at low $n$,
because the sublevels are nearly degenerate and every collisional process mixes
them, but it is assumed, not measured, and this thesis has not measured it.

\subsection{The check, run}
\label{sec:bundling_check}

The script that quantifies this was written and, until now, never executed.
Running it produced a result and also produced three defects in the script
itself, which are reported here because two of them would have made its verdict
worthless.

\paragraph{What the script could not do.} It promised to test the bundled
shells against the pipeline's own $\ell$-mixing array, and that is impossible:
the array is identically zero for states $36$ to $42$, in both directions,
because $\ell$-mixing was computed for the resolved block alone. The file
contains nothing to test. The comparison has to be made against the rate
\emph{function} rather than the rate table, and the function is the one
\texttt{compute\_lmix.py} used to build the resolved block, so it is the model's
own rate rather than a second one. Calling it at $n = 2$ to $8$ reproduces the
stored array to $1.0000$, which is a wiring check and not a physical
confirmation; it establishes only that the rate carried up to the bundled
shells is the same function with the same arguments.

\paragraph{The two defects that would have inverted the answer.} The script
carried its own Pengelly--Seaton expression, independent of the pipeline's,
which disagrees with the rate the model actually uses by a factor of $95$ at
$n = 2$ rising to $5.6\times10^{3}$ at $n = 8$, growing roughly as $n^{2}$. Used
as written, it would have reported the mixing rate two to three orders of
magnitude below its true value and could have returned a verdict of marginal
where the answer is not close. It is deleted rather than corrected: a
verification script has no business maintaining a second implementation of a
rate the pipeline already computes. The script also fell back on a synthesised
logarithmic temperature grid of its own when the pipeline's grid files were
absent,
printing a warning and continuing, which is the failure mode described in
Chapter~\ref{ch:validation}: every ratio would have carried a correct-looking
label attached to conditions the model was never evaluated at. It now raises.

\paragraph{The result.} The criterion is that the $\ell$-mixing rate out of the
$np$ substate, the substate the Lyman channel empties fastest, must exceed that
substate's total radiative rate. Over the seven bundled shells and all $400$
grid points, the ratio is above the well-mixed threshold of $10$ at every one of
the $2800$ combinations. The worst case is $n = 9$ at
$\Te = \SI{1}{\electronvolt}$, $\nel = \SI{e12}{\per\cubic\centi\metre}$, where
the mixing rate is $\SI{1.155e10}{\per\second}$ against
$A(9p) = \SI{7.51e6}{\per\second}$, a margin of $1538$, that is, a factor $154$
clear of the threshold (\texttt{verify\_bundling\_psm20.py},
\texttt{outputs/bundling/bundling\_report.txt}).

Two inputs to that comparison were checked rather than assumed. The Einstein
coefficients $A(np\to1s)$ for $n = 9$ to $15$ are hardcoded in the script,
because the model's own radiative file resolves $\ell$ only to $n = 8$ and has
nothing to read above it. Since $A(np\to1s)\,n^{3}$ tends to a constant for
hydrogen, the model's resolved values and the hardcoded table must lie on one
curve: they do, with a step of \SI{0.43}{\percent} across the boundary at
$n = 8$ to $9$ and a monotone approach to $4.13\times10^{9}$ at $n = 15$. And
the $\ell$-mixing coefficient carries a Debye cutoff frozen at
$\nel = \SI{e14}{\per\cubic\centi\metre}$, because the array has no density
axis. That freezing is not small: re-evaluating the cutoff at the local density
changes the $n = 15$ rate by a factor $20$ upward at
\SI{e12}{\per\cubic\centi\metre} and a factor $30$ downward at
\SI{e15}{\per\cubic\centi\metre}. Against the radiative criterion it does not
move the verdict, because the worst point sits at the lowest density, where the
local cutoff raises the rate: the worst ratio is $1538$ with the freezing and
$3.26\times10^{4}$ without it. At the highest densities the same correction
lowers the rate, and that matters for the criterion that follows.

\paragraph{What this closes and what it does not.} The criterion above compares
mixing with radiative decay, and radiative decay is not how a high shell
empties at these densities. The total loss out of a bundled state, $-L_{kk}$,
is \SIrange{89.5}{99.1}{\percent} collisional transfer to other shells,
\SIrange{0.9}{10.5}{\percent} ionisation and at most \SI{0.19}{\percent}
radiative (\texttt{verify\_bundling\_total\_loss.py},
\texttt{validation/bundling\_total\_loss/}). Against that total the margin is
not three orders of magnitude. With the frozen cutoff the model is built with,
mixing exceeds total escape at every one of the $2800$ combinations, by $26$
at $n=9$ falling to $2.85$ at $n=14$, the minimum, all at
\SI{1}{\electronvolt} and \SI{e12}{\per\cubic\centi\metre}, and $288$
combinations sit below $10$. With the cutoff evaluated at the local density
the ratio falls below unity at $121$ of the $2800$, in every bundled shell, at
the two highest densities and $\Te \le \SI{2.95}{\electronvolt}$; at
\SI{1}{\electronvolt} and \SI{e15}{\per\cubic\centi\metre} it runs from $0.86$
at $n=9$ to $0.09$ at $n=14$. The comment in \texttt{compute\_lmix.py} that the
cutoff factor varies by about \SI{30}{\percent} across the grid understates the
variation already at $n=2$, where the local-to-frozen ratio of $F$ runs from
$1.96$ at \SI{e12}{\per\cubic\centi\metre} to $0.53$ at
\SI{e15}{\per\cubic\centi\metre} at \SI{1}{\electronvolt}
(\texttt{validation/bundling\_total\_loss/}, RESULT~4), and fails outright
above $n=9$, where $U_m > 1$ and $F$ falls as $\nel^{-3/2}$; it is reported
here and left in place. What survives is narrower than an
earlier draft claimed. Where mixing wins, the $\ell$-distribution is
statistical whatever the feeds do. Where collisional transfer wins, at high
density, the distribution is set by the feeds, and $n$-changing collisions
between adjacent shells, which carry \SIrange{71}{95}{\percent} of the
collisional loss, are close to $\ell$-blind at high $n$, so the statistical
assumption is plausible there for a different reason, one this thesis has not
tested. It does not close the terminal-shell over-population above, which is a
different defect with a different cause, and it does not make the bundled
block's $\ell$-distribution a computed quantity. It remains assumed, now with
two numbers attached and a stated region where the first of them does not
license it.

%-----------------------------------------------------------------------------
\section{Remaining assumptions, and the sensitivities that are known}
\label{sec:other_assumptions}
% QUESTION: what else was assumed, and which of those assumptions has actually
% been tested?

Four assumptions remain. Each is stated with what is known about it, rather than
listed as a disclaimer, and two of the four have been tested.

\paragraph{Maxwellian electrons: untested.} Every rate coefficient in
Chapter~\ref{ch:atoms} is an average of a cross section over a Maxwell--Boltzmann
electron energy distribution. Ionisation and excitation are driven by the
high-energy tail of that distribution, which holds a small minority of the
electrons. A \SI{1}{\electronvolt} plasma next to a recycling surface has no
guarantee of a Maxwellian tail, and a depleted or enhanced tail would change the
rates that set both clocks. Not tested here.

\paragraph{$T_i = \Te$ in the $\ell$-mixing rates: untested.} The
proton-impact $\ell$-mixing rates of Section~\ref{sec:lmixing} are evaluated at
an ion temperature equal to the electron temperature. In a detaching divertor
the ions are frequently colder than the electrons. Not tested here.

\paragraph{The step is instantaneous: tested, and safe for a reason that is
easy to get wrong.} The error maps impose a sudden temperature step, whereas a
real ELM has a finite rise time. The tempting ratio
$\taurel/\taudrive \approx 10^{-5}$ is about the wrong quantity, the fast
excited-manifold step error (Section~\ref{sec:persistence}); the plateau error
is governed by the slow ground-state response, for which the relevant ratio is
$\tauslow/\taudrive$. At the worst point of the grid that
ratio is $2332$, but that point sits at \SI{1}{\electronvolt} and the figure is
not the margin this result runs on. Over the $448$ defended pairs the ratio has
median $0.815$, and over the $45$ that break down it runs from $1.95$ to $288$,
median $11.2$ (\texttt{validation/divertor\_map/}). The step idealisation is
therefore safe by one to two orders of magnitude where a plateau is reported,
not by three. The governing quantity is in any case
$D_{\rm e} = \tauslow/t_{\rm ramp}$ rather than $\tauslow/\taudrive$: a linear
ramp reaches $D_{\rm e}(1-\mathrm{e}^{-1/D_{\rm e}})$ of the step's plateau
(\texttt{validation/ramp\_plateau/}), so any ramp with
$\taurel \ll t_{\rm ramp} \ll \tauslow$ reaches the same plateau and the
idealisation \emph{is} safe. It is safe because of the slow ratio, and the fast
ratio must never be quoted for or against this result.

\paragraph{The step is in temperature only: tested, and the correction goes the
wrong way.} A real ELM raises $\nel$ as well as $\Te$. The joint map over $4968$
cases (Section~\ref{sec:joint_step}) leaves the exact decomposition unchanged,
with the gain generalising to an additive gradient accurate to
\SI{0.2}{\percent} at the median; over $\Te\ge\SI{2}{\electronvolt}$ the
breakdown fraction moves from $10.0$ to \SI{11.3}{\percent} when the density
steps up one index and the worst case from $0.1748$ to $0.2206$. The earlier
benchmark spot check, a bound falling by a factor $4$, had the right number and
the wrong sign for the grid, because the benchmark and the worst point lie on
opposite sides of the density crest; the fixed-density map understates the
worst case.

\medskip
Two sensitivities of the reported census belong in the same place, because both
are properties of a reporting choice rather than of the physics, and both would
otherwise be mistaken for measurements.

The step size is one interval of the temperature grid, which displaces the
reservoir by $|\ln\xscale|$ between $0.124$ and $0.682$ across the grid. A
four-interval step gives $0.516$ to $2.563$ and raises $\epsplat$ at the ridge
from \SI{18.1}{\percent} to \SI{81.1}{\percent}. On step size alone, therefore,
the reported numbers \emph{understate} the error rather than overstate it.

The \SI{10}{\percent} breakdown threshold is doing real work: at \SI{5}{\percent},
\SI{10}{\percent} and \SI{20}{\percent} the count over the same $680$ points is
$348$, $202$ and $61$ (Section~\ref{sec:persistence}), so the count is not robust
to the threshold, the worst case is, and ``$N$ points break down'' must carry both.

%-----------------------------------------------------------------------------
\section{Scope: five questions this thesis leaves open}
\label{sec:scope_open}
% QUESTION: which questions are open, why can none of them be closed from
% inside a 0-D atomic model, and what would close each?

The objections computed in this chapter were all answerable from inside a
$0$-D collisional--radiative matrix. Five are not, and they are gathered here
rather than left scattered so that a reader can see at once what this thesis
does not determine. Each is stated with the reason it is out of reach and with
the calculation that would settle it. None is a defect of the implementation;
each is a boundary of the model class.

\paragraph{1. The $r_1$ deficit is unexplained.} This model's ground-fed
population coefficient sits a factor $8.3$ below Fujimoto's Table~4.1(b) at
$p=3$, the shell the mechanism runs on, and a factor $1.8$ to $2.0$ below it at
$\log_{10}(\nel/\si{\per\cubic\metre}) = 21$; truncation (\SI{1.2}{\percent}),
Lyman trapping and the source's own statistical-$\ell$ closure
(\SI{0.42}{\percent}) have each been eliminated (Section~\ref{sec:r1_deficit}).
What would settle it is an $\ell$-resolved published tabulation or the atomic
data underlying the tabulated values, neither of which was available. The
deficit is reported as open, with its consequence for $\ffrac{3}-\ffrac{4}$
computed in Section~\ref{sec:fujimoto}.

\paragraph{2. There is no convergence study above $n_{\max} = 15$.} The state
space is cut at $n = 15$ because the atomic data stop there: the quantum
cross sections extend to $n = 10$, shells $11$ to $15$ are already semiclassical
Vriens--Smeets values, and no radiative or recombination coefficients exist
above $15$ in any form. Truncation has been tested downward, and it moves
$\ffrac{3}-\ffrac{4}$ by about \SI{1}{\percent} with the sign changing across
the grid (Section~\ref{sec:convergence}). What is untested is the sensitivity of
the ridge \emph{position} to $n_{\max}$, which would require rebuilding the rate
pipeline at $n_{\max} = 20$ rather than running a validation script: new
excitation, ionisation, radiative and recombination coefficients for five
further shells, and a re-indexing of every array built on the $43$-state layout.
That is a rate-generation exercise, not an analysis one.

\paragraph{3. The ground-state population is solved for, not supplied.} Every
magnitude in Chapter~\ref{ch:results} is conditional on $\ngs$ coming from
$\Lmat\nvec + \Svec\nion = 0$; in a divertor it is set by recycling, on a
transport timescale four orders of magnitude shorter than the local ionisation
balance (Section~\ref{sec:transport}). The open-boundary test of
Section~\ref{sec:attacks} bounds the sensitivity without removing the
conditional, and what would settle it is stated in Section~\ref{sec:transport}.

\paragraph{4. There are no molecular channels in $\Lmat$.} At the onset of
detachment a large fraction of the Balmer light is molecular in origin. The
dilution estimate of Section~\ref{sec:molecules} says how far the cold-corner
number would move under one specific dilution picture, not what a molecular
model would give. What would settle it is a state space containing
$\mathrm{H}_2$, $\mathrm{H}_2^+$ and $\mathrm{H}^-$ with their rate set, a
different model rather than an extension of this one.

\paragraph{5. The electrons are assumed Maxwellian.} The reasons are given in
Section~\ref{sec:other_assumptions}. The rate integrals as built do not accept
an arbitrary electron energy distribution, so this is not a parameter that can
be varied within the present code, and the $T_i = \Te$ assumption in the
$\ell$-mixing rates is untested for the same reason.

\medskip

\noindent Items 3 and 4 are named in Chapter~\ref{ch:conclusions} among the
further work. Items 1 and 2 are bounded in their consequences rather than in
their causes: the thesis states what each would change if it went the wrong
way, and for item 1 that stake is computed in Section~\ref{sec:fujimoto}.

%-----------------------------------------------------------------------------
\section{What is left standing}
\label{sec:limits_summary}
% QUESTION: after all of that, what may this thesis actually claim?

\subsection{What this model cannot say}

It cannot say anything quantitative about the cold edge of its own grid. Below
\SI{2}{\electronvolt} the optically thin plateau error is inflated by a factor
of $2.5$ to $3.3$ by neglected Lyman trapping; the size of the correction
depends on a slab thickness the model does not contain; and below
\SI{1.15}{\electronvolt} even the location of the ridge is undetermined.

It cannot say what a real divertor does, because recycling renews the
ground-state population there on a transport timescale the model does not
contain: \SI{72}{\micro\second} against \SI{233}{\milli\second} at the worst
point, a factor $3.2\times10^{3}$, below the scope boundary; inside the defended
range one to three orders rather than four, with a median $\tauesc/\tauslow$ of
$0.0079$ over the $45$ defended pairs (Table~\ref{tab:transport_selection}); and
$3.6$ at the benchmark, where transport is the \emph{slower} of the two.

It cannot say what an instrument records at the onset of detachment, because a
large fraction of that light is molecular in origin and there are no molecules
in $\Lmat$.

It cannot name the density at which its own ridge sits as a region of a tokamak.
The ridge falls between the modelled upstream separatrix and the detached
divertor band, and belongs to neither.

And it cannot establish the absolute normalisation of its ground-fed
population coefficient against an independent calculation: $r_1$ sits a factor
$8.3$ below the standard tabulation at $p = 3$, the disagreement is
unexplained, and its stake for $\ffrac{3}-\ffrac{4}$ is computed in
Section~\ref{sec:fujimoto} (Sections~\ref{sec:r1_deficit}
and~\ref{sec:scope_open}).

\subsection{What this model can say, with the objections computed}

Above \SI{2}{\electronvolt} the breakdown census is invariant under Lyman
trapping across a twentyfold range of assumed slab thickness, $45$ of $448$ in
every run with the worst case moving by \SI{0.5}{\percent}
(Section~\ref{sec:opacity}). Closing the open boundary moves the census by two
points out of $680$ and the worst case by one part in $325$; the factor $2300$
that would have refuted the result at the cold corner measures $15$, and at the
ridge the factor required is $20$ against a measured $1.01$
(Section~\ref{sec:attacks}). The ridge mechanism was verified over
$\Te \ge \SI{2}{\electronvolt}$, where trapping is now shown to be irrelevant,
so the agreement with Griem's boundary criterion is not an artifact of neglected
opacity. At citable divertor densities above \SI{2}{\electronvolt} the
\SI{100}{\micro\second} time-averaged error never reaches \SI{10}{\percent} over
$108$ points, worst case \SI{7.2}{\percent} (Section~\ref{sec:not_detachment}).

\subsection{The bridge}

The pattern is consistent, and it is why this chapter is not an apology. Every
objection that could be computed inside a $0$-D atomic model was computed, and
the defensible claim came through all of them with its numbers unchanged. The
objections that remain, neutral transport and molecular chemistry, are the
ones that cannot be computed inside a $0$-D atomic model at all, and they have
been stated as conditionals rather than dissolved into caveats.

That is the boundary of the claim. What is inside it is a quantitative statement
about when a spectroscopist may trust a lookup table and when she may not, and
Chapter~\ref{ch:conclusions} puts it into the form she can use.
 -- those labels are reused below, so every existing
%  cross-reference from Chapters 3 and 5 continues to resolve.
%
%  Sources for every number, in order of appearance:
%    outputs/findings_10_four_agent_review.md   ss 3.1, 4.1-4.4, 7.3, 8, 10,
%                                               and the ADDENDUM (A.1-A.6)
%    validation/lyman_trapping/lyman_trapping.txt    (run 9-10 Sep 2026,
%                                               src/validation/verify_lyman_trapping.py)
%    thesis_tex/chapter5_C5D.tex:284    (the "It is not detachment" wording)
%    outputs/findings_09_central_quantity_misnamed.md   ss 8, 9
%=============================================================================

\chapter{What This Model Cannot Say}
\label{ch:limits}

Chapter~\ref{ch:results} handed the reader two things: a map of how wrong a
tabulated $H_\alpha/H_\beta$ inversion can be while the ground state is still
catching up, and a mechanism that explains the map's shape. The map has a ridge
in it, the ridge has a closed-form explanation, and the explanation was tested
three separate ways. At this point in the argument the reader believes a
result.

This chapter's job is to bound that belief without destroying it. Every number
in Chapter~\ref{ch:results} came out of one $43$-state matrix describing one
kind of atom under one set of assumptions, and some of those assumptions are
badly violated in exactly the corner of the map where the numbers are largest.
An examiner will find that corner. It is better to arrive there first, with
arithmetic.

\medskip
\noindent\textbf{The standard this chapter holds itself to.} A limitation that
has been \emph{computed} is worth more than a limitation that has been hedged.
Where an objection could be turned into a calculation, it was, and the
calculation is reported here with the number that would have refuted the
result. Two objections were expected to be fatal and were not; they are given
the same space as the ones that bite. Where an objection cannot be computed
inside a model with no spatial dimension, it is stated as a \emph{conditional},
``\emph{if} this holds, the error is as computed'', and not softened into a
caveat sentence. Where a check has simply not been run, this chapter says so.

\medskip
\noindent\textbf{Three abbreviations recur below} and are expanded here for
this chapter: \emph{CR} for collisional--radiative, the class of model built in
Chapter~\ref{ch:theory}; \emph{QSS} for quasi-steady-state, the Bodenstein-style
approximation that eliminates the fast excited levels; and \emph{CRE} for
collisional--radiative equilibrium, the fully settled steady state that the
diagnostic lookup table assumes and that this thesis measures the cost of
assuming. An \emph{ELM} (edge-localised mode) is the repetitive burst of
heat and particles that a tokamak edge throws at the divertor every few
milliseconds, and it is the transient every ``event duration'' below refers to;
$\taudrive = \SI{100}{\micro\second}$ throughout, the duration used in
Chapter~\ref{ch:results} and the aggressive end of the range; the
ITER-pedestal extrapolation is \SI{506}{\micro\second}
\cite[Sec.~5]{Loarte2003} (Section~\ref{sec:persistence}).

%-----------------------------------------------------------------------------
\section{Optical thickness: when the plasma cannot let its own light out}
\label{sec:opacity}
% QUESTION: the model lets every emitted photon escape the plasma -- where is
% that false, what does it do to the error map, and how much of the map
% survives it?

\subsection{The question, and the picture behind it}

Chapter~\ref{ch:atoms} gave every downward transition an Einstein $A$
coefficient: the probability per second that an excited atom drops to a lower
level and emits a photon. Those coefficients entered the CR matrix $\Lmat$
bare. Entering bare means one specific physical assumption: that every photon
emitted anywhere in the plasma leaves it and never comes back. A plasma for
which this is true is called \emph{optically thin}.

A chemical engineer has met the failure of this assumption in another costume.
Consider a reaction whose product is removed by desorption from a catalyst
surface. If the product leaves the reactor, the desorption step runs at its
intrinsic rate. If the product cannot leave, if it re-adsorbs somewhere else
in the bed and is put back onto a site, then the \emph{effective} desorption
rate is smaller than the intrinsic one, and the whole network downstream of that
step shifts. Radiation trapping is that, with photons.

Here is the specific chain. The strongest transition in hydrogen is
Lyman-$\alpha$: an atom in the $2p$ level drops to the ground state $1s$ and
emits an ultraviolet photon. If the gas around it contains many ground-state
atoms, that photon has a good chance of being absorbed by one of them, which
puts \emph{that} atom back into $2p$. Nothing has left the plasma; a
$2p$ population has simply been moved from one place to another. Averaged over
the volume, the $2p$ level decays more slowly than its $A$ coefficient says. The
standard way to write this down is to multiply the coefficient by an
\emph{escape factor} $\Theta_P$ between $0$ and $1$: $\Theta_P = 1$ is a
transparent plasma, $\Theta_P = 10^{-5}$ is one in which the transition is
effectively switched off. The formulation used here is Holstein's population
escape factor~\cite{Holstein1947}.

\textbf{Why this reaches the populations, and therefore the result.} If $2p$
lives longer, atoms pile up in $n=2$. Atoms in $n=2$ are only
\SI{10.2}{\electronvolt} from the ground state but only
\SI{3.4}{\electronvolt} from the continuum, so they are far easier to ionise
than ground-state atoms. A longer $n=2$ lifetime therefore means more
\emph{stepwise} ionisation, ionisation via an excited level rather than
directly out of the ground state, which means fewer neutrals and a faster
turnover of the ground-state reservoir. The ground-state reservoir is the whole
mechanism of this thesis. Trapping is not a cosmetic correction to a line
intensity; it acts on the slow clock.

\subsection{Where the assumption fails, in numbers}

The formalism needed is one line. A photon travelling a distance $D$ through a
gas of absorbers is attenuated by $e^{-\tau}$, where the \emph{optical depth}
$\tau = \sigma_0\,\ngs\,D$ is dimensionless: $\ngs$ is the ground-state density
in \si{\per\cubic\centi\metre}, $D$ the path length in \si{\centi\metre}, and
$\sigma_0$ the absorption cross section at line centre in
\si{\square\centi\metre}. For a Doppler-broadened line,
$\sigma_0 = 2.654\times10^{-2}\,f/(\sqrt{\pi}\,\Delta\nu_D)$, with $f$ the
oscillator strength (a dimensionless measure of how strongly the transition
couples to light) and $\Delta\nu_D$ the Doppler width set by the temperature of
the absorbing atoms. $\tau \ll 1$ is transparent; $\tau \gg 1$ is opaque.

The number that matters is $\ngs$, not $\nel$, and Chapter~\ref{ch:theory}
already established the awkward fact that these two run in opposite directions:
the plasma is most neutral where it is coldest. The worst point of the error map
is $\Te = \SI{1.0}{\electronvolt}$,
$\nel = 5.18\times10^{13}\,\si{\per\cubic\centi\metre}$, the maximum of
$\epsplat$ over all $680$ window-resolvable points of the map, cold edge
included. There the model's own CR solution gives $\ngs/\nion = 28.3$, so
quasi-neutrality fixes
$\ngs = 1.47\times10^{15}\,\si{\per\cubic\centi\metre}$. At
\SI{1}{\electronvolt} that gives
$\sigma_0 = 5.47\times10^{-14}\,\si{\square\centi\metre}$, and the Lyman-$\alpha$
optical depth per centimetre of path is

\begin{center}
\begin{tabular}{lr}
  \toprule
  $\Te$ (eV) & $\tau_{\rm Ly\alpha}$ per \si{\centi\metre} at
              $\nel = 5.18\times10^{13}\,\si{\per\cubic\centi\metre}$ \\
  \midrule
  1.00 & 80.3 \\
  1.15 & 10.6 \\
  1.60 & 0.224 \\
  2.02 & 0.0255 \\
  2.95 & 0.00163 \\
  \bottomrule
\end{tabular}
\end{center}

\noindent An optical depth of $80$ per centimetre means a Lyman-$\alpha$ photon
travels about a tenth of a millimetre before it is reabsorbed. Over
\SI{5}{\centi\metre} of plasma the line-centre depth
over the half slab is $201$ and the population escape factor is
$1.2\times10^{-3}$; roughly one photon in eight hundred gets out. Chapter~%
\ref{ch:theory} quotes $3\times10^{-5}$ at a different grid point, a different
slab thickness and a different density, so the two are not the same quantity and
the agreement in order of magnitude should not be read as a reproduction of
one by the other.
The blunt reading is that the effective $A_{2p\to1s}$ belonging in $\Lmat$ at
that point is wrong, self-consistently, by about $1.8$ orders of magnitude, or
$2.9$ orders on the thin one-shot estimate
(\texttt{verify\_lyman\_optical\_depth.py}).

Two rows down the same table the problem has evaporated. At
\SI{2.02}{\electronvolt} the depth per centimetre is $0.0255$, and at the
benchmark point (\benchTe, \benchne) it is $0.0039$; a photon there crosses
metres of plasma before anything happens to it. \textbf{The measured quantity
is not ``how opaque is this plasma'' but ``how fast does opacity switch on as
the plasma cools'',} and the answer is a factor of $5\times10^{4}$ between
\SI{1}{\electronvolt} and \SI{2.95}{\electronvolt}.

That is why the results of Chapter~\ref{ch:results} were declared for
$\Te \gtrsim \SI{2}{\electronvolt}$ in the first place
(Section~\ref{sec:two_references}). It is also, uncomfortably, where the
headline was living. Recounting the \SI{100}{\micro\second} breakdown census against optical depth
rather than against temperature: of the $202$ points at which the
\SI{100}{\micro\second} time-averaged estimate of the error exceeds \SI{10}{\percent}, an extremum count over
the $680$ window-resolvable points of the full grid, $157$ lie below
\SI{2}{\electronvolt}, $85$ have
$\tau_{\rm Ly\alpha}(\SI{5}{\centi\metre}) > 1$, and $15$ have it above $100$,
on the half-slab depth used above.

An earlier draft gave $105$ and $26$, a convention error rather than a
recomputation: those counts are reproduced only by combining the full-path
depth $\kappa_0 D$ with a Doppler width $\sqrt{kT_n/m}$, short of the correct
$\sqrt{2kT_n/m}$ by $\sqrt{2}$. With the correct width the census is $85$ and
$15$ on the half-slab depth quoted here and $100$ and $23$ on the full-path
depth. A deuterium cross-section with the full-path depth also returns $105$
and $26$, since the isotope mass enters the width as the same $\sqrt{2}$, so
the two errors are indistinguishable from the counts alone. All seven
conventions are tabulated in \texttt{validation/lyman\_optical\_depth/}.
Nothing in the defended range depends on the choice: every crossing lies below
\SI{2}{\electronvolt} except two, both at $D = \SI{20}{\centi\metre}$.

\subsection{What we expected trapping to do, and what it did}

The expectation, stated before the calculation was run, was that trapping would
demolish the map: an effective $A$ coefficient wrong by about three orders of
magnitude at the worst point ought to invalidate everything computed there and
nearby, and the project's own records warned that under strong trapping the
fast relaxation time $\taurel$ rises to \SI{2.05}{\micro\second} and the
timescale separation $\Msep$ falls by a factor $158$. That expectation was
wrong.

So the calculation was done. The Lyman series was given a self-consistent
Holstein escape factor and the entire error map rebuilt. Self-consistent matters
here: $\Theta_P$ depends on $\ngs$, and $\ngs$ depends on $\Theta_P$ through the
CR balance, so the two were iterated to a fixed point rather than evaluated once.
The fixed point converged at all $400$ grid points for every slab thickness
tested, and non-convergence was made to raise an error rather than silently
return the last iterate. All $14$ Lyman channels were trapped, the seven
resolved $n\mathrm{P}\to1\mathrm{S}$ transitions and the seven bundled ones,
not Lyman-$\alpha$ alone.

Three gates were passed before any result was looked at. The oscillator
strengths were derived from this repository's own $A$ coefficients by inverting
the Einstein relation rather than imported, and come out as $f = 0.4162$,
$0.0791$, $0.0290$, $0.0139$ for Lyman-$\alpha$ through Lyman-$\delta$, the
accepted literature values to four figures. The derived Lyman-$\alpha$
$\sigma_0$ agrees with an independent implementation to \SI{0.059}{\percent},
after a first run that failed at \SI{1156}{\percent} on a $4\pi$ factor; and
rebuilding the matrix with trapping switched off reproduces the canonical
$\Lmat$ exactly, so every trapped number below is a difference against the real
matrix (Section~\ref{sec:severe_checks}).

\textbf{One new quantity had to be invented to do this at all, and it is not in
the model.} The escape factor depends on how far a photon must travel to get
out, so it depends on the size of the plasma, the slab thickness $D$. A model
with no spatial dimension has no $D$. It was therefore swept, from
\SI{1}{\centi\metre} to \SI{20}{\centi\metre}, and never fixed; every trapped
number below carries the $D$ it was computed at. Reporting a single trapped
number without its $D$ would be reporting an assumption as a measurement.

\begin{table}[tbp]
  \centering
  \caption[Breakdown census under Lyman trapping]{Effect of radiation
  trapping on the \SI{100}{\micro\second} breakdown census. Each entry counts the grid points at
  which the time-averaged estimate of the CR-equilibrium-reference error,
  averaged over a chosen duration $\taudrive = \SI{100}{\micro\second}$, exceeds
  \SI{10}{\percent}, out of the points at which a
  quasi-steady-state window is resolvable; ``worst'' is the largest estimate
  within that scope. Rows differ only in the assumed slab thickness $D$ used to
  compute the self-consistent Holstein escape factor applied to all $14$ Lyman
  transitions. The first row is the optically thin matrix used throughout
  Chapter~\ref{ch:results}. Above \SI{2}{\electronvolt} the count does not move
  at all across a twentyfold range of $D$. The $D = \SI{20}{\centi\metre}$ row
  has a slightly smaller denominator because two points lose their resolvable
  window under trapping. Source:
  \texttt{validation/lyman\_trapping/lyman\_trapping.txt}.}
  \label{tab:trapping_census}
  \footnotesize
  \begin{tabular}{lccc}
    \toprule
    matrix & all window-resolvable & $\Te \ge \SI{2}{\electronvolt}$ &
          $\Te \ge \SI{2}{\electronvolt}$, $\nel \ge 10^{14}$ \\
    \midrule
    optically thin (canonical)          & $202/680$, worst $0.3868$ & $45/448$, worst $0.1748$ & $0/108$, worst $0.0717$ \\
    trapped, $D = \SI{1}{\centi\metre}$  & $176/680$, worst $0.2925$ & $45/448$, worst $0.1748$ & $0/108$, worst $0.0712$ \\
    trapped, $D = \SI{5}{\centi\metre}$  & $170/680$, worst $0.2550$ & $45/448$, worst $0.1746$ & $0/108$, worst $0.0694$ \\
    trapped, $D = \SI{20}{\centi\metre}$ & $160/678$, worst $0.2237$ & $45/446$, worst $0.1740$ & $0/106$, worst $0.0627$ \\
    \bottomrule
  \end{tabular}
\end{table}

\subsection{Above two electron-volts, nothing moves}

Table~\ref{tab:trapping_census} contains the strongest robustness statement in
this thesis, and it is the opposite of what was expected.

\textbf{Above \SI{2}{\electronvolt} the breakdown count does not change at all,
$45$ of $448$ points in every run, and the worst case within that set
moves by \SI{0.5}{\percent}, from $0.1748$ to $0.1740$, across a twentyfold
range in assumed slab thickness.} The escape factor at the benchmark point
(\benchTe, \benchne) runs from $0.9986$ at $D = \SI{1}{\centi\metre}$ to
$0.973$ at $D = \SI{20}{\centi\metre}$. Even at the thickest slab tested,
\SI{97}{\percent} of Lyman-$\alpha$ photons still leave: the plasma is optically
thin there in the sense that actually matters, and assuming so costs less than
the resolution of the map.

This converts a disclaimer into a measurement. The
$\Te \gtrsim \SI{2}{\electronvolt}$ restriction was chosen on an escape-factor
argument before any of this was recomputed, and it sits almost exactly at the
temperature above which trapping stops mattering: a guess that has now been
checked, and survived. The defensible claim of Chapter~\ref{ch:results}, $45$
of $448$ points above \SI{2}{\electronvolt} at \SI{100}{\micro\second} and $24$
of $448$ at \SI{506}{\micro\second} (Section~\ref{sec:persistence}), is
invariant under the one physical process most likely to have invalidated it.

\subsection{Below two electron-volts, the cold-edge number is inflated
threefold}

The same calculation is much less kind to the cold edge, and it should be,
because that is where the optical depth was $80$ per centimetre.
Table~\ref{tab:trapping_worstpoint} follows the worst point of the optically
thin map through the sweep.

\begin{table}[tbp]
  \centering
  \caption[Lyman trapping at the cold-edge worst point]{The cold-edge worst
  point of the optically thin error map, at $\Te = \SI{1.0}{\electronvolt}$ and
  $\nel = 5.18\times10^{13}\,\si{\per\cubic\centi\metre}$ under a heating step,
  is the maximum of $\epsplat$ over all $680$ window-resolvable grid points; it
  is traced here as a function of the assumed slab thickness $D$. $\epsplat$ is the plateau
  error, the fractional error a lookup table makes while the ground state is
  stale; the ELM bound is its average over an event of
  $\taudrive = \SI{100}{\micro\second}$. The plateau error falls by a factor
  $2.5$ to $3.3$, but the diagnostic's intrinsic sensitivity
  $|\ffrac{3}-\ffrac{4}|$ barely moves: what collapses is how far the ground
  state is displaced, not how strongly the line ratio responds to it. Source:
  \texttt{validation/lyman\_trapping/lyman\_trapping.txt}.}
  \label{tab:trapping_worstpoint}
  \small
  \begin{tabular}{lrrrrrrr}
    \toprule
    $D$ (cm) & $\epsplat$ & ELM bound & $\tauslow$ (s) & $\Msep$ &
    $\ffrac{3}$ & $\ffrac{4}$ & $|\ffrac{3}-\ffrac{4}|$ \\
    \midrule
    $0$ (thin) & $0.38690$ & $0.38682$ & $2.3324\times10^{-1}$ & $3.74\times10^{7}$ & $0.7728$ & $0.2862$ & $0.4866$ \\
    $1$        & $0.15748$ & $0.15739$ & $9.0651\times10^{-2}$ & $2.75\times10^{6}$ & $0.6643$ & $0.2055$ & $0.4589$ \\
    $5$        & $0.13274$ & $0.13260$ & $4.6940\times10^{-2}$ & $6.13\times10^{5}$ & $0.5967$ & $0.1793$ & $0.4174$ \\
    $20$       & $0.11570$ & $0.11550$ & $2.8500\times10^{-2}$ & $2.14\times10^{5}$ & $0.5139$ & $0.1501$ & $0.3639$ \\
    \bottomrule
  \end{tabular}
\end{table}

\textbf{The \SI{38.7}{\percent} becomes \SIrange{11.6}{15.7}{\percent}}, a
factor of $2.5$ to $3.3$, and which value inside that range it takes is set
by a slab thickness the model does not contain. The point still exceeds the
\SI{10}{\percent} threshold at every thickness tested, so the qualitative
statement that the cold corner breaks down survives intact; the \emph{number}
does not, and it must never appear without its trapped range and without the
optical depth that produced it.

The route by which the number falls is itself a result, and it is not the
obvious route. The obvious guess is that trapping desensitises the diagnostic,
that a trapped plasma simply responds less. It does not. The sensitivity
$|\ffrac{3}-\ffrac{4}|$, which is the factor in the exact decomposition of
Section~\ref{sec:mechanism} measuring how strongly the line ratio answers to its
reservoir, moves by only \SI{6}{\percent} at $D = \SI{1}{\centi\metre}$ and
\SI{25}{\percent} at $D = \SI{20}{\centi\metre}$. What collapses is the other
factor: the displacement.

The chain is the one set out at the start of this section, now with numbers on
it. Trapping lengthens the $n=2$ lifetime; the longer lifetime raises stepwise
ionisation; the ground-state population falls; and the slow clock $\tauslow$
shortens by a factor $8$ at the worst point. At $\Te = \SI{1.0}{\electronvolt}$,
$\nel = 1.93\times10^{13}\,\si{\per\cubic\centi\metre}$, $\ngs$ drops from
$2.99\times10^{14}$ to $7.26\times10^{13}\,\si{\per\cubic\centi\metre}$ between
$D = 1$ and \SI{20}{\centi\metre}. A reservoir that turns over faster is a less
stale reservoir, so $|\Delta\ln\xscale|$ roughly halves. This is exactly the
two-axis attribution of Section~\ref{sec:temperature_mechanism} arriving from a
third direction: trapping acts on the displacement, not on the sensitivity.

\subsection{Below 1.15 electron-volts the ridge location is not determined}

One more thing moves at the cold edge, and it is a location rather than a
magnitude. Taking the density at which $\epsplat$ is maximal, row by row: at
$\Te \ge \SI{1.6}{\electronvolt}$ that maximum falls on
$\nel = 1.93\times10^{13}\,\si{\per\cubic\centi\metre}$ in every run, optically
thin and trapped alike, and the values themselves barely register the change:
at the benchmark temperature $\epsplat = 0.12166$, $0.12166$, $0.12165$,
$0.12163$ for $D = 0, 1, 5$ and \SI{20}{\centi\metre}, with
$|\ffrac{3}-\ffrac{4}|$ changing in the fifth decimal place. Below
\SI{1.15}{\electronvolt} the maximum wanders across three density columns, a
factor of $7$, as a function of an assumed parameter. \textbf{The two coldest
rows of the grid cannot support a claim about where the ridge is at all.}

That result cuts in the mechanism's favour, which is worth saying explicitly
because it is easy to miss. The ridge attribution of
Section~\ref{sec:mechanism} was verified over $\Te \ge \SI{2}{\electronvolt}$,
and $\Te \ge \SI{2}{\electronvolt}$ is precisely the range in which trapping has
now been shown to be irrelevant. The agreement with Griem's boundary criterion
is therefore not an artifact of neglected opacity; it was established where
opacity does not act.

\subsection{A contradiction resolved, and what is still open}

Two statements about trapping had been circulating in this project's records:
that $\taurel$ is insensitive to Lyman-$\alpha$ trapping to within
\SI{1}{\percent}, and that under strong trapping $\taurel$ rises to
\SI{2.05}{\micro\second} with $\Msep$ falling $158$-fold. They look
contradictory. Measured directly, both are true and neither had stated its
scope: at the benchmark point (\benchTe, \benchne) $\taurel$ moves from
$2.2769\times10^{-9}$ to $2.2947\times10^{-9}\,\si{\second}$ at
$D = \SI{20}{\centi\metre}$, a change of \SI{0.78}{\percent}; at
$\Te = \SI{1.0}{\electronvolt}$,
$\nel = 1.93\times10^{13}\,\si{\per\cubic\centi\metre}$ it moves from
$8.8842\times10^{-9}$ to $2.4289\times10^{-7}\,\si{\second}$, a factor of $27.3$.
The ``within \SI{1}{\percent}'' claim is a benchmark-point claim and fails by
more than an order of magnitude at the cold corner. This is the same failure
mode as quoting any other extremum without the set it is an extremum over, and
neither sentence may travel again without its regime attached.

What this calculation does not settle should be as clear as what it does. The
geometry is assumed, not derived: a homogeneous slab, uniform emission, an
isotropic radiation field, a Doppler line profile, evaluated at the plasma
centre. A $0$-D model cannot check any of those. One assumption inside it
\emph{was} checked and is not load-bearing: the temperature of the absorbing
atoms, which sets the Doppler width. Fixing it at \SI{3}{\electronvolt}, the
value a molecule dissociating by the Franck--Condon process would give, instead
of at $\Te$, changes the census at $D = \SI{5}{\centi\metre}$ from $170/680$ and
a worst case of $0.2550$ to $173/680$ and $0.2653$. Continuum lowering is
untouched here and remains open. Opacity in the Balmer lines themselves, which
would act on the observed ratio directly rather than through the level
populations, has been computed: the largest $H_\alpha$ optical depth above
\SI{2}{\electronvolt} is $0.010$, and the worst cold cell, $[0,7]$, reaches
$0.214$ (Section~\ref{sec:balmer_equivalence},
\texttt{validation/balmer\_optical\_depth/}).

%-----------------------------------------------------------------------------
\section{Transport: the atoms do not stay where the model puts them}
\label{sec:transport}
% QUESTION: the slow clock is the ionisation time of a stationary atom -- what
% happens to the result when the atom leaves before it ionises?

\subsection{The question}

This is the objection that cannot be answered inside this model, and it is a
more serious threat to the cold-edge number than anything in
Section~\ref{sec:opacity}. It is also the simplest to state.

\subsection{The picture: a reactor with no walls}

The model of Chapter~\ref{ch:theory} is a batch reactor. A fixed parcel of gas
is given a temperature, and its composition evolves until it stops changing. The
slow eigenvalue of that system, $\tauslow$, is how long the ground-state
population takes to reach its new value, and it is long, because ionising an
atom out of the ground state is slow.

A divertor is not a batch reactor. It is closer to a continuous-flow reactor
with an enormous throughput, in which the residence time of a molecule is far
shorter than the time it needs to react. Neutral atoms are born at the target
surface when ions strike it and recombine there, which is \emph{recycling},
and they fly back into the plasma. They do not wait around to be ionised
locally; the ones that do not ionise leave, and fresh ones arrive continuously.

The comparison is not close. At the worst point of the map $\tauslow$ is
\SI{233}{\milli\second}: nearly a quarter of a second for a stationary atom to
be ionised. A deuterium atom at \SIrange{1}{3}{\electronvolt} would cross a
\SI{10}{\centi\metre} plasma ballistically in \SIrange{6}{10}{\micro\second}, but
it does not travel ballistically: resonant charge exchange with the surrounding
ions randomises its direction after about \SI{1.7}{\centi\metre} at this
density, so escape is diffusive and takes closer to \SI{72}{\micro\second}. That
correction runs in this model's favour. The threshold at which the worst point
would fall below \SI{10}{\percent} is a ground-state renewal time of about
\SI{26}{\micro\second}, so the result survives by a factor of $2.8$ rather than
failing outright. It survives only because charge exchange traps the neutrals,
which is a thin margin to rest on and is stated as one.
That is one grid point, and it is one this thesis does not defend:
$[0,4]$ lies below the \SI{2}{\electronvolt} scope boundary.
Section~\ref{sec:transport_selection} asks the same question at the points that
do carry the result, and gets a worse answer.

\begin{quote}
\textbf{In a divertor, the ground-state neutral density at a given point is not
set by local ionisation balance. It is set by recycling.}
\end{quote}

\subsection{Why this bears on the result specifically}

The error of Chapter~\ref{ch:results} exists for one reason: after a temperature
step, the ground-state reservoir is \emph{stale}. Its population has not moved
while the temperature has, the excited levels equilibrate against a reservoir
that belongs to the old conditions, and the lookup table, which assumes the
reservoir finished settling, reads the wrong answer. Everything hangs on the
reservoir not being refreshed.

Neutral transport refreshes it. Any process that replaces the ground-state
population on a timescale shorter than $\tauslow$ resets the staleness, and this
model contains no such process; the \SI{26}{\micro\second} renewal threshold
quoted above is the map's own measure of how fast that process would have to be.

\subsection{The assumption fails hardest where the result is generated}
\label{sec:transport_selection}

The worst point was the wrong place to test transport: $[0,4]$ lies below the
scope boundary of Section~\ref{sec:scope}. The pairs that carry the defended
claim are the $45$ above \SI{2}{\electronvolt} whose \SI{100}{\micro\second}
estimate exceeds \SI{10}{\percent} (Table~\ref{tab:elm_census}), and there the
answer is different.

Comparing the neutral escape time with $\tauslow$ pair by pair over the grid
gives Table~\ref{tab:transport_selection}. The escape model is the one used
above, stated as a formula: the atom moves at $v=\sqrt{2k\Tn/m_{\mathrm D}}$
with $\Tn=\Te$, resonant charge exchange with the ambient ions randomises its
direction after $\lambda = v/(\nel\langle\sigma v\rangle_{\mathrm{CX}})$ with
$\langle\sigma v\rangle_{\mathrm{CX}} = \SI{1.1e-8}{\cubic\centi\metre\per\second}$,
and escape from a boundary a distance $L$ away takes $L/v$ when
$\lambda \ge L$ and $4L^{2}/\pi^{2}D$ with $D = v\lambda/3$ otherwise. At
$[0,4]$ it returns $\lambda = \SI{1.72}{\centi\metre}$ and
$\tauesc = \SI{72.3}{\micro\second}$, the values used above
(\texttt{verify\_transport\_selection.py}).

\begin{table}[htbp]
\centering
\caption{Ratio of neutral escape time to $\tauslow$, over the $45$ defended
breakdown pairs and over the other $635$ analysed pairs. $L$ is the distance
from the emitting parcel to the plasma boundary. ``Satisfies'' means
$\tauesc > \tauslow$, that is, the atom is ionised before it leaves, which is
what the closed parcel assumes.}
\label{tab:transport_selection}
\begin{tabular}{llrrr}
  \toprule
  $L$ & set & median $\tauesc/\tauslow$ & largest & satisfies \\
  \midrule
  \SI{10}{\centi\metre} & the 45 breakdown pairs & $0.0079$ & $0.159$ & $0/45$ \\
  \SI{10}{\centi\metre} & the other 635 & $0.0471$ & $384$ & $168/635$ \\
  \SI{20}{\centi\metre} & the 45 breakdown pairs & $0.0192$ & $0.636$ & $0/45$ \\
  \SI{20}{\centi\metre} & the other 635 & $0.129$ & $1535$ & $227/635$ \\
  \bottomrule
\end{tabular}
\end{table}

\textbf{Not one of the $45$ pairs that carry the defended result satisfies the
assumption used to compute it, at either scale length, while a quarter to a
third of the pairs that do not break down satisfy it comfortably.}

A count of zero is worth little if the $45$ sit just under the line, so the
distance matters. The closest of them is $[18,4]$ heating,
$\Te = \SI{2.33}{\electronvolt}$, $\nel = \SI{5.18e13}{\per\cubic\centi\metre}$,
at $\tauesc/\tauslow = 0.636$ for $L = \SI{20}{\centi\metre}$: its escape time
would have to rise by a factor $1.57$ for the parcel to close. The median of the
$45$ would have to rise by a factor $52$. One pair is marginal and the rest are
not.

Three things were varied to see whether the conclusion depends on a choice.
Tripling $\langle\sigma v\rangle_{\mathrm{CX}}$, which lengthens every diffusive
escape time in proportion, moves the count from $0$ to $3$ of $45$; dividing it
by three leaves it at $0$. Replacing $\sqrt{2k\Tn/m}$ by the Maxwellian mean
speed $\sqrt{8k\Tn/\pi m}$ shortens every escape time and leaves the count at
$0$. Doubling $L$ from $10$ to \SI{20}{\centi\metre} leaves it at $0$. The
finding is not an artifact of any of the three.

\paragraph{The selection is structural, not incidental.} A pair enters the
census because $\tauslow$ is long, since the single-slow-mode estimate of
Eq.~\eqref{eq:elm_bound} rises with $\tauslow$ at fixed $\taudrive$. A long
$\tauslow$ means slow ionisation out of the ground state. Slow ionisation is
exactly the condition under which neutral transport, whose timescale does not
depend on $\tauslow$ at all, sets the ground-state density instead of the local
atomic physics. The census therefore selects for the regime in which its own
assumption is weakest, and it does so by construction rather than by accident.
That is why testing the assumption at the cold corner gave the reassuring answer
and testing it inside the defended scope does not.

\subsection{What this reaches, and what it does not}
\label{sec:transport_partition}

The finding above is about magnitudes; it does not touch the algebra.

\textbf{Untouched by transport.} The two-channel split of
Eq.~\eqref{eq:two_supplies}, exact to a relative residual of
$3.075\times10^{-14}$; the unit-width logistic form of the ground-fed fraction;
the bound $\max|f_3-f_4| = \tanh(|\Delta|/4)$ and its corollary
$|\mathrm{d}\ln\Robs/\mathrm{d}\ln\bdep| < 1$; the sensitivity map
$\overline{S}$; and the accuracy of the excited-state closure,
$6.73\times10^{-9}$ at the worst point. These are properties of the operator
$\Lmat(\Te,\nel)$ at a fixed operating point, or theorems about its structure.
$\overline{S}$ says how the line ratio responds to whatever the reservoir does;
it is silent on what made the reservoir do it. Transport changes the reservoir's
history and cannot reach any of them.

\textbf{Conditional on the closed parcel.} The reservoir gain $G$, which is
a secant along the CRE curve and therefore assumes the ground
state responds through collisional--radiative balance alone; $\epsplat$, which
is built from it; the breakdown census; and every time-averaged magnitude in
Chapter~\ref{ch:results}. If recycling renews $\ngs$ on a timescale shorter than
$\tauslow$, $G$ is not the gain the plasma has. These numbers are not upper
estimates. A fast-exchanging reservoir with a steady source freezes the
closed-parcel error rather than removing it, and a rising source enlarges it: over
the defended set the true error exceeds the closed-parcel plateau at $165$ of the
$180$ (pair, $L$, source-multiplier) combinations tested
(\texttt{verify\_open\_reservoir.py}). The closed-parcel value is a reference
magnitude conditional on the reservoir's history, not a bound on it.

The thesis therefore leads with the structure and the bound and reports the
magnitudes as conditional; Section~\ref{sec:transport_selection} adds that the
condition is measured to fail at the points that matter, by one to two orders
of magnitude.

\subsection{The claim, restated as the conditional it is}

The honest form is therefore not a caveat but a conditional, and it should be
written this way wherever the number appears:

\begin{quote}
\emph{If} the ground-state neutral density is frozen for the duration of the
event, the error is as computed.
\end{quote}

That conditional is close to exactly true for the excited manifold, which
equilibrates in nanoseconds, and it is the standard implicit assumption of every
$0$-D collisional--radiative table in use, the object this thesis set out to
test. But it is an assumption about the plasma, not about the atoms, and no
eigenvalue analysis inside a $0$-D matrix can validate it: the two timescales
are measured to differ by four orders of magnitude and assumed not to interact.
What would settle it is a two-region or transport-coupled calculation, an edge
code such as SOLPS (Scrape-Off Layer Plasma Simulation) supplying $\ngs$ from a
recycling influx; that is outside the scope of this thesis and is named in
Chapter~\ref{ch:conclusions} among the further work.

%-----------------------------------------------------------------------------
\section{No molecular channels, in a gas that is mostly molecules}
\label{sec:molecules}
% QUESTION: the model contains only atomic hydrogen -- at the conditions where
% the error is largest, is that a defensible description of the gas?

\subsection{The question, and why a referee asks it first}

Every state in Chapter~\ref{ch:atoms}'s state space is a level of a single
hydrogen \emph{atom}, plus a reservoir of bare protons. There is no
$\mathrm{H}_2$ (or $\mathrm{D}_2$), no molecular ion $\mathrm{D}_2^+$, no
negative ion $\mathrm{D}^-$, and no reactions between them. A search of the rate
assembly for any of these returns nothing at all.

\subsection{The picture: a third feed stream}

Chapter~\ref{ch:theory} fed every excited level from two directions, up from
the ground state by electron collisions and down from the ion continuum by
recombination, and the error mechanism is the difference in how the $n=3$ and
$n=4$ shells split their supply between those two feed streams.

Molecules open a third stream. In a cold, dense, mostly neutral gas, hydrogen
atoms form $\mathrm{H}_2$; the molecules are ionised to $\mathrm{H}_2^+$, or
exchange charge, and the resulting molecular ions recombine with electrons and
break apart, leaving one of the fragments in an \emph{excited} atomic level. The
whole family of such routes is called MAR, molecular activated recombination.
The important feature for this thesis is not that MAR exists but where it
delivers: it populates $n=3$ directly, without the atom ever having passed
through the ground-state reservoir whose staleness is the entire mechanism here.
A supply channel that bypasses the reservoir does not inherit the reservoir's
lag.

\subsection{How much gas is involved}

At the worst point of the map the model's own equilibrium solution is
\SI{96.6}{\percent} neutral: $\ngs = 1.47\times10^{15}$ against
$\nel = 5.18\times10^{13}\,\si{\per\cubic\centi\metre}$. A hydrogen gas at
\SI{1}{\electronvolt} that is \SI{96.6}{\percent} neutral is a gas in which
molecules form readily. The atomic description is not being applied at the edge
of its validity there; it is being applied outside it.

The size of what is being left out is documented experimentally.
\emph{Detachment} is the operating regime in which the plasma cools enough that
the ion flux to the target collapses, and the regime this diagnostic is chiefly
used to identify. In spectroscopic inversions of divertor Balmer emission at its
onset, molecularly-induced emission is found to account for up to
\SIrange{60}{70}{\percent} of the $D_\alpha$ light and
\SIrange{10}{20}{\percent} of $D_\gamma$~\cite{Wijkamp2023,Verhaegh2021}. An
$H_\alpha/H_\beta$ ratio computed from atomic processes alone is not the ratio
such an instrument records.

\subsection{What survives it, and what does not}

Two consequences follow, and they point in opposite directions, so both must be
stated.

The \emph{numerical} cold-edge results are atomic-model results and nothing
more. Adding a third supply channel that feeds $n=3$ preferentially would change
$\ffrac{3}-\ffrac{4}$, and could plausibly change it a great deal; this work
offers no direct calculation of how much. A bound can nonetheless be
constructed, from the emissivity fractions cited in this section. For a given
upper level the molecular fraction of the emissivity equals the molecular
fraction of the population, because both channels emit through the same
transition probability, and molecular-assisted recombination bypasses the ground
state entirely, so it dilutes the ground-fed fraction as
$\ffrac{p} \to \ffrac{p}(1-\phi_p)$. Taking $\phi_3 = 0.60$ and $\phi_4 = 0.30$
from the reported \SIrange{60}{70}{\percent} of $\mathrm{D}\alpha$ and
\SIrange{10}{20}{\percent} of $\mathrm{D}\gamma$, the difference of ground-fed
fractions at the cold corner falls from $0.4866$ to $0.1088$ and at the density
crest from $0.4580$ to $0.0998$, and the secant sensitivity $\overline{S}$
falls with it, from $0.4822$ to $0.0903$ and from $0.4552$ to $0.0903$.

The plateau error does not fall in the same proportion, because
$\epsplat = |\mathrm{e}^{\overline{S}\Lambda}-1|$ is not linear in
$\overline{S}$; treating it as linear understates the reduction. Imposing the
same fractions on the model's own operator, with a non-negative source
depositing into $n=3$ and $n=4$ so that $(\phi_3,\phi_4)$ is reached exactly at
the plateau, the cold-corner plateau error falls from \SI{38.7}{\percent} to
\SI{6.3}{\percent}, a \textbf{factor of $6.1$}; at the upper end of the
reported ranges, $(\phi_3,\phi_4) = (0.70,0.45)$, it falls to \SI{4.1}{\percent},
a factor of $9.4$. At the crest the same two targets give factors of $5.4$ and
$8.1$, and at the benchmark $3.2$ and $4.4$
(\texttt{validation/molecular\_channel/}).

Two features of that calculation are worth stating, because neither is visible
in the dilution formula. The two fractions cannot be chosen independently: a
source large enough to make $\phi_3 = 0.60$ already drives more than $0.15$ of
$n=4$ through collisional $3\to4$ transfer, so $(\phi_3,\phi_4) = (0.60,0.15)$
requires a negative $n=4$ source and is not realisable on this operator. And
the reduction is not monotone in $\phi_3$. The difference
$\ffrac{3}(1-\phi_3) - \ffrac{4}(1-\phi_4)$ passes through zero near
$\phi_3 = 0.74$ at $\phi_4 = 0.30$, beyond which $n=4$ is the more ground-fed
of the two and the error grows again with $\overline{S}$ carrying the opposite
sign. At the crest, where that zero is reachable with non-negative sources, a
two-source dilution at $\phi_4 = 0.30$ brings the plateau error down to
\SI{0.35}{\percent} at $\phi_3 = 0.75$, past the zero at $\phi_3 = 0.7355$; a
single $n=3$ source there gives \SI{0.30}{\percent} at $\phi_3 = 0.70$ and
\SI{1.1}{\percent} again at $\phi_3 = 0.80$, with $\overline{S}$ negative. A
molecular channel strong enough cancels the diagnostic's sensitivity to the
reservoir rather than merely reducing it, and past the cancellation the error
returns.

The bound is crude, it inherits every uncertainty in the quoted fractions, and
it applies only where those fractions do.

The \emph{structural} results, the list of Section~\ref{sec:transport_partition},
are theorems about whatever matrix $\Lmat$ is supplied and hold for one
containing molecular states; what such a matrix changes is the numerical value
of the ingredients, not the form of the result.

What is \emph{not} available is a calculation with molecules in it: the factors
above are the model's own operator driven by an imposed source of the published
strength, not a solution of a matrix containing molecular states. This is an
atomic model, its cold-edge numbers are atomic-model numbers, and the bound
above is an estimate of how wrong they are rather than a correction to them.

%-----------------------------------------------------------------------------
\section{Where the ridge sits depends on what the model assumed about neutrals}
\label{sec:ridge_conditional}
% QUESTION: is the density at which the error peaks a property of hydrogen, or
% a property of this model's own ionisation balance?

\subsection{The question}

Section~\ref{sec:mechanism} found an interior maximum, a ridge, in the
error map at $\nel \approx 2\times10^{13}\,\si{\per\cubic\centi\metre}$, and
showed that it moves with the shell pair as Griem's criterion for local
thermodynamic equilibrium (LTE) predicts. The mechanism is robust; three
independent tests agreed. But a mechanism and a location are different claims,
and only the first of them has been established.

\subsection{The picture: the ridge is where the two feeds are balanced}

The ridge exists because the $n=3$ and $n=4$ shells switch from
recombination-fed to ground-fed at different reservoir strengths, so their
difference vanishes when either feed dominates and peaks in between. The
controlling variable is therefore not the electron density directly; it is
\emph{how much neutral gas there is to feed from}. The electron density enters
only because, in this model, it determines the neutral density through the CR
ionisation balance.

That is the conditionality. Change what the model assumes about neutrals and the
ridge moves, whatever the electron density does.

\subsection{The measurement}

Within the two-channel split the ground-fed fraction is exactly
$\ffrac{p} = a_1 n_g/(a_0 + a_1 n_g)$, where $n_g$ is the ground-state density
feeding the excitation channel and $a_0$, $a_1$ are the network responses of
Section~\ref{sec:R_depends_on_b}. Scaling $n_g$ is therefore an exact operation
on the ridge, not an approximation to one. Doing it, at
$\Te = \SI{2.02}{\electronvolt}$:

\begin{center}
\begin{tabular}{lr}
  \toprule
  $\ngs$ scaling & density at which $|\ffrac{3}-\ffrac{4}|$ peaks \\
  \midrule
  $\times 0.01$ & $\le 10^{12}\,\si{\per\cubic\centi\metre}$ (off grid) \\
  $\times 0.1$  & $2.28\times10^{12}\,\si{\per\cubic\centi\metre}$ \\
  $\times 1$ (the model's own CRE value) & $1.68\times10^{13}\,\si{\per\cubic\centi\metre}$ \\
  $\times 10$   & $2.48\times10^{14}\,\si{\per\cubic\centi\metre}$ \\
  $\times 100$  & $\ge 10^{15}\,\si{\per\cubic\centi\metre}$ (off grid) \\
  \bottomrule
\end{tabular}
\end{center}

\noindent \textbf{One decade in the assumed neutral density moves the ridge by
roughly one decade in electron density.} The response is close to
proportional, and the whole three-decade span of the density grid is traversed
by a two-decade change in $\ngs$.

The agreement with Griem's criterion at unit scaling, $1.93\times10^{13}$
measured against $2.05\times10^{13}\,\si{\per\cubic\centi\metre}$ predicted, is
therefore an agreement evaluated at \emph{this model's} collisional--radiative
ionisation balance, which contains no transport and no recycling. In a real
divertor the neutral density is set by recycling, which is
Section~\ref{sec:transport}'s objection arriving in a different place. The
sentence ``the ridge sits at
$\nel \approx 2\times10^{13}\,\si{\per\cubic\centi\metre}$'' is a statement about
an ionisation balance, not about a divertor, and should be quoted with that
condition attached or not quoted.

The ridge is also one line pair's and not the Balmer series': the $(3,5)$ pair
of $H_\alpha/H_\gamma$ has its worst density at
$7.2\times10^{12}\,\si{\per\cubic\centi\metre}$, a factor $2.68$ lower, and the
form of words this imposes is stated in Section~\ref{sec:error_maps}.

%-----------------------------------------------------------------------------
\section{The model does not identify the crest with detachment}
\label{sec:not_detachment}
% QUESTION: can the density at which the ridge sits be identified with a named
% operating regime of a tokamak divertor?

\subsection{The question, and the temptation}

The identification is tempting. The ridge sits where ionising and recombining
supplies compete; detachment, the regime ITER must operate in to survive its
exhaust power and the transition this diagnostic is chiefly used to identify
(Section~\ref{sec:molecules}), is where ionising and recombining plasma compete.
The words match.

\subsection{The arithmetic}

The numbers do not. Section~\ref{sec:withdrawn} withdrew the identification:
the ridge, at $1.93\times10^{19}\,\si{\per\cubic\metre}$, sits a factor $5.2$
below the detached-divertor band of published edge
modelling~\cite{Guillemaut2011} and a factor $1.6$ below the modelled upstream
separatrix density~\cite{Stangeby2023,Pitts2019}, in a range neither source
characterises.

The comparison band deserves a word, because a factor of $1.6$ is not a large
margin to be arguing over. Read from the primary source, the band is
\SIrange{3e19}{6e19}{\per\cubic\metre} across the $53$ ITER SOLPS-4.3 cases
\cite[Sec.~3.1.1]{Pitts2019}, and it is the density \emph{at the outer
midplane}, which that paper states explicitly where it plots the dependence
\cite[Sec.~3.2.2]{Pitts2019}. So the quantity is settled: outer midplane, not
X-point. What is not settled by a single factor is where in the band to
measure from. The ridge sits $1.6$ times below the \emph{bottom} of it and
$3.1$ times below the top, and the honest statement is the range rather than
either end. On the same flux tube the outboard-midplane and X-point densities
differ by roughly \SI{50}{\percent}, which is comparable to the width of the
band itself. The conclusion does not turn on any of this, since the ridge is
below the band on every reading, but the factor should not be quoted more
precisely than the band it is measured against.
The repository recorded \cite{Stangeby2023} as ``Nucl.\ Fusion 62 (2022)'' in
\texttt{src/validation/grid.py}; the DOI was right, the volume and year wrong.
It is Nuclear Fusion \textbf{63}(1) 016017 (2023), companion part A at 63(1)
016016 \cite{Stangeby2023a}; \cite{Lore2022} in the same block of that file
genuinely is volume 62.

\subsection{What that leaves}

The honest statement is uncomfortable and should be made anyway. \textbf{The
structure this work found is real and reproducible, and it sits at a density
this thesis cannot identify with any named region of a tokamak.} It is not
detachment; it is not the upstream separatrix; and calling it either would be a
claim about edge transport that a $0$-D atomic model cannot support. The ridge
is a property of how two shells of hydrogen divide their supply. It needs no
divertor geometry to produce it, and it acquires none by being described in one.

The ridge and the detached band do not overlap: restricted to
$\Te \ge \SI{2}{\electronvolt}$ and $\nel \ge 10^{14}\,\si{\per\cubic\centi\metre}$
the census is $0$ of $108$ with a worst estimate of \SI{7.2}{\percent}
(Table~\ref{tab:trapping_census}), so where the diagnostic is most used the
transient error is real but modest.

%-----------------------------------------------------------------------------
\section{The open boundary: an attack that was expected to be fatal}
\label{sec:open_system}
% QUESTION: does holding the ion density fixed manufacture the slow clock that
% the whole result rests on?

\subsection{The question}

Section~\ref{sec:rate_equation} was explicit that the ion density $\nion$ is
held fixed: the recombination feed $\Svec\nion$ enters as a constant source term
a feed stream into a CSTR in the reactor picture, rather than as a
matrix element, and the $43$-state system has no $44$th state for ionised atoms
to accumulate in. Atoms ionise and vanish from the accounting. The slow
eigenvalue $\lambda_0$ is therefore the CR ionisation rate of ground-state
hydrogen against a prescribed reservoir, not the equilibration rate of a closed
system.

\subsection{The falsifier, and what happened}

If the long $\tauslow$ values were an artifact of the open boundary, so would be
the claim that the error outlasts an ELM. The falsifier named in
Section~\ref{sec:attacks} before the test was run is a closure factor of $2300$
at $[0,4]$, the value of $\tauslow/\taudrive$ there. Closing the manifold with
a $44$th ion state gives a closure factor with minimum $1.00$, median $1.00$
and maximum $41.4$ over the grid, moves the ELM census from $202$ to $200$ of
$680$ and the worst case by one part in $325$, and leaves the count above
\SI{2}{\electronvolt} at $45$ before and $45$ after (Table~\ref{tab:closure}).
The factor is large only where $\tauslow$ already exceeds the event duration by
three orders of magnitude and is unity wherever $\tauslow$ approaches
$\taudrive$, so where the boundary condition matters the outcome does not
depend on it. The measured factor at $[0,4]$ is $15$ against the $2300$
required, a cold-corner margin; inside the defended range, at the ridge, it is
$1.01$ against $20$, and at the benchmark $\tauslow$ moves by one part in
$10^{3}$. The attack was survived everywhere it was pressed.

\subsection{The caveat that must not be softened}

The closure adds an ion reservoir. It does \emph{not} close the electron
balance, since $\nel$ is held fixed while the ion population moves, and it adds
no transport whatever. It is a test of the boundary condition on the atomic
manifold, and nothing more. Section~\ref{sec:transport} is the objection this
test does not answer, and no amount of closing the atomic system will answer it.

The two index labels for this point, $[1,4]$ in Table~\ref{tab:closure} and
$[0,4]$ in the census, are the post- and pre-step operators of one heating
step: the census rows carry the pre-step index, every quoted timescale belongs
to the post-step operator that governs the relaxation (Chapter~\ref{ch:results}),
and the two must not be interchanged. Diagonalising each gives
$\tauslow = \SI{0.421012}{\second}$ at $\Lmat[0,4]$ and
\SI{0.233239}{\second} at $\Lmat[1,4]$, so the quoted \SI{0.23324}{\second}
belongs to $\Lmat[1,4]$.

%-----------------------------------------------------------------------------
\section{An external benchmark this model fails, and what the failure measures}
\label{sec:r1_deficit}
% QUESTION: where does the comparison against published atomic data break down,
% and does the failure touch the mechanism?

\subsection{The question}

Chapter~\ref{ch:validation} presented the comparison against Fujimoto's
published population coefficients as the strongest external check this model
passes. It is, for one of the two coefficients. The other disagrees at low
principal quantum number, and this section reports what that disagreement
measures. The short answer is that it is unexplained, that it is not the
$\ell$~closure, and that it is direct evidence about the coefficients
Chapter~\ref{ch:results} uses.

\subsection{The picture: two coefficients, one for each feed}

Chapter~\ref{ch:theory} split every excited population into two pieces
corresponding to the two supply channels. The population coefficient $r_0(p)$
measures how strongly level $p$ is fed from the ion continuum by recombination;
$r_1(p)$ measures how strongly it is fed from the ground state by excitation.
The published tabulation gives both, computed independently, decades ago, by a
different method. Agreement is a real test; disagreement is a real failure.

\subsection{The disagreement, and what it measures}

The recombination-fed coefficients agree. At
$\nel = 10^{12}\,\si{\per\cubic\centi\metre}$, $r_0$ matches to
\SI{1.3}{\percent} at $p=2$ and to \SIrange{12}{14}{\percent} at $p=3$ and
$p=4$. The ground-fed coefficients at low principal quantum number do not: this
model's $r_1$ sits a factor $8.3$ low at $p=3$ and a factor $4.4$ low at $p=4$.

Those are the two shells whose difference is the mechanism of
Chapter~\ref{ch:results}, so the disagreement had to be understood rather than
noted.

Two candidate explanations were eliminated before the source was consulted. If
the excitation rates feeding the ground-fed channel were wrong, the deficit
would follow, but those rates were benchmarked against
R-matrix-with-pseudostates calculations and the dominant $\Delta\ell = 1$
channels agree at a median of $-\SI{5.0}{\percent}$ over $120$ comparisons. And
a normalisation error is excluded by $r_0$ itself: $r_0$ and $r_1$ carry the
same $p$-dependent factor $Z(p)$, so an error large enough to move $r_1$ by a
factor of $8$ would have moved $r_0$ with it.

\subsection{The source's closure, imposed and tested}

Table~4.1(b) is indexed by $p = 2, 3, 4, 5, 7, 10, 15$. Those are principal
quantum numbers, and the table carries no $\ell$ label anywhere. The published
coefficients are shell quantities computed with the $\ell$~substructure
bundled.

This model resolves $\ell$ below $n=9$, so the closures differ. An earlier
draft of this section treated that difference as the explanation, on the
evidence that recomputing the model with proton $\ell$-mixing removed moves
$r_1(2)$ by a factor $118$ and $r_1(3)$ by a factor $3.35$ at
$\nel = 10^{12}\,\si{\per\cubic\centi\metre}$. That argument does not survive
its own test. Removing $\ell$-mixing is the limit in which the process is
absent; the tabulation assumes the limit in which it is infinitely fast, and
those are opposite ends rather than a bracket. Imposing the source's own closure
directly, by bundling this model's operator under statistical $\ell$~weights,
moves $r_1(3)$ by \SI{0.42}{\percent} and makes the $p=2$ agreement worse
(\texttt{verify\_fujimoto\_bundle.py}, stamped in
\texttt{validation/fujimoto\_bundle/}).

The disagreement therefore does \emph{not} measure the $\ell$~closure, and this
thesis does not explain it. An earlier draft of this chapter said the deficit
did not reach $\ffrac{3}-\ffrac{4}$ because that difference never passes
through a bundled quantity; that was wrong, and bundling is not the protection.
Section~\ref{sec:fujimoto} makes $a_p = r_1(p)\,\Zsb{p}/\Zsb{1}$, so a deficit
in $r_1(3)$ and $r_1(4)$ is a deficit in $a_3$ and $a_4$ themselves, and the
only protection is that $\Delta$ and the cap depend on the ratio $a_3/a_4$, in
which a deficit common to both shells cancels. It is not common at
\SI{e12}{\per\cubic\centi\metre}, $8.3$ against $4.4$, and nearly so at
\SI{e15}{\per\cubic\centi\metre}, $1.8$ against $1.9$. Dividing $a_3$ and
$a_4$ by those ratios at every grid point
(\texttt{validation/fujimoto\_r1\_consequence/}) moves $\Delta$ by $+0.64$ and
$-0.08$; over the $280$ points above \SI{2}{\electronvolt} the cap moves by
\SI{+29}{\percent} and \SI{-4}{\percent} at the median, and the sensitivity at
the operating point by \SI{+30}{\percent} and \SI{+11}{\percent} at the median,
with ranges of $-58$ to \SI{+524}{\percent} and $-38$ to \SI{+68}{\percent}.
The disagreement is direct evidence about the coefficients
Chapter~\ref{ch:results} uses; what is open is whether the deficit lies in this
model or in the tabulation, and the stake is stated in
Section~\ref{sec:fujimoto}. The open status is carried in
Section~\ref{sec:scope_open}.

\subsection{Two limitations of the benchmark, stated precisely}

Table~4.1(b) is computed at $T_\mathrm{e} = \SI{1.28e5}{\kelvin}$, which is
\SI{11.03}{\electronvolt}. The grid used here reaches \SI{10}{\electronvolt},
so the comparison is made at the grid edge and \SI{9.3}{\percent} below the
tabulated temperature.

There is no second usable case. The companion Table~4.1(a) is at
$T_\mathrm{e} = \SI{e3}{\kelvin}$, which is \SI{0.0862}{\electronvolt}, two
orders of magnitude below this grid. The external check is therefore one
temperature at one edge, and it should not be described as more than that.

\subsection{What the benchmark establishes, and what it does not}

What the benchmark does establish is the recombination-fed channel, which
agrees, and that is worth stating as a pass rather than leaving implicit beside
a disagreement. What it does not establish is the absolute normalisation of the
ground-fed channel against an independent calculation, because no tabulation
resolved in $\ell$ was available to compare against.

The sensitivity of the density maximum to the ground-fed supply remains a
separate and open question, and Section~\ref{sec:ridge_conditional} treats it
on its own terms: the location responds roughly decade-for-decade to the assumed
neutral density, which is a stronger reason to quote it as a range than
anything the benchmark shows.

%-----------------------------------------------------------------------------
\section{The top of the level ladder, and what the bundling costs}
\label{sec:nmax}
% QUESTION: the model stops at n=15 and bundles the levels above n=8 -- what
% rests on the part that is approximated?

\subsection{The question}

A hydrogen atom has infinitely many bound levels, crowding together towards the
ionisation limit. A computation cannot have infinitely many states, so a
boundary was chosen: this model tracks levels up to $n_{\max} = 15$ and
discards the rest. Above $n = 8$ it also stops resolving orbital angular
momentum $\ell$, treating each shell as a single bundled state populated
statistically. Both are choices, and this section states what they cost.

\subsection{Three consequences, in decreasing order of severity}

\paragraph{The terminal shell is over-populated.} Truncation removes from the
terminal shell's balance the collisional excitation out of it to $n > n_{\max}$,
a way out, and the de-excitation and cascade into it from those levels; an
earlier draft said the shell had ``nothing above it to cascade out to'', which
names the wrong process, since cascade is downward. Against published values
this model's ground-fed coefficient at $n = 15$ is high by a factor of $4.9$ to
$6.2$, while at $n = 10$ it agrees to \SIrange{4}{20}{\percent}; the budget is
in Section~\ref{sec:convergence}. \emph{No result in this thesis rests on the
top shell}: the observable is built from $n = 3$ and $n = 4$, both fully
$\ell$-resolved and far from the truncation.

\paragraph{The bundling nonetheless enters the result indirectly.} The
recombination flux cascading down from the bundled shells feeds the
recombination-fed channel of both $n=3$ and $n=4$, one of the two channels the
mechanism compares, which is why Section~\ref{sec:ridge_alternatives} lists
truncation among the alternative explanations for the ridge that this work does
not exclude.

\paragraph{There is no $\ell$-mixing inside the bundled block.} The
proton-impact $\ell$-mixing that Section~\ref{sec:lmixing} showed to be
essential for the $n = 2$ manifold is present only for $n = 2$--$8$; the
$\ell$-mixing rate array is identically zero above that. Within the bundled
shells the $\ell$-distribution is assumed statistical rather than computed. At
high $n$ that assumption is far more defensible than it would be at low $n$,
because the sublevels are nearly degenerate and every collisional process mixes
them, but it is assumed, not measured, and this thesis has not measured it.

\subsection{The check, run}
\label{sec:bundling_check}

The script that quantifies this was written and, until now, never executed.
Running it produced a result and also produced three defects in the script
itself, which are reported here because two of them would have made its verdict
worthless.

\paragraph{What the script could not do.} It promised to test the bundled
shells against the pipeline's own $\ell$-mixing array, and that is impossible:
the array is identically zero for states $36$ to $42$, in both directions,
because $\ell$-mixing was computed for the resolved block alone. The file
contains nothing to test. The comparison has to be made against the rate
\emph{function} rather than the rate table, and the function is the one
\texttt{compute\_lmix.py} used to build the resolved block, so it is the model's
own rate rather than a second one. Calling it at $n = 2$ to $8$ reproduces the
stored array to $1.0000$, which is a wiring check and not a physical
confirmation; it establishes only that the rate carried up to the bundled
shells is the same function with the same arguments.

\paragraph{The two defects that would have inverted the answer.} The script
carried its own Pengelly--Seaton expression, independent of the pipeline's,
which disagrees with the rate the model actually uses by a factor of $95$ at
$n = 2$ rising to $5.6\times10^{3}$ at $n = 8$, growing roughly as $n^{2}$. Used
as written, it would have reported the mixing rate two to three orders of
magnitude below its true value and could have returned a verdict of marginal
where the answer is not close. It is deleted rather than corrected: a
verification script has no business maintaining a second implementation of a
rate the pipeline already computes. The script also fell back on a synthesised
logarithmic temperature grid of its own when the pipeline's grid files were
absent,
printing a warning and continuing, which is the failure mode described in
Chapter~\ref{ch:validation}: every ratio would have carried a correct-looking
label attached to conditions the model was never evaluated at. It now raises.

\paragraph{The result.} The criterion is that the $\ell$-mixing rate out of the
$np$ substate, the substate the Lyman channel empties fastest, must exceed that
substate's total radiative rate. Over the seven bundled shells and all $400$
grid points, the ratio is above the well-mixed threshold of $10$ at every one of
the $2800$ combinations. The worst case is $n = 9$ at
$\Te = \SI{1}{\electronvolt}$, $\nel = \SI{e12}{\per\cubic\centi\metre}$, where
the mixing rate is $\SI{1.155e10}{\per\second}$ against
$A(9p) = \SI{7.51e6}{\per\second}$, a margin of $1538$, that is, a factor $154$
clear of the threshold (\texttt{verify\_bundling\_psm20.py},
\texttt{outputs/bundling/bundling\_report.txt}).

Two inputs to that comparison were checked rather than assumed. The Einstein
coefficients $A(np\to1s)$ for $n = 9$ to $15$ are hardcoded in the script,
because the model's own radiative file resolves $\ell$ only to $n = 8$ and has
nothing to read above it. Since $A(np\to1s)\,n^{3}$ tends to a constant for
hydrogen, the model's resolved values and the hardcoded table must lie on one
curve: they do, with a step of \SI{0.43}{\percent} across the boundary at
$n = 8$ to $9$ and a monotone approach to $4.13\times10^{9}$ at $n = 15$. And
the $\ell$-mixing coefficient carries a Debye cutoff frozen at
$\nel = \SI{e14}{\per\cubic\centi\metre}$, because the array has no density
axis. That freezing is not small: re-evaluating the cutoff at the local density
changes the $n = 15$ rate by a factor $20$ upward at
\SI{e12}{\per\cubic\centi\metre} and a factor $30$ downward at
\SI{e15}{\per\cubic\centi\metre}. Against the radiative criterion it does not
move the verdict, because the worst point sits at the lowest density, where the
local cutoff raises the rate: the worst ratio is $1538$ with the freezing and
$3.26\times10^{4}$ without it. At the highest densities the same correction
lowers the rate, and that matters for the criterion that follows.

\paragraph{What this closes and what it does not.} The criterion above compares
mixing with radiative decay, and radiative decay is not how a high shell
empties at these densities. The total loss out of a bundled state, $-L_{kk}$,
is \SIrange{89.5}{99.1}{\percent} collisional transfer to other shells,
\SIrange{0.9}{10.5}{\percent} ionisation and at most \SI{0.19}{\percent}
radiative (\texttt{verify\_bundling\_total\_loss.py},
\texttt{validation/bundling\_total\_loss/}). Against that total the margin is
not three orders of magnitude. With the frozen cutoff the model is built with,
mixing exceeds total escape at every one of the $2800$ combinations, by $26$
at $n=9$ falling to $2.85$ at $n=14$, the minimum, all at
\SI{1}{\electronvolt} and \SI{e12}{\per\cubic\centi\metre}, and $288$
combinations sit below $10$. With the cutoff evaluated at the local density
the ratio falls below unity at $121$ of the $2800$, in every bundled shell, at
the two highest densities and $\Te \le \SI{2.95}{\electronvolt}$; at
\SI{1}{\electronvolt} and \SI{e15}{\per\cubic\centi\metre} it runs from $0.86$
at $n=9$ to $0.09$ at $n=14$. The comment in \texttt{compute\_lmix.py} that the
cutoff factor varies by about \SI{30}{\percent} across the grid understates the
variation already at $n=2$, where the local-to-frozen ratio of $F$ runs from
$1.96$ at \SI{e12}{\per\cubic\centi\metre} to $0.53$ at
\SI{e15}{\per\cubic\centi\metre} at \SI{1}{\electronvolt}
(\texttt{validation/bundling\_total\_loss/}, RESULT~4), and fails outright
above $n=9$, where $U_m > 1$ and $F$ falls as $\nel^{-3/2}$; it is reported
here and left in place. What survives is narrower than an
earlier draft claimed. Where mixing wins, the $\ell$-distribution is
statistical whatever the feeds do. Where collisional transfer wins, at high
density, the distribution is set by the feeds, and $n$-changing collisions
between adjacent shells, which carry \SIrange{71}{95}{\percent} of the
collisional loss, are close to $\ell$-blind at high $n$, so the statistical
assumption is plausible there for a different reason, one this thesis has not
tested. It does not close the terminal-shell over-population above, which is a
different defect with a different cause, and it does not make the bundled
block's $\ell$-distribution a computed quantity. It remains assumed, now with
two numbers attached and a stated region where the first of them does not
license it.

%-----------------------------------------------------------------------------
\section{Remaining assumptions, and the sensitivities that are known}
\label{sec:other_assumptions}
% QUESTION: what else was assumed, and which of those assumptions has actually
% been tested?

Four assumptions remain. Each is stated with what is known about it, rather than
listed as a disclaimer, and two of the four have been tested.

\paragraph{Maxwellian electrons: untested.} Every rate coefficient in
Chapter~\ref{ch:atoms} is an average of a cross section over a Maxwell--Boltzmann
electron energy distribution. Ionisation and excitation are driven by the
high-energy tail of that distribution, which holds a small minority of the
electrons. A \SI{1}{\electronvolt} plasma next to a recycling surface has no
guarantee of a Maxwellian tail, and a depleted or enhanced tail would change the
rates that set both clocks. Not tested here.

\paragraph{$T_i = \Te$ in the $\ell$-mixing rates: untested.} The
proton-impact $\ell$-mixing rates of Section~\ref{sec:lmixing} are evaluated at
an ion temperature equal to the electron temperature. In a detaching divertor
the ions are frequently colder than the electrons. Not tested here.

\paragraph{The step is instantaneous: tested, and safe for a reason that is
easy to get wrong.} The error maps impose a sudden temperature step, whereas a
real ELM has a finite rise time. The tempting ratio
$\taurel/\taudrive \approx 10^{-5}$ is about the wrong quantity, the fast
excited-manifold step error (Section~\ref{sec:persistence}); the plateau error
is governed by the slow ground-state response, for which the relevant ratio is
$\tauslow/\taudrive$. At the worst point of the grid that
ratio is $2332$, but that point sits at \SI{1}{\electronvolt} and the figure is
not the margin this result runs on. Over the $448$ defended pairs the ratio has
median $0.815$, and over the $45$ that break down it runs from $1.95$ to $288$,
median $11.2$ (\texttt{validation/divertor\_map/}). The step idealisation is
therefore safe by one to two orders of magnitude where a plateau is reported,
not by three. The governing quantity is in any case
$D_{\rm e} = \tauslow/t_{\rm ramp}$ rather than $\tauslow/\taudrive$: a linear
ramp reaches $D_{\rm e}(1-\mathrm{e}^{-1/D_{\rm e}})$ of the step's plateau
(\texttt{validation/ramp\_plateau/}), so any ramp with
$\taurel \ll t_{\rm ramp} \ll \tauslow$ reaches the same plateau and the
idealisation \emph{is} safe. It is safe because of the slow ratio, and the fast
ratio must never be quoted for or against this result.

\paragraph{The step is in temperature only: tested, and the correction goes the
wrong way.} A real ELM raises $\nel$ as well as $\Te$. The joint map over $4968$
cases (Section~\ref{sec:joint_step}) leaves the exact decomposition unchanged,
with the gain generalising to an additive gradient accurate to
\SI{0.2}{\percent} at the median; over $\Te\ge\SI{2}{\electronvolt}$ the
breakdown fraction moves from $10.0$ to \SI{11.3}{\percent} when the density
steps up one index and the worst case from $0.1748$ to $0.2206$. The earlier
benchmark spot check, a bound falling by a factor $4$, had the right number and
the wrong sign for the grid, because the benchmark and the worst point lie on
opposite sides of the density crest; the fixed-density map understates the
worst case.

\medskip
Two sensitivities of the reported census belong in the same place, because both
are properties of a reporting choice rather than of the physics, and both would
otherwise be mistaken for measurements.

The step size is one interval of the temperature grid, which displaces the
reservoir by $|\ln\xscale|$ between $0.124$ and $0.682$ across the grid. A
four-interval step gives $0.516$ to $2.563$ and raises $\epsplat$ at the ridge
from \SI{18.1}{\percent} to \SI{81.1}{\percent}. On step size alone, therefore,
the reported numbers \emph{understate} the error rather than overstate it.

The \SI{10}{\percent} breakdown threshold is doing real work: at \SI{5}{\percent},
\SI{10}{\percent} and \SI{20}{\percent} the count over the same $680$ points is
$348$, $202$ and $61$ (Section~\ref{sec:persistence}), so the count is not robust
to the threshold, the worst case is, and ``$N$ points break down'' must carry both.

%-----------------------------------------------------------------------------
\section{Scope: five questions this thesis leaves open}
\label{sec:scope_open}
% QUESTION: which questions are open, why can none of them be closed from
% inside a 0-D atomic model, and what would close each?

The objections computed in this chapter were all answerable from inside a
$0$-D collisional--radiative matrix. Five are not, and they are gathered here
rather than left scattered so that a reader can see at once what this thesis
does not determine. Each is stated with the reason it is out of reach and with
the calculation that would settle it. None is a defect of the implementation;
each is a boundary of the model class.

\paragraph{1. The $r_1$ deficit is unexplained.} This model's ground-fed
population coefficient sits a factor $8.3$ below Fujimoto's Table~4.1(b) at
$p=3$, the shell the mechanism runs on, and a factor $1.8$ to $2.0$ below it at
$\log_{10}(\nel/\si{\per\cubic\metre}) = 21$; truncation (\SI{1.2}{\percent}),
Lyman trapping and the source's own statistical-$\ell$ closure
(\SI{0.42}{\percent}) have each been eliminated (Section~\ref{sec:r1_deficit}).
What would settle it is an $\ell$-resolved published tabulation or the atomic
data underlying the tabulated values, neither of which was available. The
deficit is reported as open, with its consequence for $\ffrac{3}-\ffrac{4}$
computed in Section~\ref{sec:fujimoto}.

\paragraph{2. There is no convergence study above $n_{\max} = 15$.} The state
space is cut at $n = 15$ because the atomic data stop there: the quantum
cross sections extend to $n = 10$, shells $11$ to $15$ are already semiclassical
Vriens--Smeets values, and no radiative or recombination coefficients exist
above $15$ in any form. Truncation has been tested downward, and it moves
$\ffrac{3}-\ffrac{4}$ by about \SI{1}{\percent} with the sign changing across
the grid (Section~\ref{sec:convergence}). What is untested is the sensitivity of
the ridge \emph{position} to $n_{\max}$, which would require rebuilding the rate
pipeline at $n_{\max} = 20$ rather than running a validation script: new
excitation, ionisation, radiative and recombination coefficients for five
further shells, and a re-indexing of every array built on the $43$-state layout.
That is a rate-generation exercise, not an analysis one.

\paragraph{3. The ground-state population is solved for, not supplied.} Every
magnitude in Chapter~\ref{ch:results} is conditional on $\ngs$ coming from
$\Lmat\nvec + \Svec\nion = 0$; in a divertor it is set by recycling, on a
transport timescale four orders of magnitude shorter than the local ionisation
balance (Section~\ref{sec:transport}). The open-boundary test of
Section~\ref{sec:attacks} bounds the sensitivity without removing the
conditional, and what would settle it is stated in Section~\ref{sec:transport}.

\paragraph{4. There are no molecular channels in $\Lmat$.} At the onset of
detachment a large fraction of the Balmer light is molecular in origin. The
dilution estimate of Section~\ref{sec:molecules} says how far the cold-corner
number would move under one specific dilution picture, not what a molecular
model would give. What would settle it is a state space containing
$\mathrm{H}_2$, $\mathrm{H}_2^+$ and $\mathrm{H}^-$ with their rate set, a
different model rather than an extension of this one.

\paragraph{5. The electrons are assumed Maxwellian.} The reasons are given in
Section~\ref{sec:other_assumptions}. The rate integrals as built do not accept
an arbitrary electron energy distribution, so this is not a parameter that can
be varied within the present code, and the $T_i = \Te$ assumption in the
$\ell$-mixing rates is untested for the same reason.

\medskip

\noindent Items 3 and 4 are named in Chapter~\ref{ch:conclusions} among the
further work. Items 1 and 2 are bounded in their consequences rather than in
their causes: the thesis states what each would change if it went the wrong
way, and for item 1 that stake is computed in Section~\ref{sec:fujimoto}.

%-----------------------------------------------------------------------------
\section{What is left standing}
\label{sec:limits_summary}
% QUESTION: after all of that, what may this thesis actually claim?

\subsection{What this model cannot say}

It cannot say anything quantitative about the cold edge of its own grid. Below
\SI{2}{\electronvolt} the optically thin plateau error is inflated by a factor
of $2.5$ to $3.3$ by neglected Lyman trapping; the size of the correction
depends on a slab thickness the model does not contain; and below
\SI{1.15}{\electronvolt} even the location of the ridge is undetermined.

It cannot say what a real divertor does, because recycling renews the
ground-state population there on a transport timescale the model does not
contain: \SI{72}{\micro\second} against \SI{233}{\milli\second} at the worst
point, a factor $3.2\times10^{3}$, below the scope boundary; inside the defended
range one to three orders rather than four, with a median $\tauesc/\tauslow$ of
$0.0079$ over the $45$ defended pairs (Table~\ref{tab:transport_selection}); and
$3.6$ at the benchmark, where transport is the \emph{slower} of the two.

It cannot say what an instrument records at the onset of detachment, because a
large fraction of that light is molecular in origin and there are no molecules
in $\Lmat$.

It cannot name the density at which its own ridge sits as a region of a tokamak.
The ridge falls between the modelled upstream separatrix and the detached
divertor band, and belongs to neither.

And it cannot establish the absolute normalisation of its ground-fed
population coefficient against an independent calculation: $r_1$ sits a factor
$8.3$ below the standard tabulation at $p = 3$, the disagreement is
unexplained, and its stake for $\ffrac{3}-\ffrac{4}$ is computed in
Section~\ref{sec:fujimoto} (Sections~\ref{sec:r1_deficit}
and~\ref{sec:scope_open}).

\subsection{What this model can say, with the objections computed}

Above \SI{2}{\electronvolt} the breakdown census is invariant under Lyman
trapping across a twentyfold range of assumed slab thickness, $45$ of $448$ in
every run with the worst case moving by \SI{0.5}{\percent}
(Section~\ref{sec:opacity}). Closing the open boundary moves the census by two
points out of $680$ and the worst case by one part in $325$; the factor $2300$
that would have refuted the result at the cold corner measures $15$, and at the
ridge the factor required is $20$ against a measured $1.01$
(Section~\ref{sec:attacks}). The ridge mechanism was verified over
$\Te \ge \SI{2}{\electronvolt}$, where trapping is now shown to be irrelevant,
so the agreement with Griem's boundary criterion is not an artifact of neglected
opacity. At citable divertor densities above \SI{2}{\electronvolt} the
\SI{100}{\micro\second} time-averaged error never reaches \SI{10}{\percent} over
$108$ points, worst case \SI{7.2}{\percent} (Section~\ref{sec:not_detachment}).

\subsection{The bridge}

The pattern is consistent, and it is why this chapter is not an apology. Every
objection that could be computed inside a $0$-D atomic model was computed, and
the defensible claim came through all of them with its numbers unchanged. The
objections that remain, neutral transport and molecular chemistry, are the
ones that cannot be computed inside a $0$-D atomic model at all, and they have
been stated as conditionals rather than dissolved into caveats.

That is the boundary of the claim. What is inside it is a quantitative statement
about when a spectroscopist may trust a lookup table and when she may not, and
Chapter~\ref{ch:conclusions} puts it into the form she can use.
